\chapter[Classical freeness at the boundary level: Arakawa--Dai]{Classical freeness at the boundary admissible level: the Arakawa--Dai theorems}\label{ch:L20}
\begingroup\setlength{\emergencystretch}{3em}\sloppy

\providecommand{\ltwentyvac}{\mathbf{1}}
\providecommand{\ltwentysl}{\mathfrak{sl}_2}
\providecommand{\ltwentyR}{R}
\providecommand{\ltwentygr}{\operatorname{gr}^F}
\providecommand{\ltwentyJ}{J_\infty}
\providecommand{\ltwentyQ}{Q}
\providecommand{\ltwentyA}{A_m}
\providecommand{\ltwentyB}{B_m}
\providecommand{\ltwentyCm}{C_m}
\providecommand{\ltwentyqq}[1]{(q)_{#1}}
\providecommand{\ltwentyord}{\operatorname{ord}}

This chapter was written against the working tree of 2026-10-13, in which
the files of the Arakawa--Dai campaign were present but not yet committed. They were committed the same day, and every file in the import closure of the umbrella is tracked (checked by the release audit of
2026-10-27). Every Lean quotation below was copied from those files, re-checked by the author on 2026-10-13, and is
checked verbatim against the committed sources by the book's automatic check (\texttt{scripts/textbook/sync.py}). The aggregate \texttt{lake build} was run on
2026-10-13 on this tree and was green (13150 jobs; not re-run for this edition), and the maintained axiom audit
\texttt{scripts/audit/ArakawaDaiAxiomAudit.lean} (160 declarations) passes with only Lean's three standard
axioms and \emph{without} \lean{VOA.pbw} (the PBW theorem of Chapter~\ref{ch:L02}; formerly a project axiom, now a theorem).

The chapter covers the seventy-two files imported by the campaign umbrella \leanpath{VOA/ArakawaDai.lean}
(and, through them, about ninety further new support files in \leanpath{VOA/Commutative/},
\leanpath{VOA/Filtered/} and \leanpath{VOA/Examples/}, some twenty-six thousand lines of new Lean in total).
All headline statements live in one file, \leanpath{VOA/Examples/AffineSl2BoundaryMainTheorems.lean}.

\textbf{Added 2026-10-27.} Two further sections document files that reached the library
later. Section~\ref{l20:sec:vir} is the \emph{Virasoro companion}: classical freeness of the minimal models $M(2,2m+3)$ for every
$m\ge0$, $\ltwentyJ(\C[T]/(T^{m+1}))\cong\ltwentygr L(c_{2,2m+3},0)$, with the full non-reduced $C_2$-algebra and the Poisson vertex
structure (files \leanpath{VOA/Filtered/VirasoroWordLi.lean}, \leanpath{VOA/Examples/VirasoroMinimalClassicalFreeness.lean},
\leanpath{VOA/Examples/VirasoroMinimalPoisson.lean}). Section~\ref{l20:sec:p3} describes thirty-two files of finite commutative algebra
left from an abandoned ``rooted Gordon conductor'' route to the boundary theorem at $p=3$; they are kept because the Deligne carrier of
Chapter~\ref{ch:L24} reuses five of them (imported by five Deligne-carrier files).

\textbf{Added 2026-10-28.} Section~\ref{l20:sec:ad11} presents the file
\leanpath{VOA/\allowbreak Examples/\allowbreak AffineSl2BoundaryTheorem11.lean}, which states Arakawa--Dai's Theorem~1.1 as one
theorem in the paper's form, \lean{arakawaDai_theorem_1_1}, and makes the presentation isomorphism canonical by theorem. This
replaces the earlier ``two-part'' wording of Theorem~1.1 throughout the chapter.

We rely on earlier chapters for the following and do not repeat them. Lean idioms (\lean{Submodule},
quotients \lean{⧸}, \lean{→ₐ[ℂ]}, \lean{≃ₐ[ℂ]}, \lean{Module.finrank}, \lean{PowerSeries}, \lean{MvPolynomial},
\lean{Function.Bijective}, coercions) are in Chapter~\ref{ch:01}. Vertex operator algebras, the mode
convention \lean{W.smul u n v} $=u_{(n)}v$ and the grading \lean{W.grading n} are in
Chapter~\ref{ch:L03}. Ideals, quotient VOAs and the maximal proper graded ideal \lean{QuotVOA.maxGradedIdeal}
are in Chapter~\ref{ch:L04} (Section~\ref{l04:sec:quot}, Definition~\ref{l04:def:maxideal},
Theorem~\ref{l04:thm:maxsimple}). The universal affine VOA $V^k(\ltwentysl)$ (the library's
\lean{affineSl2VOA_closed k hk}), its currents, PBW basis, Wakimoto realisation and Verma modules are in
Chapter~\ref{ch:L08} (Definition~\ref{l08:def:Vk}). Admissible levels and the earlier, \emph{conditional},
boundary $C_2$ programme are in Chapter~\ref{ch:L09} (Sections~\ref{l09:sec:admissible} and
\ref{l09:sec:boundary}). Li's filtration $F^\bullet V$, its exact layers \lean{liWeightLayer W p N} and the
associated graded algebra \lean{liAssociatedGraded W}, finite jets and classical freeness are in
Chapter~\ref{ch:L11} (Sections~\ref{l11:sec:li} and \ref{l11:sec:jets}). Zhu's $C_2$-algebra
$R_V=V/C_2(V)$ (the library's \lean{A₂ W}) with its Poisson structure is in Chapter~\ref{ch:L12}
(Section~\ref{l12:sec:c2}). Quasi-lisse algebras and the example $L_{-4/3}(\ltwentysl)$ are in
Chapter~\ref{ch:L16}.

\paragraph{What this chapter is about, in one paragraph.} Fix an integer $m\ge1$, put $p=2m+1$ and
$k=-2+2/p$ (so $k=-4/3,-8/5,-12/7,\dots$). These are the \emph{boundary admissible} levels of affine
$\ltwentysl$. Let $V=L_k(\ltwentysl)$ be the simple affine vertex algebra at level $k$. Its $C_2$-algebra
$R_V$ is a finitely generated Poisson algebra; the arc algebra $\ltwentyJ R_V$ is a Poisson vertex algebra;
and there is a canonical surjection of Poisson vertex algebras $\ltwentyJ R_V\twoheadrightarrow\ltwentygr V$
onto Li's associated graded algebra. Arakawa and Dai (arXiv:2609.33348) prove that this surjection is an
isomorphism (\emph{classical freeness}, their Theorem~1.1), using the surjection
$\C[e,h,f]/(e\ltwentyQ^m,h\ltwentyQ^m,f\ltwentyQ^m)\twoheadrightarrow R_V$ with $\ltwentyQ=h^2+4ef$ (which then is an
isomorphism), and that the vacuum
character of $V$ is given by an explicit fermionic sum (their Theorem~1.2), which equals the product
$(q^p;q^p)_\infty^3/(q;q)_\infty^3$. The library proves this, for every $m\ge1$, about the
\emph{actual} simple quotient of the \emph{actual} universal affine VOA, with no hypothesis besides $m\ge1$, in the
following exact form. The character theorem is proved as stated. Theorem~1.1 is proved as \emph{one} Lean theorem,
\lean{arakawaDai_theorem_1_1} (Theorem~\ref{l20:thm:ad11}, added 2026-10-28), whose six conjuncts say: (1) the
canonical map $\C[e,h,f]/(e\ltwentyQ^m,h\ltwentyQ^m,f\ltwentyQ^m)\to R_V$ is bijective, (2) with $e,h,f$ sent to the
$C_2$-classes of the three currents; (3) an explicitly constructed algebra $\ltwentyJ R_V$ with an arc
$R_V\to(\ltwentyJ R_V)[[t]]$ has the universal property of the arc algebra of $R_V$; (4) the Poisson vertex algebra
isomorphism $\ltwentyJ R_V\cong\ltwentygr V$ restricts on $R_V\subset\ltwentyJ R_V$ to the canonical map
$R_V\to\ltwentygr V$, (5) and is the \emph{unique} algebra map with that restriction commuting with the translations, i.e.\ it
\emph{is} the paper's canonical map; (6) it preserves the conformal grading in both directions. One honest caveat
remains: the library has no functor $\ltwentyJ$ on abstract commutative algebras, so $\ltwentyJ R_V$ is \emph{realised}
as the presented arc algebra $\ltwentyJ\bigl(\C[e,h,f]/(e\ltwentyQ^m,h\ltwentyQ^m,f\ltwentyQ^m)\bigr)$ together with a
universal arc, and it is characterised up to unique isomorphism by the proved universal property (3). The paper itself
asserts only the surjection in (1); the bijection is a consequence.

\paragraph{The findings of this chapter, in brief.}
\begin{enumerate}
\item \textbf{Both theorems are proved unconditionally} (Theorem~1.1 as the single theorem
  \lean{arakawaDai_theorem_1_1}), with only Lean's
  three standard axioms. (\lean{VOA.pbw} is now a theorem, so ``\lean{pbw}-free'' below records history, not an axiom
  difference.) The ten headline
  declarations (Section~\ref{l20:sec:headlines}) and the paper-form theorem carry only the hypothesis $1\le m$ (besides the
  evident arguments $N$, $x$, $A$ of the statements). They concern
  \lean{simple (level m) (level_noncritical m)}, the quotient of $V^k(\ltwentysl)$ by its maximal proper
  graded ideal, so ``simple'' is not an assumption but a construction.
\item \textbf{The Lean proof is more self-contained than the paper.} Arakawa--Dai use two inputs from
  the literature: the surjection $\C[e,h,f]/(e\ltwentyQ^m,h\ltwentyQ^m,f\ltwentyQ^m)\twoheadrightarrow R_V$
  (from the singular vector, citing Arakawa's earlier work) and, to close the squeeze, the
  Kac--Wakimoto product formula for the character of $L_k(\ltwentysl)$. The library proves \emph{both}.
  The relations come from a singular vector whose existence is derived from the Kac--Kazhdan determinant
  (Section~\ref{l20:sec:m0}); the character lower bound comes from an elementary but decisive linear
  algebra fact, ``kernel dimension of a polynomial matrix at a point is at most the order of vanishing of
  its determinant there'', together with an exact count of determinant zeros (Section~\ref{l20:sec:m6}).
  In particular the previously open presentation of Chapter~\ref{ch:L09}, Section~\ref{l09:sec:boundary},
  is now an unconditional theorem (\lean{baseC2AlgEquiv}); it does \emph{not} prove that the singular
  vector generates the maximal ideal, and does not need to.
\item \textbf{Statements match the paper.} The character is the unshifted, unflavoured series
  $\sum_N\dim V_N\,q^N$, as in Theorem~1.2 of the paper (we checked the paper's displayed formula). The
  PVA isomorphism is stated as a structure (\lean{PoissonVertexAlgebra.Equiv}) that preserves translation
  and every non-negative $\lambda$-bracket coefficient in both directions, and a separate theorem says it
  preserves conformal weight in both directions.
\item \textbf{Points a referee should know} (all in \texttt{warning} boxes): the Poisson vertex algebra on
  the arc side is built on the \emph{presented} arc algebra $\ltwentyJ(\C[e,h,f]/I_m)$, which \emph{serves as}
  $\ltwentyJ R_V$ through the universal arc \lean{c2UniversalArc} and its proved universal property
  (\lean{c2UniversalArc_universal}); there is no $\ltwentyJ$ functor on abstract algebras. The type of
  \lean{actualC2ArcEquiv} is a bijection of sets for each coefficient algebra; its naturality in the coefficient
  algebra is a separate theorem (\lean{actualC2ArcEquiv_natural}), and together they say that the presented arc algebra
  represents the arc functor of $R_V$ (the bundled textbook form is \lean{c2UniversalArc_universal}).
  The type of \lean{baseC2AlgEquiv} is an abstract isomorphism; it is canonical \emph{by theorem}
  (\lean{baseC2AlgEquiv_apply}, \lean{canonicalC2Map_mk_X}). A few docstrings are
  stale (they still call the now-proved character equality ``remaining''). These are listed in
  Section~\ref{l20:sec:defects}.
\item \textbf{The Virasoro companion is also unconditional} (Section~\ref{l20:sec:vir}). For every $m\in\N$ (no hypothesis;
  $m=0$ is $M(2,3)$, $m=1$ is Lee--Yang) the canonical map $\ltwentyJ(\C[T]/(T^{m+1}))\to\ltwentygr L(c_{2,2m+3},0)$ is bijective
  (\lean{VirasoroClassicalFreeness.classicallyFree}), $R_{L}\cong\C[T]/(T^{m+1})$ canonically, and the map is an isomorphism of
  Poisson vertex algebras from the arc algebra with the \emph{zero} $\lambda$-bracket, which is forced: every mode of
  $\ltwentygr L(c_{2,2m+3},0)$ vanishes. Along the way Li's filtration of the universal Virasoro VOA is shown to be the word filtration,
  for every $c$. The injectivity input is a leading-symbol theorem (\lean{virSymbol_twoSeries}), which on its own is \emph{not}
  classical freeness (Section~\ref{l20:sec:vir} explains the bridge).
\item \textbf{The $p=3$ conductor files prove no theorem about a VOA} (Section~\ref{l20:sec:p3}). Their terminal bound is conditional on
  two named hypotheses and a datum \lean{D : P3LinearResidualization W E} that are never instantiated (and \lean{E} is
  connected to no VOA); the files are kept for the finite lemmas that Chapter~\ref{ch:L24} reuses.
\end{enumerate}

%======================================================================
\section{The mathematics}\label{l20:sec:math}
%======================================================================

\subsection{Boundary admissible levels}

Recall from Chapter~\ref{ch:L09} (Section~\ref{l09:sec:admissible}) that a level $k$ of affine $\ltwentysl$
is \emph{admissible} if $k+2=p/q$ with $p,q$ coprime positive integers and $p\ge2$. The admissible levels
are those for which Kac and Wakimoto found modular-invariant families of representations. Among them, the
\emph{boundary} admissible levels are those with the smallest possible numerator, $p=2$:
\begin{equation*}
k+2=\frac{2}{q},\qquad q\ \text{odd},
\end{equation*}
so $k=-2+2/q$; the case $q=1$ gives $k=0$ (the trivial VOA $L_0(\ltwentysl)=\C$), and the genuinely
interesting cases are $q=2m+1\ge3$. The library and this chapter write $p=2m+1$ for this odd denominator
(the paper does the same; do not confuse it with the numerator of Chapter~\ref{ch:L09}, which is $2$ here).
So
\begin{equation*}
p=2m+1,\qquad k=k_m=-2+\frac{2}{p}=-\frac{4m}{2m+1},\qquad m=1,2,3,\dots
\end{equation*}
For example $k_1=-4/3$ (the example of Chapter~\ref{ch:L16}) and $k_2=-8/5$. These levels are
\emph{non-critical} ($k+2\neq0$), so the Sugawara construction makes $V^k(\ltwentysl)$ a vertex operator
algebra with central charge $c=3k/(k+2)=3(-2p+2)/2=3-3p=-6m$ (so $c=-6,-12,-18,\dots$).

The central charge is negative and $V=L_k(\ltwentysl)$ is not $C_2$-cofinite (Chapter~\ref{ch:L16}
proves this at $m=1$): $R_V$ is infinite-dimensional, so $V$ is not ``rational'' in the usual sense.
Nevertheless boundary admissible representations have unusually simple characters: products of theta
functions. This simplicity is the subject of the chapter.

\begin{physics}
\textbf{Argyres--Douglas theories (heuristic dictionary, not used in any proof).} Beem, Lemos, Liendo,
Peelaers, Rastelli and van Rees associated to every four-dimensional $\mathcal N=2$ superconformal field
theory a vertex operator algebra (its ``chiral algebra''), with $c_{2d}=-12c_{4d}$. The vacuum character of
that VOA is the \emph{Schur index} of the 4d theory, and its associated variety
$X_V=\operatorname{Specm}R_V$ is (conjecturally, and in many cases provably) the \emph{Higgs branch}.
For the Argyres--Douglas theories of type $(A_1,D_{2m+1})$ the chiral algebra is
$L_{-4m/(2m+1)}(\ltwentysl)$, i.e.\ exactly the boundary admissible algebras of this chapter (Cordova--Shao,
Creutzig). The Higgs branch is $\C^2/\Z_2$, the nilpotent cone of $\ltwentysl$, which is the reduced variety
of $\C[e,h,f]/(e\ltwentyQ^m,h\ltwentyQ^m,f\ltwentyQ^m)$: the relations say ``$\ltwentyQ^m$ times every
generator is zero'', which set-theoretically forces $\ltwentyQ=0$, the cone. The non-reducedness (nilpotent
directions $\ltwentyQ,\ltwentyQ^2,\dots,\ltwentyQ^{m}$, non-zero in $R_V$ because their degree $2j\le2m$ is below that of the relations) is invisible to the variety but
is part of the ring $R_V$; it is detected by the character.
\end{physics}

\subsection{The $C_2$-algebra, Li's filtration, arcs and classical freeness}

For a VOA $V$ of CFT type, Zhu's $C_2$-algebra $R_V=V/C_2(V)$ is a commutative Poisson algebra under
$[a]\cdot[b]=[a_{(-1)}b]$, $\{[a],[b]\}=[a_{(0)}b]$ (Chapter~\ref{ch:L12}, Theorem~\ref{l12:thm:poisson}).
Li's filtration $V=F^0V\supseteq F^1V=C_2(V)\supseteq F^2V\supseteq\cdots$ (Chapter~\ref{ch:L11},
equation~\eqref{l11:eq:lidef}) counts derivatives, and its associated graded
$\ltwentygr V=\bigoplus_pF^pV/F^{p+1}V$ is a commutative algebra under the $(-1)$-product with a derivation
$\partial$ induced by $L_{-1}$ and with $\lambda$-brackets induced by the non-negative products:
\begin{equation}\label{l20:eq:grbracket}
\{\bar a_\lambda\bar b\}=\sum_{n\ge0}\frac{\lambda^n}{n!}\,\overline{a_{(n)}b},\qquad
\bar a\in F^p/F^{p+1},\ \bar b\in F^q/F^{q+1},\ \overline{a_{(n)}b}\in F^{p+q-n}/F^{p+q-n+1}.
\end{equation}
This makes $\ltwentygr V$ a \emph{Poisson vertex algebra} (Definition~\ref{l20:def:pva} below). Its degree-zero
part is $R_V$.

If $R=\C[x_1,\dots,x_g]/I$, the \emph{arc algebra} $\ltwentyJ R$ is the polynomial ring in the variables
$x_i^{(n)}$, $1\le i\le g$, $n\ge0$, modulo the ideal generated by the coefficients of $t^n$ in
$f\bigl(\sum_n x_1^{(n)}t^n,\dots,\sum_nx_g^{(n)}t^n\bigr)$, $f\in I$. Its $\C$-points are the arcs
$\operatorname{Spec}\C[[t]]\to\operatorname{Spec}R$, i.e.\ algebra maps $R\to\C[[t]]$. If $R$ is Poisson, then
$\ltwentyJ R$ carries a unique Poisson vertex algebra structure with $\partial x^{(n)}=(n+1)x^{(n+1)}$ and
$\{x_\lambda y\}=\{x,y\}$ for $x,y\in R$. The map $R=\ltwentygr_0V\hookrightarrow\ltwentygr V$ extends uniquely
to a homomorphism of differential algebras
\begin{equation}\label{l20:eq:canonical}
\pi_V\colon\ltwentyJ R_V\longrightarrow\ltwentygr V,\qquad x^{(n)}\longmapsto\frac{\partial^n\bar x}{n!}
=\overline{x_{(-n-1)}\ltwentyvac},
\end{equation}
which is surjective (Li) and a homomorphism of Poisson vertex algebras (Arakawa). $V$ is \emph{classically
free} if $\pi_V$ is an isomorphism. Classical freeness fails in general (for many rational VOAs, such as most
Virasoro minimal models, $\pi_V$ has a kernel); it holds for free-field algebras, for universal affine
algebras, and for a few distinguished simple quotients. Andrews, van Ekeren and Heluani conjectured it for
the boundary admissible $L_k(\ltwentysl)$, and Arakawa--Dai proved it.

\begin{physics}
$\ltwentygr V$ is the \emph{classical limit} of the chiral algebra: replace the operator product expansion by
its leading, commutative part, and remember only the singular part of the OPE as a Poisson bracket.
$\ltwentyJ R_V$ is the algebra of functions on the space of \emph{classical field configurations}
$x(t)=\sum x^{(n)}t^n$ on the formal disc with values in the associated variety. Classical freeness says the
quantum theory has no relations beyond those forced classically, \emph{and their derivatives}: the vacuum
module, graded by conformal weight, has the same size as the space of classical fields on the (possibly
non-reduced) Higgs branch. In 4d language: the Schur index equals the Hilbert series of the arc space of the
Higgs branch \emph{scheme} (a ``Higgs branch reconstruction'' statement). This is a heuristic slogan; the
proof below is pure algebra.
\end{physics}

\subsection{The two theorems}

\begin{theorem}[Arakawa--Dai, Theorem 1.1, with the presentation as a corollary]\label{l20:thm:paper11}
Let $m\ge1$, $p=2m+1$, $k=-2+2/p$, $V=L_k(\ltwentysl)$, $\ltwentyQ=h^2+4ef\in\C[e,h,f]$. Then:
\begin{enumerate}
\item $R_V\cong\ltwentyA:=\C[e,h,f]/(e\ltwentyQ^m,h\ltwentyQ^m,f\ltwentyQ^m)$, via $e,h,f\mapsto$ the
  classes of the three currents;
\item the canonical map $\pi_V\colon\ltwentyJ R_V\to\ltwentygr V$ is an isomorphism of Poisson vertex
  algebras preserving the conformal grading.
\end{enumerate}
\end{theorem}
\noindent\textbf{Attribution.} The paper's displayed Theorem~1.1 is part~2. For part~1 the paper states only the
\emph{surjection} $\ltwentyA\twoheadrightarrow R_V$ (citing earlier work); that it is an isomorphism is a corollary of
part~2 (the argument is in the proof of Theorem~\ref{l20:thm:c2arc}), which the library proves as
\lean{baseC2AlgEquiv}. So part~1 is presented here as a consequence of the paper's theorem, not as a claim of the paper.
Both parts together are the single Lean theorem \lean{arakawaDai_theorem_1_1}, presented as Theorem~\ref{l20:thm:ad11}
in Section~\ref{l20:sec:headlines}.

\begin{physics}
Part 1 says the Higgs branch chiral ring of the $(A_1,D_{2m+1})$ theory is exactly
$\C[e,h,f]/(x\ltwentyQ^m)$, which is non-reduced ($\ltwentyQ\neq0$ but $\ltwentyQ^{m+1}=0$); part 2 says the whole chiral algebra is, classically, nothing but
field configurations on that scheme. The Lean statements of both parts are in Section~\ref{l20:sec:headlines}.
\end{physics}

\begin{theorem}[Arakawa--Dai, Theorem 1.2]\label{l20:thm:paper12}
With the same notation, and $(q)_n=\prod_{j=1}^n(1-q^j)$, $(q)_\infty=\prod_{j\ge1}(1-q^j)$,
\begin{equation}\label{l20:eq:fermionic}
\sum_{N\ge0}\dim V_N\,q^N=\frac{1}{(q)_\infty^2}\sum_{n_1\ge n_2\ge\cdots\ge n_m\ge0}
\frac{q^{\sum_{j=1}^mn_j(n_j+1)}\,(q)_{n_m}^2}{(q)_{n_1-n_2}\cdots(q)_{n_{m-1}-n_m}\,(q)_{2n_m+1}}
=\frac{(q^p;q^p)_\infty^3}{(q;q)_\infty^3}.
\end{equation}
\end{theorem}

The character is \emph{unshifted} (no $q^{-c/24}$) and \emph{unflavoured} (no fugacity $z^{h_0}$ for the
$\ltwentysl$-charge). For $m=1$ the product is $(1-q^3)^3(1-q^6)^3\cdots/(1-q)^3(1-q^2)^3\cdots$, whose
first coefficients are $1,3,9,19,42,81,155,\dots$, while $V^k(\ltwentysl)$ has $1,3,9,22,51,108,221,\dots$
(Example~\ref{l20:ex:numbers}).

\subsection{The proof strategy}

The proof is a \emph{Hilbert-series squeeze}. It compares two graded vector spaces through one map:
\begin{equation*}
\ltwentyB:=\ltwentyJ\ltwentyA\ \xrightarrow{\ \ \pi\ \ }\ \ltwentygr V,\qquad
\dim(\ltwentygr V)_N=\dim V_N .
\end{equation*}
Here $\pi$ is the composite of $\ltwentyJ\ltwentyA\twoheadrightarrow\ltwentyJ R_V$ (from part~1's surjection)
with $\pi_V$; it is surjective, so $\dim V_N\le\dim(\ltwentyB)_N$. If one shows
\begin{equation}\label{l20:eq:squeeze}
\dim(\ltwentyB)_N\ \le\ [q^N]\frac{(q^p;q^p)^3_\infty}{(q)^3_\infty}\ \le\ \dim V_N\qquad(N\ge0),
\end{equation}
then all three numbers are equal, $\pi$ is injective in every conformal weight, hence bijective; since
$\ltwentyB\to\ltwentyJ R_V\to\ltwentygr V$ is then injective, also $\ltwentyA\to R_V$ is injective (restrict to
constant arcs, which embed $\ltwentyA$ in $\ltwentyB$; details in the proof of Theorem~\ref{l20:thm:c2arc}), which is
part~1; and the character is the product, which with a $q$-series identity is Theorem~1.2. The work splits
into the following nine steps, which this chapter treats in order. We name each step by its content and
refer to it by that name throughout.

\begin{enumerate}
\item[(1)] \emph{The boundary relations hold, and the currents generate} (Section~\ref{l20:sec:m0}): the relations
  $e\ltwentyQ^m, h\ltwentyQ^m, f\ltwentyQ^m$ lie in the maximal ideal, so they hold in $R_V$, and
  $\ltwentyA\twoheadrightarrow R_V$.
\item[(2)] \emph{Degeneration to the auxiliary arc algebra} (Section~\ref{l20:sec:m1}): $\ltwentyB$ degenerates to
  the arc algebra $\ltwentyCm=J(\C[e,h,f]/(x(ef)^m))$, which is at least as large in every weight.
\item[(3)] \emph{Restricted dual and Gordon divisibility} (Section~\ref{l20:sec:m23}): a functional model of the
  restricted dual of $\ltwentyCm$ and the divisibility conditions of Gordon type.
\item[(4)] \emph{The partition filtration} (Section~\ref{l20:sec:m23}): a pair/partition filtration whose
  associated graded pieces give a symmetric Hilbert series bound.
\item[(5)] \emph{The $q$-series identities} (Section~\ref{l20:sec:m4}): the fermionic chain equals the product.
\item[(6)] \emph{The Hilbert-series upper bound} (Section~\ref{l20:sec:m5}): steps (2)--(5) combine to
  $\dim(\ltwentyB)_N\le[q^N]\,\text{product}$.
\item[(7)] \emph{The lower bound for the character} (Section~\ref{l20:sec:m6}): $[q^N]\,\text{product}\le\dim V_N$,
  from the Gram rank of $L_k$, the root order of the vacuum determinant, the cancellation against the reflected
  determinant, and the identification of the boundary resonances with the Weyl correction.
\item[(8)] \emph{The conformal Hilbert squeeze} (Section~\ref{l20:sec:m7}): every conformal piece of
  $\ltwentyB\to\ltwentygr V$ is injective.
\item[(9)] \emph{The Poisson-vertex-algebra structure and the intrinsic arc equivalence}
  (Section~\ref{l20:sec:m8}): the Poisson vertex algebra structure, the inverse map, the conformal grading and the
  $C_2$ presentation.
\end{enumerate}
The project's campaign ledger \texttt{Notes/ARAKAWA\_DAI\_2609\_33348\_CAMPAIGN.md} labels these steps
\texttt{M0}, \ldots, \texttt{M8} in the same order; those labels occur in some file and declaration names, but we do
not use them.

\begin{warning}[title={Where the library departs from the paper}]
The paper obtains the second inequality in \eqref{l20:eq:squeeze} by quoting the Kac--Wakimoto character
formula for boundary admissible representations. The library does \emph{not} quote it: the step \emph{the lower bound for the character} (Section~\ref{l20:sec:m6}) proves the
lower bound from scratch. Likewise the paper quotes the surjection $\ltwentyA\twoheadrightarrow R_V$; the
library proves it (Section~\ref{l20:sec:m0}). The library also never uses (and never proves) the statement that the maximal ideal
of $V^k(\ltwentysl)$ is generated by one singular vector (the Adamovi\'c--Milas theorem quoted as a hypothesis
in Chapter~\ref{ch:L09}). This makes the formalisation strictly more self-contained than the paper; it is not
a discrepancy in the statements.
\end{warning}

\begin{physics}
The squeeze has a transparent physical meaning. The upper bound counts \emph{classical} states: monomials in
the modes of the currents $e(z),h(z),f(z)$ modulo the classical constraints $x\ltwentyQ^m=0$ and their
derivatives. These are counted by a \emph{fermionic} (``quasi-particle'') formula: particles of three
colours with an exclusion statistics encoded in $q^{\sum n_j(n_j+1)}/(q)_{n_j-n_{j+1}}$, in the spirit of
Feigin--Stoyanovsky and of the Andrews--Gordon identities. The lower bound counts \emph{quantum} states: the
universal vacuum module minus the null states, and the null states are located by the zeros of the Kac
determinant (the Shapovalov form) as a function of the level. Equality says every classical state is realised
quantum mechanically and vice versa.
\end{physics}

%======================================================================
\section{The objects in the headline statements}\label{l20:sec:objects}
%======================================================================

This section defines, mathematically and by their Lean source, every object that appears in the ten
headline statements of Section~\ref{l20:sec:headlines}. We go from the level, through the VOA and its
character, the $q$-series, the commutative algebra of arcs, Poisson vertex algebras, to the canonical
map.

\subsection{The level and the VOA}

\begin{definition}[Boundary denominator and level]\label{l20:def:level}
For $n\in\N$ put $\mathrm{denominator}(n)=2n+1$ and $\mathrm{level}(n)=-2+2/(2n+1)\in\C$.
\end{definition}
\noindent File \leanpath{VOA/\allowbreak Examples/\allowbreak AffineSl2BoundarySingularFamily.lean}
(Chapter~\ref{ch:L09}).
\begin{leancode}
def denominator (n : ℕ) : ℕ := 2 * n + 1

def level (n : ℕ) : ℂ := -2 + 2 / (denominator n : ℂ)

theorem level_noncritical (n : ℕ) : level n + 2 ≠ 0 := by
\end{leancode}

\begin{dictionary}
Both live in the namespace \lean{VOA.AffineSl2BoundarySingularFamily}. \lean{(denominator n : ℂ)} is the
cast of the natural number $2n+1$ into $\C$; division in \lean{ℂ} is the field division (no truncation, since
the cast is non-zero). So \lean{level n} is exactly $k_n=-2+2/(2n+1)$. At $n=0$ it is $0$; the headline
theorems all assume $1\le m$, so $n=0$ never matters. \lean{level_noncritical} is the proof that
$k_n+2=2/(2n+1)\neq0$; it is needed because the universal VOA is only constructed at non-critical level
(Chapter~\ref{ch:L08}). The paper's $p$ is \lean{denominator m}. The Lean convention agrees with the book's.
\end{dictionary}

\begin{physics}
$k+2$ is the ``shifted level'' (level plus dual Coxeter number $h^\vee=2$); the Sugawara stress tensor has a
factor $1/(k+2)$. Boundary admissible means $k+2=2/p$: in the Argyres--Douglas dictionary $p=2m+1$ labels the
theory $(A_1,D_{2m+1})$.
\end{physics}

\begin{status}
Definitions; \lean{level_noncritical} proved, axioms \lean{[propext, Classical.choice, Quot.sound]} (audited
as a dependency of the headlines).
\end{status}

\begin{definition}[Universal and simple affine VOA; the boundary vacuum algebra]\label{l20:def:simple}
For $k\neq-2$, $\mathrm{universal}(k)=V^k(\ltwentysl)$ is the universal affine VOA with its Sugawara conformal
vector, and $\mathrm{simple}(k)=V^k(\ltwentysl)/\mathfrak m$, where $\mathfrak m$ is the maximal proper graded
ideal (the sum of all ideals that are graded and do not contain the vacuum). The boundary vacuum algebra is
$\mathrm{vacuum}(m)=\mathrm{simple}(k_m)=L_{k_m}(\ltwentysl)$.
\end{definition}
\noindent File \leanpath{VOA/\allowbreak Examples/\allowbreak AffineSl2GeneralRigidityBasis.lean} (Chapter~\ref{ch:L09}).
\begin{leancode}
noncomputable abbrev universal := affineSl2VOA_closed k hk
noncomputable abbrev simple := QuotVOA.maxQuotientVOA (universal k hk)
\end{leancode}
\noindent File \leanpath{VOA/\allowbreak Core/\allowbreak QuotientVOA.lean} (Chapter~\ref{ch:L04}).
\begin{leancode}
def ProperGradedIdeals (W : VertexOperatorAlgebra 𝒱) : Set (Submodule ℂ 𝒱) :=
  {I | W.IsIdeal I ∧ IsGradedIdeal W I ∧ W.vacuum ∉ I}

def maxGradedIdeal (W : VertexOperatorAlgebra 𝒱) : Submodule ℂ 𝒱 :=
  sSup (ProperGradedIdeals W)

def maxQuotientVOA (W : VertexOperatorAlgebra 𝒱) :
    VertexOperatorAlgebra (𝒱 ⧸ maxGradedIdeal W) :=
  quotientVOA W (maxGradedIdeal W) (maxGradedIdeal_isIdeal W) (maxGradedIdeal_isGradedIdeal W)
\end{leancode}
\noindent File \leanpath{VOA/\allowbreak Examples/\allowbreak AffineSl2BoundaryVacuumCharacter.lean}.
\begin{leancode}
abbrev vacuum (m : ℕ) := simple (level m) (level_noncritical m)
\end{leancode}

\begin{dictionary}
In \leanpath{AffineSl2GeneralRigidityBasis.lean} the variables \lean{(k : ℂ) (hk : k + 2 ≠ 0)} are section
variables, so \lean{universal k hk} and \lean{simple k hk} take them as arguments.
\lean{affineSl2VOA_closed k hk} is $V^k(\ltwentysl)$ of Chapter~\ref{ch:L08}, Definition~\ref{l08:def:Vk}
(underlying space \lean{WakModule k}, the vacuum module, with its proved PBW basis). \lean{sSup} of a set of
submodules is the submodule they span. \lean{W.IsIdeal I} (two-sided ideal) and \lean{IsGradedIdeal W I}
(the sum of its intersections with the weight spaces) are the notions of Chapter~\ref{ch:L04},
Section~\ref{l04:sec:quot}. Chapter~\ref{ch:L04} proves that this supremum is again a proper graded ideal
(\lean{maxGradedIdeal_mem}), that the quotient is simple (\lean{isSimple_maxQuotientVOA'}, here
\lean{simple_isSimple}), and that every proper graded ideal is contained in it. For a VOA of CFT type every
ideal is graded and proper ideals avoid the vacuum, so $\mathfrak m$ is \emph{the} maximal ideal and
\lean{simple k hk} is \emph{the} simple quotient $L_k(\ltwentysl)$. The non-vanishing of the quotient
(\lean{simple_vacuum_ne_zero}) is proved from the PBW basis. \lean{abbrev} means the name is a transparent
abbreviation: \lean{vacuum m} \emph{is} \lean{simple (level m) (level_noncritical m)}. The name ``vacuum''
refers to the vacuum module; it is the whole VOA, not the vacuum vector.

\textbf{Could it be vacuous?} A trap would be to define $L_k$ as a quotient by an ideal generated by a chosen
singular vector, which is $L_k$ only by a theorem (Adamovi\'c--Milas). Here the quotient is by the genuine
maximal ideal, defined without reference to any vector; so every statement below is about the actual
simple VOA.
\end{dictionary}

\begin{physics}
$V^k(\ltwentysl)$ is the chiral algebra generated by three currents $e(z),h(z),f(z)$ with the $\ltwentysl$
current algebra OPE at level $k$; $L_k(\ltwentysl)$ is the physical theory in which every null state (a state
orthogonal to everything under the Shapovalov form) has been set to zero. Only $L_k$ can be a consistent
CFT, and only $L_k$ is the chiral algebra of an Argyres--Douglas theory.
\end{physics}

\begin{status}
Definitions; the simplicity and non-triviality theorems are proved in Chapters~\ref{ch:L04},
\ref{ch:L09} without \lean{VOA.pbw}.
\end{status}

\subsection{The character and the product}

\begin{definition}[Vacuum character and boundary product]\label{l20:def:character}
$\mathrm{character}(m)=\sum_{N\ge0}\dim_\C L_{k_m}(\ltwentysl)_N\,q^N\in\Q[[q]]$, and
$\mathrm{product}(m)=(q^p;q^p)_\infty^3\cdot(q;q)_\infty^{-3}$ with $p=2m+1$.
\end{definition}
\noindent File \leanpath{VOA/\allowbreak Examples/\allowbreak AffineSl2BoundaryVacuumCharacter.lean}.
\begin{leancode}
def character (m : ℕ) : PowerSeries ℚ :=
  PowerSeries.mk (fun N => (Module.finrank ℂ ((vacuum m).grading N) : ℚ))

def product (m : ℕ) : PowerSeries ℚ :=
  (PowerSeries.expand (denominator m) (by dsimp [denominator]; omega) euler) ^ 3 *
    reciprocalEuler ^ 3
\end{leancode}

\begin{dictionary}
\lean{PowerSeries ℚ} is $\Q[[q]]$; \lean{PowerSeries.mk a} is $\sum_Na_Nq^N$. \lean{(vacuum m).grading N} is
the weight-$N$ subspace $V_N$ (the integer $N$ coerced from \lean{ℕ} to \lean{ℤ}); negative weights are
omitted, and they are zero anyway because $V$ is of CFT type. \lean{Module.finrank ℂ} is the dimension (which
is finite, Chapter~\ref{ch:L08}), cast to $\Q$. So \lean{character m} is the unshifted, unflavoured vacuum
character, defined from the actual dimensions and \emph{independently} of any formula. This is the
ordinary character of the paper's Theorem~1.2.

\lean{euler} and \lean{reciprocalEuler} are $(q;q)_\infty$ and its inverse (Definition~\ref{l20:def:euler}).
\lean{PowerSeries.expand p hp g} is Mathlib's substitution $q\mapsto q^p$, i.e.\ $g(q^p)$; the proof
argument \lean{(by dsimp [denominator]; omega)} only certifies $p\neq0$. Hence
\lean{product m} $=(q^p;q^p)^3_\infty/(q;q)^3_\infty$, the right-hand side of \eqref{l20:eq:fermionic}.
\end{dictionary}

\begin{physics}
\lean{character m} is the vacuum character $\operatorname{tr}_Vq^{L_0}$ (the Schur index of the
$(A_1,D_{2m+1})$ theory with flavour fugacity set to one and without the Casimir-energy prefactor). The
product is the plethystic exponential $\operatorname{PE}\bigl[3(q-q^p)/((1-q)(1-q^p))\bigr]$: three bosonic
fields of weight one (the currents, the factor $1/(q)^3_\infty$) with a ``fermionic'' constraint at every
multiple of $p$.
\end{physics}

\begin{status}
Definitions.
\end{status}

\begin{definition}[$q$-Pochhammer symbols and the Euler product]\label{l20:def:euler}
$(a;q)_n=\prod_{j=0}^{n-1}(1-aq^j)$, $(q)_n=(q;q)_n$, and $(q^p;q^p)_\infty=\prod_{j\ge1}(1-q^{pj})$ is $(q)_\infty$
with $q$ replaced by $q^p$. In $\Q[[q]]$ (or $R[[q]]$ for any commutative ring)
$(q)_\infty$ is the power series whose $d$-th coefficient is the $d$-th coefficient of $(q)_d$; and
$(q)_\infty^{-1}$, $(q)_n^{-1}$ are the inverses of these power series (which exist because the constant term
is $1$).
\end{definition}
\noindent File \leanpath{VOA/\allowbreak Commutative/\allowbreak QPochhammer.lean}.
\begin{leancode}
def pochhammer (a q : R) (n : ℕ) : R :=
  ∏ j ∈ Finset.range n, (1 - a * q ^ j)

def qFactorial (q : R) (n : ℕ) : R := pochhammer q q n

noncomputable def euler : PowerSeries R :=
  PowerSeries.mk fun d => PowerSeries.coeff d (qFactorial (PowerSeries.X : PowerSeries R) d)

theorem coeff_euler (d n : ℕ) (h : d ≤ n) :
    PowerSeries.coeff d (euler : PowerSeries R) =
      PowerSeries.coeff d (qFactorial (PowerSeries.X : PowerSeries R) n) := by

noncomputable def reciprocalQFactorial (n : ℕ) : PowerSeries R :=
  PowerSeries.invOfUnit (qFactorial PowerSeries.X n) 1

noncomputable def reciprocalEuler : PowerSeries R := PowerSeries.invOfUnit euler 1

@[simp] theorem euler_mul_reciprocal :
    (euler : PowerSeries R) * reciprocalEuler = 1 := by
\end{leancode}

\begin{dictionary}
All in namespace \lean{VOA.QSeries}. \lean{PowerSeries.X} is $q$. \lean{qFactorial X n} is
$\prod_{j=0}^{n-1}(1-q\cdot q^j)=(1-q)\cdots(1-q^n)$. The infinite product is defined
coefficientwise by a sufficiently long truncation; \lean{coeff_euler} certifies that any truncation $n\ge d$
gives the same $d$-th coefficient, which is exactly why the definition is the infinite product (the factor
$1-q^{j}$ with $j>d$ does not affect $q^d$). \lean{PowerSeries.invOfUnit f u} is Mathlib's inverse of a power
series whose constant term is the unit $u$ (here $1$); \lean{euler_mul_reciprocal} certifies
$(q)_\infty\cdot(q)_\infty^{-1}=1$. The Lean convention agrees with the book's.
\end{dictionary}

\begin{physics}
$1/(q)_\infty=\sum_N p(N)q^N$ is the partition function of one chiral boson (one oscillator $a_{-n}$ for
each $n\ge1$): states are partitions of the energy.
\end{physics}

\begin{status}
Definitions and elementary lemmas; audited (\lean{coeff_euler}, \lean{coeff_reciprocalEuler}): standard
axioms only.
\end{status}

\begin{definition}[Locally finite sums of power series]\label{l20:def:lfsum}
For a family $(F_i)_{i\in\iota}$ in $R[[q]]$ such that each coefficient $[q^n]F_i$ is non-zero for only
finitely many $i$, $\sum_iF_i$ is the power series with $n$-th coefficient $\sum_i[q^n]F_i$.
\end{definition}
\noindent File \leanpath{VOA/\allowbreak Commutative/\allowbreak LocallyFinitePowerSeriesSum.lean}.
\begin{leancode}
def IsCoefficientwiseFinite (F : ι → PowerSeries R) : Prop :=
  ∀ n, (Function.support fun i => PowerSeries.coeff n (F i)).Finite

def sum (F : ι → PowerSeries R) : PowerSeries R :=
  PowerSeries.mk fun n => ∑ᶠ i, PowerSeries.coeff n (F i)
\end{leancode}

\begin{dictionary}
\lean{∑ᶠ i, g i} is Mathlib's \lean{finsum}: the sum of $g$ over its support \emph{if the support is finite},
and $0$ otherwise. So \lean{sum F} is total (defined for every family), and agrees with the mathematical sum
only when \lean{IsCoefficientwiseFinite F} holds. \textbf{Vacuity check:} every sum used in the headlines
(\lean{fermionicTrace}, \lean{hilbertBoundSeries}, \lean{multiplicitySource}) has its local finiteness proved
in the library (\lean{fermionicTrace_finite}, \lean{hilbertBoundSeries_finite},
\lean{multiplicityTerm_finite}), and more to the point the summands there are bounded below in $q$-degree by
a quantity growing with the index, so each coefficient is a finite sum on paper too.
\end{dictionary}

\begin{physics}
Fermionic character formulas are sums over particle numbers; this is the bookkeeping device that makes such
infinite sums of $q$-series literal algebraic objects.
\end{physics}

\begin{status}
Definitions.
\end{status}

\begin{definition}[The fermionic chain and the Hilbert-bound series]\label{l20:def:chain}
For $\ell,n\ge0$ define $K_\ell(n)\in\Q[[q]]$ by
\begin{equation*}
K_0(n)=\frac{(q)_n^2}{(q)_{2n+1}},\qquad K_{\ell+1}(n)=\sum_{j=0}^n\frac{q^{j(j+1)}}{(q)_{n-j}}K_\ell(j),
\end{equation*}
and the \emph{fermionic trace} $T_m=\sum_{n\ge0}q^{n(n+1)}K_{m-1}(n)$. Unwinding the recursion,
$T_m=\sum_{n_1\ge\cdots\ge n_m\ge0}q^{\sum n_j(n_j+1)}(q)_{n_m}^2/\bigl((q)_{n_1-n_2}\cdots
(q)_{n_{m-1}-n_m}(q)_{2n_m+1}\bigr)$, the sum in \eqref{l20:eq:fermionic}. The \emph{Hilbert-bound series}
$S_m$ (equation (5.12) of the paper) is the same chain with $K_0$ replaced by the three-species residual sum
\begin{equation*}
\mathrm{res}(n)=\sum_{s\ge0}q^{(n+1)s}\,\frac{1}{(q)_s}\,E(n+s+1)\,E(n+1)\,\frac{1}{(q)_n},\qquad
E(A)=\sum_{j\ge0}\frac{q^{Aj}}{(q)_j}.
\end{equation*}
\end{definition}
\noindent File \leanpath{VOA/\allowbreak Commutative/\allowbreak QChainTrace.lean}.
\begin{leancode}
def fermionicKernel : ℕ → ℕ → PowerSeries K
  | 0, n => qFactorial PowerSeries.X n ^ 2 * reciprocalQFactorial (2 * n + 1)
  | ℓ + 1, n => ∑ j ∈ Finset.range (n + 1),
      PowerSeries.X ^ (j * (j + 1)) * reciprocalQFactorial (n - j) * fermionicKernel ℓ j

def fermionicTrace (m : ℕ) : PowerSeries K :=
  VOA.LocallyFinitePowerSeriesSum.sum (fun n =>
    PowerSeries.X ^ (n * (n + 1)) * fermionicKernel (m - 1) n)
\end{leancode}
\noindent Files \leanpath{VOA/\allowbreak Commutative/\allowbreak QBinomialResidual.lean},
\leanpath{VOA/\allowbreak Commutative/\allowbreak QEuler.lean},
\leanpath{VOA/\allowbreak Commutative/\allowbreak QSeriesSum.lean}.
\begin{leancode}
def residualSum (n : ℕ) : PowerSeries R :=
  powerShiftSum (n + 1) fun s =>
    reciprocalQFactorial s * eulerSum (n + s + 1) *
      eulerSum (n + 1) * reciprocalQFactorial n

def eulerSum (A : ℕ) : PowerSeries R :=
  powerShiftSum A reciprocalQFactorial

def powerShiftSum (A : ℕ) (B : ℕ → PowerSeries R) : PowerSeries R :=
  VOA.LocallyFinitePowerSeriesSum.sum (fun j => PowerSeries.X ^ (A * j) * B j)
\end{leancode}
\noindent File \leanpath{VOA/\allowbreak Commutative/\allowbreak QHilbertProduct.lean}.
\begin{leancode}
def hilbertBoundKernel : ℕ → ℕ → PowerSeries K
  | 0, n => residualSum n
  | ℓ + 1, n => ∑ j ∈ Finset.range (n + 1),
      PowerSeries.X ^ (j * (j + 1)) * reciprocalQFactorial (n - j) * hilbertBoundKernel ℓ j

def hilbertBoundSeries (m : ℕ) : PowerSeries K :=
  VOA.LocallyFinitePowerSeriesSum.sum (fun n =>
    PowerSeries.X ^ (n * (n + 1)) * hilbertBoundKernel (m - 1) n)
\end{leancode}

\begin{dictionary}
\lean{K} is any field (here $\Q$). \lean{Finset.range (n + 1)} is $\{0,\dots,n\}$. \lean{n - j} with
$j\le n$ is honest subtraction. \lean{m - 1} is truncated subtraction in \lean{ℕ} ($0-1=0$); this only matters
at $m=0$, which the headlines exclude. For $m\ge1$, \lean{fermionicTrace m} sums over $n=n_1$ and
\lean{fermionicKernel (m - 1) n} contributes the $m-1$ further nested sums over $n_2,\dots,n_m$, as in the paper. \lean{fermionicKernel 0 n} is $(q)_n^2\cdot(q)_{2n+1}^{-1}$, the ``terminal
$q$-factorial ratio'' of \eqref{l20:eq:fermionic}. So \lean{fermionicTrace m} is literally the sum
$T_m$ of the paper, and
\begin{equation*}
\frac{1}{(q)_\infty^2}\,T_m\ =\ \text{the middle expression of \eqref{l20:eq:fermionic}}
\ =\ \texttt{reciprocalEuler\^{}2 * fermionicTrace m}.
\end{equation*}
\lean{hilbertBoundSeries m} is the paper's $S_m$: the same chain, with the terminal factor replaced by the
three residual species (two Euler sums and one $q$-binomial-like sum). It is the series that the arc
calculation produces; the identity $S_m=T_m/(q)_\infty^2$ is proved among the $q$-series identities (Theorem~\ref{l20:thm:qseries}).
\end{dictionary}

\begin{physics}
In a quasi-particle basis, $n_j$ counts particles of ``charge at least $j$''; $q^{n_j(n_j+1)}$ is the
exclusion energy cost (each new particle pays for the ones already present), and $1/(q)_{n_j-n_{j+1}}$ is the
free motion of the $n_j-n_{j+1}$ particles of charge exactly $j$. The terminal factor and the two reciprocal
Euler products are the residual free bosons. This is the classic form of fermionic sums for
parafermion-like theories (Lepowsky--Primc, Feigin--Stoyanovsky).
\end{physics}

\begin{status}
Definitions; local finiteness proved.
\end{status}

\subsection{The base algebra, arcs, and the presented arc algebra}

\begin{definition}[The base presentation]\label{l20:def:base}
$\mathrm{BasePoly}=\C[x_0,x_1,x_2]$ with $x_0=e$, $x_1=h$, $x_2=f$; $\ltwentyQ=h^2+4ef$; the base relations
are $x_i\ltwentyQ^m$ ($i=0,1,2$); $I_m$ is the ideal they generate and $\ltwentyA=\C[e,h,f]/I_m$.
\end{definition}
\noindent File \leanpath{VOA/\allowbreak Commutative/\allowbreak BoundaryAuxiliaryInitial.lean}.
\begin{leancode}
abbrev BasePoly (R : Type*) [CommRing R] := MvPolynomial (Fin 3) R
abbrev JetPoly (R : Type*) [CommRing R] := VOA.InfiniteJet.JetPoly R (Fin 3)

def casimir : BasePoly R := MvPolynomial.X 1 ^ 2 + 4 * MvPolynomial.X 0 * MvPolynomial.X 2

def baseRelation (m : ℕ) (i : Fin 3) : BasePoly R := MvPolynomial.X i * casimir ^ m

def baseIdeal (m : ℕ) : Ideal (BasePoly R) := Ideal.span (Set.range (baseRelation (R := R) m))
\end{leancode}

\begin{dictionary}
Namespace \lean{VOA.BoundaryAuxiliary}. \lean{MvPolynomial (Fin 3) R} is the polynomial ring over $R$ in
three variables \lean{X 0}, \lean{X 1}, \lean{X 2}; the file docstring fixes ``coordinates are
\lean{(e,h,f) = (0,1,2)}''. \lean{Set.range (baseRelation m)} is the three-element set
$\{e\ltwentyQ^m,h\ltwentyQ^m,f\ltwentyQ^m\}$ and \lean{Ideal.span} the ideal it generates. Headlines use
\lean{(R := ℂ)}. So \lean{BasePoly ℂ ⧸ baseIdeal m} is $\ltwentyA$. It has the same three generators (with the factors
written in the other order, $\ltwentyQ^m x_i$, and the same variable order $(e,h,f)=(0,1,2)$) as the ideal
\lean{AffineSl2BoundaryC2.boundaryIdeal m} of Chapter~\ref{ch:L09}, hence is the same ideal; the library records only the
inclusion it needs, \lean{BoundaryArcLiMap.baseIdeal_le_boundaryIdeal}.
\end{dictionary}

\begin{physics}
$e,h,f$ are the classical values of the three currents; $\ltwentyQ=h^2+4ef$ is the classical quadratic
Casimir (the classical limit of the Sugawara field, up to normalisation). The relations say that
$\ltwentyQ^m$ annihilates the whole adjoint representation: the Higgs branch is the nilpotent cone
$\ltwentyQ=0$, thickened to order $m$.
\end{physics}

\begin{status}
Definitions.
\end{status}

\begin{definition}[Jet polynomials, the arc ideal, the presented arc algebra]\label{l20:def:jet}
Let $\sigma$ be a set of generators. $\mathrm{JetPoly}(\sigma)=\C[x_{i,n}: i\in\sigma,n\ge0]$. For $f\in\C[x_i]$
put $f(\mathbf x(t))=f\bigl(\sum_nx_{i,n}t^n\bigr)$ and let $f_{[n]}$ be its coefficient of $t^n$. For an ideal
$I\subseteq\C[x_i]$, the \emph{arc ideal} $J_\infty I\subseteq\mathrm{JetPoly}$ is generated by all
$f_{[n]}$, $f\in I$, $n\ge0$, and the \emph{presented arc algebra} is $\mathrm{JetPoly}/J_\infty I$.
\end{definition}
\noindent File \leanpath{VOA/\allowbreak Commutative/\allowbreak InfiniteJet.lean}.
\begin{leancode}
abbrev JetPoly (R : Type uR) [CommSemiring R] (σ : Type uσ) :=
  MvPolynomial (σ × ℕ) R

noncomputable def relationSet (I : Ideal (MvPolynomial σ R)) : Set (JetPoly R σ) :=
  {f | ∃ (m : ℕ) (g : FiniteJet.JetPoly R σ m),
    g ∈ FiniteJet.jetIdeal (R := R) m I ∧ ofFinite (R := R) m g = f}

noncomputable def jetIdeal (I : Ideal (MvPolynomial σ R)) : Ideal (JetPoly R σ) :=
  Ideal.span (relationSet (R := R) I)

abbrev PresentedJet (I : Ideal (MvPolynomial σ R)) :=
  JetPoly R σ ⧸ jetIdeal (R := R) I

noncomputable def presentedJetMk (I : Ideal (MvPolynomial σ R)) :
    JetPoly R σ →ₐ[R] PresentedJet (R := R) I :=
  Ideal.Quotient.mkₐ R (jetIdeal (R := R) I)
\end{leancode}
\noindent File \leanpath{VOA/\allowbreak Commutative/\allowbreak InfiniteJetGenerators.lean}.
\begin{leancode}
def generatorSeries (i : σ) : PowerSeries (JetPoly R σ) :=
  PowerSeries.mk fun n => MvPolynomial.X (i, n)

def relationCoefficient (f : MvPolynomial σ R) (n : ℕ) : JetPoly R σ :=
  PowerSeries.coeff n (MvPolynomial.aeval generatorSeries f)

def coefficientIdeal (S : Set (MvPolynomial σ R)) : Ideal (JetPoly R σ) :=
  Ideal.span {p | ∃ f ∈ S, ∃ n : ℕ, relationCoefficient f n = p}

theorem jetIdeal_span_eq_coefficientIdeal (S : Set (MvPolynomial σ R)) :
    jetIdeal (Ideal.span S) = coefficientIdeal S := by
\end{leancode}
\noindent File \leanpath{VOA/\allowbreak Commutative/\allowbreak BoundaryAuxiliaryInitial.lean}.
\begin{leancode}
def sourceIdeal (m : ℕ) : Ideal (JetPoly R) :=
  VOA.InfiniteJet.coefficientIdeal (Set.range (baseRelation (R := R) m))

abbrev Source (m : ℕ) := JetPoly R ⧸ sourceIdeal (R := R) m

theorem sourceIdeal_eq_jetIdeal (m : ℕ) :
    sourceIdeal (R := R) m = VOA.InfiniteJet.jetIdeal (baseIdeal (R := R) m) :=
\end{leancode}

\begin{dictionary}
The variable \lean{MvPolynomial.X (i, n)} is $x_{i,n}$, the coefficient of $t^n$ of the generic arc
$x_i(t)$; so $x_{i,n}$ corresponds to $x_i^{(n)}$ of \eqref{l20:eq:canonical}, the \emph{divided} derivative
$\partial^nx_i/n!$ (the ``Hasse--Schmidt'' normalisation, which needs no division by $n!$ in the ring).
\lean{MvPolynomial.aeval generatorSeries f} substitutes $x_i\mapsto\sum_nx_{i,n}t^n$ into $f$, giving
$f(\mathbf x(t))\in\mathrm{JetPoly}[[t]]$; \lean{relationCoefficient f n} is $f_{[n]}$.

The \emph{definition} of the arc ideal (\lean{jetIdeal}) goes through finite jets: \lean{FiniteJet.jetIdeal m I}
is the order-$m$ jet ideal of Chapter~\ref{ch:L11} (Definition~\ref{l11:def:jets}), \lean{ofFinite m} is the
inclusion of order-$\le m$ jet variables into all jet variables, and \lean{jetIdeal I} is generated by the
images of all finite jet ideals. This avoids any colimit. The theorem
\lean{jetIdeal_span_eq_coefficientIdeal} proves that this agrees with the ``textbook'' definition: the arc
ideal of $(S)$ is generated by the coefficients $f_{[n]}$, $f\in S$, $n\ge0$ (multiples of the generators need
not be added: that is the content of the theorem). \lean{sourceIdeal m} is this coefficient ideal for the three
boundary relations and \lean{sourceIdeal_eq_jetIdeal} identifies it with \lean{jetIdeal (baseIdeal m)}. So
\lean{Source m} $=$ \lean{PresentedJet (baseIdeal m)} $=\ltwentyJ\ltwentyA=\ltwentyB$, two names for the same
quotient (the first is used in the Hilbert-series part, the second in the PVA part).

\textbf{Is it the arc algebra?} Theorem~\ref{l20:thm:arc} below proves the universal property: algebra maps
$\ltwentyJ\ltwentyA\to A$ correspond bijectively to algebra maps $\ltwentyA\to A[[t]]$. So it is the
coordinate ring of the arc scheme, not merely something that looks like it.
\end{dictionary}

\begin{physics}
$x_{i,n}$ is the classical value of the mode $x_{i,(-n-1)}$, i.e.\ of the $n$-th Taylor coefficient of the
field $x_i(z)$ at the origin; a state $x_{i_1,(-n_1-1)}\cdots x_{i_r,(-n_r-1)}\ltwentyvac$ becomes the monomial
$x_{i_1,n_1}\cdots x_{i_r,n_r}$. The arc ideal says that the constraints $f(x)=0$ on the target hold for the
whole field $x(z)$, i.e.\ for all Taylor coefficients of $f(x(z))$.
\end{physics}

\begin{status}
Definitions; \lean{jetIdeal_span_eq_coefficientIdeal} and \lean{sourceIdeal_eq_jetIdeal} proved, standard axioms.
\end{status}

\begin{theorem}[The arc universal property]\label{l20:thm:arc}
For every commutative $R$-algebra $A$, restriction along the universal arc (the tautological algebra map
$R[x]/I\to(\ltwentyJ(R[x]/I))[[t]]$, $x_i\mapsto\sum_nx_{i,n}t^n$) gives a bijection
$\operatorname{Hom}_{R\text{-alg}}(\ltwentyJ(R[x]/I),A)\cong\operatorname{Hom}_{R\text{-alg}}(R[x]/I,A[[t]])$.
If $e\colon R[x]/I\cong B$ is an isomorphism of algebras, composing with $e$ gives
$\operatorname{Hom}(\ltwentyJ(R[x]/I),A)\cong\operatorname{Hom}(B,A[[t]])$.
\end{theorem}
\noindent File \leanpath{VOA/\allowbreak Commutative/\allowbreak InfiniteJetArc.lean}.
\begin{leancode}
def seriesJetEval (x : σ → PowerSeries A) : JetPoly R σ →ₐ[R] A :=
  MvPolynomial.aeval fun v => PowerSeries.coeff v.2 (x v.1)

def toArc (I : Ideal (MvPolynomial σ R)) (f : PresentedJet I →ₐ[R] A) :
    (MvPolynomial σ R ⧸ I) →ₐ[R] PowerSeries A :=
  Ideal.Quotient.liftₐ I (MvPolynomial.aeval (pointSeries I f))
    (pointSeries_relations I f)

def arcEquiv (I : Ideal (MvPolynomial σ R)) :
    (PresentedJet I →ₐ[R] A) ≃ ((MvPolynomial σ R ⧸ I) →ₐ[R] PowerSeries A) where
  toFun := toArc I
  invFun := ofArc I
  left_inv := ofArc_toArc I
  right_inv := toArc_ofArc I

theorem existsUnique_arcLift (I : Ideal (MvPolynomial σ R))
    (u : (MvPolynomial σ R ⧸ I) →ₐ[R] PowerSeries A) :
    ∃! f : PresentedJet I →ₐ[R] A,
      (PowerSeries.mapAlgHom f).comp (universalArc I) = u := by
\end{leancode}
\noindent File \leanpath{VOA/\allowbreak Commutative/\allowbreak InfiniteJetBaseEmbedding.lean}.
\begin{leancode}
def arcEquivOfBaseAlgebraEquiv {A B : Type*} [CommRing A] [Algebra R A]
    [CommRing B] [Algebra R B] (I : Ideal (MvPolynomial σ R))
    (e : (MvPolynomial σ R ⧸ I) ≃ₐ[R] B) :
    (PresentedJet I →ₐ[R] A) ≃ (B →ₐ[R] PowerSeries A) where
  toFun f := (toArc I f).comp e.symm.toAlgHom
  invFun u := ofArc I (u.comp e.toAlgHom)
  left_inv f := by simp [AlgHom.comp_assoc]
  right_inv u := by simp [AlgHom.comp_assoc]

theorem baseToJet_injective (I : Ideal (MvPolynomial σ R)) : Function.Injective (baseToJet I) := by
\end{leancode}

\begin{dictionary}
\lean{α ≃ β} is a bijection with explicit inverse (Mathlib's \lean{Equiv}). \lean{pointSeries I f i} is the
power series $\sum_nf(x_{i,n})t^n$; \lean{toArc I f} sends $x_i$ to it; \lean{ofArc I u} sends $x_{i,n}$ to the
$t^n$-coefficient of $u(x_i)$. \lean{universalArc I} is the tautological arc $R[x]/I\to(\ltwentyJ)[[t]]$,
$x_i\mapsto\sum_nx_{i,n}t^n$; \lean{existsUnique_arcLift} is the universal property in its textbook form
(``$\exists!$''). The file also proves naturality in $A$ (\lean{toArc_natural}) and functoriality in the base
(\lean{arcMap}, \lean{arcMap_comp}). \lean{baseToJet I} is the inclusion $R[x]/I\to\ltwentyJ(R[x]/I)$ of
``constant arcs'' ($x_i\mapsto x_{i,0}$), with a retraction, hence injective. No characteristic-zero or
derivation hypothesis is used (Hasse--Schmidt normalisation).
\end{dictionary}

\textbf{Proof on paper.} An algebra map $\phi\colon\ltwentyJ\to A$ is determined by the images $a_{i,n}$ of
the $x_{i,n}$, subject to the vanishing of all $\phi(f_{[n]})$; and $\phi(f_{[n]})$ is the $t^n$-coefficient of
$f\bigl(\sum_na_{i,n}t^n\bigr)$. So $\phi$ is the same thing as a tuple of power series $a_i(t)$ with
$f(a(t))=0$ for $f\in I$, i.e.\ an algebra map $R[x]/I\to A[[t]]$. \qed

\begin{physics}
The arc space is the space of classical field configurations on the formal disc; ``points with values in
$A$'' are families of such configurations parametrised by $\operatorname{Spec}A$.
\end{physics}

\begin{status}
Proved, standard axioms. \lean{arcEquivOfBaseAlgebraEquiv} is a construction (\lean{def}); see the warning
after Theorem~\ref{l20:thm:c2arc}.
\end{status}

\subsection{Poisson vertex algebras}

\begin{definition}[Poisson vertex algebra, in mode form]\label{l20:def:pva}
A Poisson vertex algebra over $\C$ is a commutative $\C$-algebra $A$ with a derivation $\partial$ and bilinear
products $a_{(n)}b$, $n\in\N$, such that for all $a,b,c$:
\begin{enumerate}
\item (locality) $a_{(n)}b=0$ for $n\gg0$;
\item (sesquilinearity) $(\partial a)_{(0)}b=0$, $(\partial a)_{(n+1)}b=-(n+1)\,a_{(n)}b$, and
  $\partial(a_{(n)}b)=(\partial a)_{(n)}b+a_{(n)}(\partial b)$;
\item (Leibniz) $a_{(n)}(bc)=(a_{(n)}b)c+b(a_{(n)}c)$;
\item (skew-symmetry) $a_{(n)}b=\sum_{j\ge0}\frac{(-1)^{n+j+1}}{j!}\partial^j(b_{(n+j)}a)$;
\item (Jacobi) $a_{(m)}(b_{(n)}c)-b_{(n)}(a_{(m)}c)=\sum_{j=0}^m\binom mj(a_{(j)}b)_{(m+n-j)}c$.
\end{enumerate}
Equivalently, $\{a_\lambda b\}=\sum_n\frac{\lambda^n}{n!}a_{(n)}b$ is a $\lambda$-bracket satisfying the usual
PVA axioms. A homomorphism is an algebra map commuting with $\partial$ and every $(n)$-product; an
equivalence is a bijective homomorphism (whose inverse is then automatically a homomorphism).
\end{definition}
\noindent File \leanpath{VOA/\allowbreak Commutative/\allowbreak PoissonVertexAlgebra.lean}.
\begin{leancode}
structure PoissonVertexAlgebra (A : Type*) [CommRing A] [Algebra ℂ A] where
  translation : Derivation ℂ A A
  mode : ℕ → A →ₗ[ℂ] A →ₗ[ℂ] A
  locality : ∀ x y, ∃ B : ℕ, ∀ n, B ≤ n → mode n x y = 0
  translation_left_zero : ∀ x y, mode 0 (translation x) y = 0
  translation_left_succ : ∀ n x y,
    mode (n + 1) (translation x) y = (-((n + 1 : ℕ) : ℂ)) • mode n x y
  translation_mode : ∀ n x y,
    translation (mode n x y) = mode n (translation x) y + mode n x (translation y)
  leibniz : ∀ n x y z, mode n x (y * z) = mode n x y * z + y * mode n x z
  skew : ∀ n x y, mode n x y = ∑ᶠ j : ℕ,
    (j.factorial : ℂ)⁻¹ • (-1 : ℂ) ^ (n + j + 1) •
      ((translation.toLinearMap ^ j) (mode (n + j) y x))
  jacobi : ∀ m n x y z,
    mode m x (mode n y z) - mode n y (mode m x z) =
      ∑ j : Fin (m + 1), (m.choose j.val : ℂ) • mode (m + n - j.val) (mode j.val x y) z
\end{leancode}
\noindent Files \leanpath{VOA/\allowbreak Commutative/\allowbreak PoissonVertexHom.lean},
\leanpath{VOA/\allowbreak Commutative/\allowbreak PoissonVertexEquiv.lean}.
\begin{leancode}
structure Hom (P : VOA.PoissonVertexAlgebra A) (Q : VOA.PoissonVertexAlgebra B) where
  toAlgHom : A →ₐ[ℂ] B
  map_translation : ∀ x, toAlgHom (P.translation x) = Q.translation (toAlgHom x)
  map_mode : ∀ n x y, toAlgHom (P.mode n x y) = Q.mode n (toAlgHom x) (toAlgHom y)

structure Equiv (P : VOA.PoissonVertexAlgebra A) (Q : VOA.PoissonVertexAlgebra B) where
  toAlgEquiv : A ≃ₐ[ℂ] B
  map_translation : ∀ x, toAlgEquiv (P.translation x) = Q.translation (toAlgEquiv x)
  map_mode : ∀ n x y, toAlgEquiv (P.mode n x y) = Q.mode n (toAlgEquiv x) (toAlgEquiv y)

def Equiv.ofBijective {P : VOA.PoissonVertexAlgebra A} {Q : VOA.PoissonVertexAlgebra B}
    (f : Hom P Q) (hf : Function.Bijective f.toAlgHom) : Equiv P Q where
  toAlgEquiv := AlgEquiv.ofBijective f.toAlgHom hf
  map_translation := f.map_translation
  map_mode := f.map_mode

def Equiv.symm {P : VOA.PoissonVertexAlgebra A} {Q : VOA.PoissonVertexAlgebra B}
    (e : Equiv P Q) : Equiv Q P where
\end{leancode}

\begin{dictionary}
The structure is a \emph{bundle of data plus laws} on a fixed commutative $\C$-algebra \lean{A} (the
product is the ring multiplication of \lean{A}). \lean{Derivation ℂ A A} is a $\C$-linear Leibniz map
$\partial$. \lean{mode n} is a bilinear map (two nested \lean{→ₗ[ℂ]}): \lean{mode n x y} is $x_{(n)}y$.
\lean{translation.toLinearMap ^ j} is $\partial^j$. In \lean{skew}, \lean{∑ᶠ} is a \lean{finsum}; by
\lean{locality} only finitely many terms are non-zero, so it is an honest finite sum. Each field matches the
corresponding item of the definition; we checked them against the $\lambda$-bracket axioms:
$\{\partial a_\lambda b\}=-\lambda\{a_\lambda b\}$ gives \lean{translation_left_zero} and
\lean{translation_left_succ}; $\{a_\lambda\partial b\}=(\lambda+\partial)\{a_\lambda b\}$ is equivalent, given
the left law, to \lean{translation_mode}; skew-symmetry $\{b_\lambda a\}=-\{a_{-\lambda-\partial}b\}$ gives
\lean{skew} after expanding $(-\lambda-\partial)^{n+j}$; the Jacobi identity
$\{a_\lambda\{b_\mu c\}\}-\{b_\mu\{a_\lambda c\}\}=\{\{a_\lambda b\}_{\lambda+\mu}c\}$ gives \lean{jacobi} by
comparing coefficients of $\lambda^m\mu^n/m!n!$. The coefficient convention $x_{(n)}y=n!\cdot[\lambda^n]$ is
the one stated in the file docstring.

\lean{Hom} is a homomorphism. \lean{Equiv} contains an algebra isomorphism \lean{toAlgEquiv} that commutes with
translation and with every mode; \lean{Equiv.symm} proves that the inverse has the same properties, and
\lean{Equiv.ofBijective} turns a bijective homomorphism into an equivalence. So
\lean{PoissonVertexAlgebra.Equiv P Q} is exactly ``$P\cong Q$ as Poisson vertex algebras''.
\textbf{Could the notion be vacuous?} The structure has a model in every case used here
(Definitions~\ref{l20:def:lipva} and \ref{l20:def:arcpva}), and the axioms are the standard ones; an
\lean{Equiv} between two given PVAs is a strong statement because both structures are fixed independently.
\end{dictionary}

\begin{physics}
A Poisson vertex algebra is the classical limit of a vertex algebra: the commutative product is the regular
part of the OPE at leading order, the $\lambda$-bracket is the singular part
$a(z)b(w)\sim\sum_n\frac{(a_{(n)}b)(w)}{(z-w)^{n+1}}$. It is the algebraic structure of classical
two-dimensional field theory (Hamiltonian PDEs, classical $W$-algebras via Drinfeld--Sokolov).
\end{physics}

\begin{status}
Definitions. \lean{Equiv.symm}, \lean{Equiv.ofBijective} are constructions whose law-fields are proved;
standard axioms.
\end{status}

\begin{definition}[The Li Poisson vertex algebra]\label{l20:def:lipva}
For a VOA $W$, $\ltwentygr W=\bigoplus_{p,N}F^pW_N/F^{p+1}W_N$ with the product of Chapter~\ref{ch:L11}, the
derivation induced by $L_{-1}$, and the modes induced by $(n)$-products as in \eqref{l20:eq:grbracket}:
on classes of degrees $(p,N)$ and $(q,M)$ the $n$-th mode lands in degree $(p+q-n,N+M-n-1)$ (Li degree by \eqref{l20:eq:grbracket}; conformal weight because
$\operatorname{wt}(a_{(n)}b)=\operatorname{wt}a+\operatorname{wt}b-n-1$), and is set to $0$
when $n>p+q$ or $n+1>N+M$.
\end{definition}
\noindent File \leanpath{VOA/\allowbreak Filtered/\allowbreak LiGradedModes.lean}.
\begin{leancode}
noncomputable def liWeightLayerModeToGraded (W : VertexOperatorAlgebra V)
    (d e : ℕ × ℕ) (n : ℕ) :
    liWeightLayer W d.1 d.2 →ₗ[ℂ]
      liWeightLayer W e.1 e.2 →ₗ[ℂ] liAssociatedGraded W :=
  if hn : n ≤ d.1 + e.1 then
    if hN : n + 1 ≤ d.2 + e.2 then
      (liWeightLayerMode W d.1 e.1 d.2 e.2 n hn hN).compr₂
        (DirectSum.lof ℂ (ℕ × ℕ)
          (fun f : ℕ × ℕ => liWeightLayer W f.1 f.2)
          (d.1 + e.1 - n, d.2 + e.2 - (n + 1)))
    else 0
  else 0

noncomputable def liAssociatedGradedMode (W : VertexOperatorAlgebra V) (n : ℕ) :
    liAssociatedGraded W →ₗ[ℂ] liAssociatedGraded W →ₗ[ℂ] liAssociatedGraded W :=
  DirectSum.toModule ℂ (ℕ × ℕ) (liAssociatedGraded W →ₗ[ℂ] liAssociatedGraded W)
    fun d => (DirectSum.toModule ℂ (ℕ × ℕ)
      (liWeightLayer W d.1 d.2 →ₗ[ℂ] liAssociatedGraded W)
      fun e => (liWeightLayerModeToGraded W d e n).flip).flip
\end{leancode}
\noindent File \leanpath{VOA/\allowbreak Filtered/\allowbreak LiGradedModeTranslation.lean}.
\begin{leancode}
theorem liGradedClass_mode (W : VertexOperatorAlgebra V) (p q N M n : ℕ)
    (hn : n ≤ p + q) (hN : n + 1 ≤ N + M)
    (x y : V) (hx : x ∈ liWeightPiece W p N) (hy : y ∈ liWeightPiece W q M) :
    liAssociatedGradedMode W n (liGradedClass W p N x hx) (liGradedClass W q M y hy) =
      liGradedClass W (p + q - n) (N + M - (n + 1)) (W.smul x (n : ℤ) y)
        (liWeightPiece_smul_nonneg_mem W hn hN hx hy) := by
\end{leancode}
\noindent File \leanpath{VOA/\allowbreak Filtered/\allowbreak LiPoissonVertexAlgebra.lean}.
\begin{leancode}
noncomputable def liPoissonVertexAlgebra (W : VertexOperatorAlgebra V) :
    PoissonVertexAlgebra (liAssociatedGraded W) where
  translation := liAssociatedGradedDerivation W
  mode := liAssociatedGradedMode W
  locality := liAssociatedGradedMode_eventually_zero W
  translation_left_zero := liAssociatedGradedMode_translation_zero W
  translation_left_succ := liAssociatedGradedMode_translation_succ W
  translation_mode := liAssociatedGradedTranslation_mode W
  leibniz := liAssociatedGradedMode_leibniz W
\end{leancode}

\begin{dictionary}
\lean{liWeightPiece W p N} $=F^pW\cap W_N$ and \lean{liWeightLayer W p N} $=F^pW_N/F^{p+1}W_N$
(Chapter~\ref{ch:L11}, Definition~\ref{l11:def:layer}); \lean{liAssociatedGraded W} is their external
direct sum, a commutative $\C$-algebra (Chapter~\ref{ch:L11}, Theorem~\ref{l11:thm:ligr}).
\lean{liWeightLayerMode} descends $x\otimes y\mapsto x_{(n)}y$ to the layers (well defined by Li's estimate
$F^p_{(n)}F^q\subseteq F^{p+q-n}$ for $n\ge0$, Chapter~\ref{ch:L11}, Theorem~\ref{l11:thm:limult});
\lean{liWeightLayerModeToGraded} inserts the result in the component $(p+q-n,N+M-n-1)$ and is $0$ when that
index would be negative; \lean{DirectSum.toModule} extends bilinearly from homogeneous pieces to the whole
direct sum. \lean{liGradedClass W p N x hx} is the class of $x\in F^pW_N$ in the layer $(p,N)$, inserted into
the direct sum, and \lean{liGradedClass_mode} is the defining formula
$\bar x_{(n)}\bar y=\overline{x_{(n)}y}$, exactly \eqref{l20:eq:grbracket}. The ``zero when the index is
negative'' convention is correct: when $n>p+q$ the target $F^{p+q-n}/F^{p+q-n+1}$ is $W/W=0$, and when
$n+1>N+M$ the target weight is negative, hence $0$ in CFT type. \lean{liAssociatedGradedDerivation W} is the
descended $L_{-1}$ (shifting bidegree by $(1,1)$) packaged as a \lean{Derivation}. The remaining two fields
(\lean{skew}, \lean{jacobi}, not shown) are discharged by the theorems \lean{liAssociatedGradedMode_skew}
and \lean{liAssociatedGradedMode_jacobi}, proved from the VOA skew-symmetry and Borcherds commutator formula.

\textbf{Relation to Chapter~\ref{ch:L11}.} That chapter's dictionary (``What is not packaged'') records that Li's
vertex Poisson structure was not defined in the library, with an update note pointing here. It now is, for every VOA,
with all PVA laws proved: \lean{liPoissonVertexAlgebra W} has no hypotheses.
\end{dictionary}

\begin{physics}
This is the classical limit of the chiral algebra $W$: commutative fields with the leading singular OPE
coefficients as Poisson brackets. A mode $x_{(n)}$ with $n\ge0$ lowers the ``number of derivatives'' by $n$:
classical brackets of derivative fields produce $\delta$-function derivatives.
\end{physics}

\begin{status}
Construction with every law proved, for every VOA; standard axioms. It is the target of all headline
statements of Theorem~1.1.
\end{status}

\begin{definition}[The arc Poisson vertex algebra]\label{l20:def:arcpva}
On $\C[x_{i,n}]$ ($i\in\{e,h,f\}$) let $\partial x_{i,n}=(n+1)x_{i,n+1}$, and let the modes be the unique ones
making a PVA with $x_{i,0}\,{}_{(0)}\,x_{j,0}=\{x_i,x_j\}_{\ltwentysl}$ and $x_{i,0}\,{}_{(n+1)}\,x_{j,0}=0$. The arc
ideal $J_\infty I_m$ is stable under $\partial$ and under all modes in both slots, so the quotient
$\ltwentyB=\ltwentyJ\ltwentyA$ inherits a PVA, \lean{arcPVA m}.
\end{definition}
\noindent File \leanpath{VOA/\allowbreak Commutative/\allowbreak BoundaryArcHamiltonian.lean}.
\begin{leancode}
def currentPoissonBracket (i j : Fin 3) : BasePoly R :=
  ![![0, -2 * MvPolynomial.X 0, MvPolynomial.X 1],
    ![2 * MvPolynomial.X 0, 0, -2 * MvPolynomial.X 2],
    ![-MvPolynomial.X 1, 2 * MvPolynomial.X 2, 0]] i j

def currentHamiltonian (i : Fin 3) : Derivation R (BasePoly R) (BasePoly R) :=
  MvPolynomial.mkDerivation R (currentPoissonBracket (R := R) i)
\end{leancode}
\noindent File \leanpath{VOA/\allowbreak Commutative/\allowbreak InfiniteJetDifferential.lean}.
\begin{leancode}
def polynomialTranslation : Derivation R (JetPoly R σ) (JetPoly R σ) :=
  MvPolynomial.mkDerivation R (fun v : σ × ℕ =>
    (v.2 + 1) • MvPolynomial.X (v.1, v.2 + 1))
\end{leancode}
\noindent File \leanpath{VOA/\allowbreak Filtered/\allowbreak BoundaryArcPoisson.lean}.
\begin{leancode}
def polynomialPVA : PoissonVertexAlgebra (JetPoly ℂ) :=
  universalJetPVA 0 (by norm_num)

def arcPVA (m : ℕ) : PoissonVertexAlgebra (InfiniteJet.PresentedJet (baseIdeal (R := ℂ) m)) :=
  PoissonVertexAlgebra.quotient polynomialPVA (InfiniteJet.jetIdeal (baseIdeal m))
    (fun f hf => InfiniteJet.polynomialTranslation_stable (baseIdeal m) f hf)
    (polynomialPVA_mode_left m) (polynomialPVA_mode_right m)

theorem arcPVA_base_zero (m : ℕ) (i j : Fin 3) :
    (arcPVA m).mode 0 (InfiniteJet.generator (baseIdeal m) i 0)
      (InfiniteJet.generator (baseIdeal m) j 0) =
      MvPolynomial.aeval (fun a => InfiniteJet.generator (baseIdeal m) a 0)
        (currentPoissonBracket (R := ℂ) i j) := by

theorem arcPVA_unique (m : ℕ)
    (P : PoissonVertexAlgebra (InfiniteJet.PresentedJet (baseIdeal (R := ℂ) m)))
    (hT : P.translation = InfiniteJet.translation (baseIdeal m))
    (hzero : ∀ i j, P.mode 0 (InfiniteJet.generator (baseIdeal m) i 0)
      (InfiniteJet.generator (baseIdeal m) j 0) =
      MvPolynomial.aeval (fun a => InfiniteJet.generator (baseIdeal m) a 0)
        (currentPoissonBracket (R := ℂ) i j))
    (hpos : ∀ i j n, P.mode (n+1) (InfiniteJet.generator (baseIdeal m) i 0)
      (InfiniteJet.generator (baseIdeal m) j 0) = 0) : P = arcPVA m := by
\end{leancode}

\begin{dictionary}
\lean{currentPoissonBracket i j} is the matrix of $\{x_i,x_j\}$ in the order $(e,h,f)$:
$\{e,h\}=-2e$, $\{e,f\}=h$, $\{h,f\}=-2f$ and antisymmetry, i.e.\ the $\ltwentysl$ Lie bracket
$[h,e]=2e$, $[h,f]=-2f$, $[e,f]=h$. \lean{mkDerivation} extends a choice of values on variables to a
derivation. \lean{polynomialTranslation} is $\partial x_{i,n}=(n+1)x_{i,n+1}$ (Hasse--Schmidt convention: since
$x_{i,n}\leftrightarrow\partial^nx_i/n!$, one has $\partial(\partial^nx/n!)=(n+1)\,\partial^{n+1}x/(n+1)!$).
\lean{InfiniteJet.generator I i n} is the class of $x_{i,n}$ in \lean{PresentedJet I};
\lean{InfiniteJet.translation I} is the descended $\partial$.

\textbf{A surprising construction, and why it is harmless.} \lean{polynomialPVA} is not built from the
formulas above but by \emph{transport}: \lean{universalJetPVA 0 _} is \lean{liPoissonVertexAlgebra} of the
universal affine VOA at level $k=0$, carried to $\C[x_{i,n}]$ along the PBW isomorphism
$\C[x_{i,n}]\cong\ltwentygr V^0(\ltwentysl)$ (proved in \leanpath{VOA/\allowbreak Filtered/\allowbreak AffineSl2UniversalLiPBW.lean}
and \leanpath{AffineSl2UniversalJetPVA.lean}). The choice of level $0$ is arbitrary: in $\ltwentygr V^k$ the
level term $x_{(1)}y=k(x|y)\ltwentyvac$ has Li degree $0+0-1<0$ and dies, so the classical bracket does not
see $k$. The theorem \lean{arcPVA_unique} removes any doubt: \emph{any} PVA structure on
\lean{PresentedJet (baseIdeal m)} with the standard translation, with $\{x_{i,0}{}_\lambda x_{j,0}\}=\{x_i,x_j\}$
and no higher terms, \emph{equals} \lean{arcPVA m}. So \lean{arcPVA m} is the canonical arc PVA of the Poisson
algebra $\ltwentyA$, whatever auxiliary model was used to build it. \lean{PoissonVertexAlgebra.quotient}
(file \leanpath{VOA/\allowbreak Commutative/\allowbreak PoissonVertexQuotient.lean}) descends a PVA to a quotient by an
ideal stable under translation and under every mode in both slots; the stability of $J_\infty I_m$ is proved
from the $\ltwentysl$-invariance of $\ltwentyQ$ (\lean{currentHamiltonian_casimir}).
\end{dictionary}

\begin{physics}
$\ltwentyJ\ltwentyA$ with this structure is the classical current algebra $\{e_\lambda f\}=h$, etc.\ (the
Kirillov--Kostant bracket extended to fields), restricted to field configurations satisfying the classical
constraints $x\ltwentyQ^m=0$. Because $\ltwentyQ$ is a Casimir, the constraints are compatible with the
Hamiltonian flows generated by the currents: they are first-class.
\end{physics}

\begin{status}
Construction with all laws proved; \lean{arcPVA_unique} proved; standard axioms.
\end{status}

\subsection{The candidate singular vector and the canonical map}

The canonical map is built from the three current states in the simple quotient. The library packages the
data of a boundary singular vector in a structure \lean{Candidate} (Chapter~\ref{ch:L09}); the canonical map
only uses its level, and the vector enters the \emph{proofs} of the relations.

\begin{definition}[Candidate; the uniform candidate]\label{l20:def:candidate}
A \emph{candidate} at level $k$ and index $n\ge1$ is a non-zero vector $v\in V^k(\ltwentysl)$ of conformal weight
$2n+1$ with $e_0v=0$, $h_0v=2v$ and $x_{(r)}v=0$ for every current $x$ and every $r>0$ (an affine singular
vector of $\ltwentysl$-weight $2$), at the level $k=k_n$. The uniform candidate is one such vector, chosen for
every $n\ge1$ from a proved existence theorem.
\end{definition}
\noindent File \leanpath{VOA/\allowbreak Examples/\allowbreak AffineSl2BoundarySingularFamily.lean} (Chapter~\ref{ch:L09}).
\begin{leancode}
structure Candidate (k : ℂ) (n : ℕ) where
  indexPositive : 1 ≤ n
  isBoundaryLevel : k = level n
  noncritical : k + 2 ≠ 0
  vector : WakModule k
  homogeneous : vector ∈ gradingZA k (denominator n : ℤ)
  horizontalRaising : J k colE 0 vector = 0
  horizontalCharge : J k colH 0 vector = (2 : ℂ) • vector
  positiveSingular : SingP k vector
  nonzero : vector ≠ 0
\end{leancode}
\noindent File \leanpath{VOA/\allowbreak Examples/\allowbreak AffineSl2AdmissiblePrimitive.lean}.
\begin{leancode}
theorem boundary_primitive_exists (n : ℕ) :
    ∃ x : WakModule (level n), x ≠ 0 ∧
      x ∈ jointSpace _ (level_noncritical n) (denominator n : ℕ) 2 ∧
      J _ colE 0 x = 0 ∧ SingP _ x := by

noncomputable def uniformCandidate (n : ℕ) (hn : 1 ≤ n) : Candidate (level n) n := by
\end{leancode}

\begin{dictionary}
\lean{WakModule k} is the underlying space of $V^k(\ltwentysl)$; \lean{gradingZA k N} its weight-$N$ space;
\lean{J k a r} the current mode $a_r$ (colour \lean{colE}, \lean{colH}, \lean{colF} for $e,h,f$) of
Chapter~\ref{ch:L08}; \lean{SingP k v} says $a_rv=0$ for all colours $a$ and all $r>0$.
\lean{jointSpace _ _ N c} is the simultaneous eigenspace $L_0=N$, $h_0=c$ (\lean{mem_jointSpace_iff}).
\lean{uniformCandidate n hn} is defined by \lean{Classical.choose} from \lean{boundary_primitive_exists n}
(the proof body is omitted above; it fills the nine fields). It is therefore a non-constructive choice.
\textbf{This does not make the headline maps non-canonical}: the canonical map (Definition~\ref{l20:def:canonical})
uses only \lean{C.noncritical} (a proof, irrelevant) and the level, never \lean{C.vector}; the vector is used
only inside proofs that the relations hold. The argument \lean{hn : 1 ≤ n} is only used to fill
\lean{indexPositive}.
\end{dictionary}

\begin{physics}
$v$ is the first null state of the vacuum module: at the boundary level the Kac determinant first vanishes
at weight $p=2m+1$, in the $\ltwentysl$-triplet $\{v,f_0v,f_0^2v\}$. For $m=1$ it is the weight-3 null vector
of $L_{-4/3}(\ltwentysl)$ of Chapter~\ref{ch:L08}, Theorem~\ref{l08:thm:sing43}.
\end{physics}

\begin{status}
\lean{boundary_primitive_exists}: proved for every $n$, standard axioms (Section~\ref{l20:sec:m0}).
\lean{uniformCandidate}: a construction from it.
\end{status}

\begin{definition}[The canonical map]\label{l20:def:canonical}
Let $V=L_{k_m}(\ltwentysl)$ and let $\bar e,\bar h,\bar f\in V_1$ be the images of the three currents. The
canonical map $\pi\colon\ltwentyB=\C[x_{i,n}]/J_\infty I_m\to\ltwentygr V$ sends $x_{i,n}$ to
$\partial^n\overline{x_i}/n!$, where $\overline{x_i}\in F^0V_1/F^1V_1$.
\end{definition}
\noindent File \leanpath{VOA/\allowbreak Filtered/\allowbreak BoundaryArcLiMap.lean}.
\begin{leancode}
abbrev simpleVOA := QuotVOA.maxQuotientVOA (affineSl2VOA_closed k C.noncritical)

def gen (i : Fin 3) : WakModule k ⧸ QuotVOA.maxGradedIdeal (affineSl2VOA_closed k C.noncritical) :=
  (QuotVOA.maxGradedIdeal (affineSl2VOA_closed k C.noncritical)).mkQ
    (currentState k (symbolToCurrentColor i))

def polynomialMap : BoundaryAuxiliary.JetPoly ℂ →ₐ[ℂ] liAssociatedGraded (simpleVOA C) :=
  liInfiniteJetPolynomialMap (simpleVOA C) (gen C) (fun _ => 1) (gen_mem C)
\end{leancode}
\noindent Files \leanpath{VOA/\allowbreak Filtered/\allowbreak InfiniteLiJetBidegree.lean},
\leanpath{VOA/\allowbreak Filtered/\allowbreak ClassicalFreeness.lean}.
\begin{leancode}
noncomputable def liInfiniteJetPolynomialMap
    (W : VertexOperatorAlgebra V)
    (gen : Fin g → V) (wt : Fin g → ℕ)
    (hgen : ∀ i, gen i ∈ W.grading (wt i : ℤ)) :
    InfiniteJet.JetPoly ℂ (Fin g) →ₐ[ℂ] liAssociatedGraded W :=
  InfiniteJet.derivationJetEval (liAssociatedGradedTranslation W)
    (fun i => liDegreeZeroClass W (wt i) ⟨gen i, hgen i⟩)

noncomputable def dividedIterate (D : A →ₗ[ℂ] A) (r : ℕ) (a : A) : A :=
  ((r.factorial : ℂ)⁻¹) • ((D ^ r) a)
\end{leancode}
\noindent File \leanpath{VOA/\allowbreak Commutative/\allowbreak InfiniteJet.lean}.
\begin{leancode}
noncomputable def derivationJetEval (D : A →ₗ[ℂ] A) (x : σ → A) :
    JetPoly ℂ σ →ₐ[ℂ] A :=
  MvPolynomial.aeval fun ir ↦ FiniteJet.dividedIterate D ir.2 (x ir.1)
\end{leancode}
\noindent Files \leanpath{VOA/\allowbreak Filtered/\allowbreak BoundaryArcConformalSqueeze.lean},
\leanpath{VOA/\allowbreak Examples/\allowbreak AffineSl2BoundaryClassicalFreeness.lean}.
\begin{leancode}
def sourceMap : Source (R := ℂ) m →ₐ[ℂ] liAssociatedGraded W :=
  Ideal.Quotient.liftₐ (sourceIdeal (R := ℂ) m) (polynomialMap W gen hgen)
    (fun f hf => hI hf)

def canonicalMap (m : ℕ) (hm : 1 ≤ m) :
    Source (R := ℂ) m →ₐ[ℂ] liAssociatedGraded (AffineSl2BoundaryCharacter.vacuum m) :=
  BoundaryArcConformalSqueeze.sourceMap
    (BoundaryArcLiMap.simpleVOA (uniformCandidate m hm))
    (BoundaryArcLiMap.gen (uniformCandidate m hm))
    (BoundaryArcLiMap.gen_mem (uniformCandidate m hm)) m
    (BoundaryArcLiMap.sourceIdeal_le_kernel (uniformCandidate m hm))
\end{leancode}
\noindent File \leanpath{VOA/\allowbreak Filtered/\allowbreak BoundarySimpleArcPoisson.lean}.
\begin{leancode}
def Candidate.arcLiPoissonHom :
    PoissonVertexAlgebra.Hom (BoundaryAuxiliary.arcPVA m)
      (liPoissonVertexAlgebra (BoundaryArcLiMap.simpleVOA C)) :=
  boundaryArcLiPoissonHom k C.noncritical
    (QuotVOA.maxGradedIdeal (affineSl2VOA_closed k C.noncritical))
    (QuotVOA.maxGradedIdeal_isIdeal _) (QuotVOA.maxGradedIdeal_isGradedIdeal _)
    m (BoundaryArcLiMap.baseIdeal_le_kernel C)

theorem Candidate.arcLiPoissonHom_mk (f : BoundaryAuxiliary.JetPoly ℂ) :
    C.arcLiPoissonHom.toAlgHom (InfiniteJet.presentedJetMk (BoundaryAuxiliary.baseIdeal m) f) =
      BoundaryArcLiMap.polynomialMap C f := rfl
\end{leancode}

\begin{dictionary}
\lean{currentState k a} is the weight-one state $a_{-1}\ltwentyvac$; \lean{symbolToCurrentColor} converts the
paper order $(e,h,f)$ to the library's colour order; \lean{mkQ} is the projection to the simple quotient. So
\lean{gen C i} is $\overline{x_i}\in L_k(\ltwentysl)_1$. \lean{liDegreeZeroClass W 1 ⟨x, _⟩} is the class of
$x\in W_1=F^0W_1$ in the layer $(0,1)$. \lean{derivationJetEval D x} sends $x_{i,n}\mapsto D^n(x_i)/n!$ (the
\lean{dividedIterate}), with $D$ the descended $L_{-1}$; so \lean{polynomialMap C} is
$x_{i,n}\mapsto\partial^n\overline{x_i}/n!$, which is \eqref{l20:eq:canonical}. \lean{Ideal.Quotient.liftₐ} factors
an algebra map through a quotient, given that the ideal is in the kernel (\lean{sourceIdeal_le_kernel}, proved
in Section~\ref{l20:sec:m0}). So \lean{canonicalMap m hm} is $\pi\colon\ltwentyB\to\ltwentygr L_{k_m}$, with source
written as \lean{Source m}. Its target \lean{liAssociatedGraded (AffineSl2BoundaryCharacter.vacuum m)} is the
same as \lean{liAssociatedGraded (simpleVOA (uniformCandidate m hm))} by unfolding the \lean{abbrev}s (the two
non-criticality proofs are equal by proof irrelevance). \lean{arcLiPoissonHom} is the same map, on the
source written as \lean{PresentedJet (baseIdeal m)} and packaged as a PVA homomorphism; \lean{_mk} certifies
that it is literally \lean{polynomialMap} on representatives (\lean{rfl}).

\textbf{Is it the paper's map?} The paper's canonical map is
$\ltwentyJ R_V\to\ltwentygr V$ precomposed (implicitly) with $\ltwentyJ\ltwentyA\to\ltwentyJ R_V$. The Lean source is
$\ltwentyJ\ltwentyA$; the identification $\ltwentyA\cong R_V$ is the separate headline \lean{baseC2AlgEquiv}, and
the arc functor carries it to $\ltwentyJ\ltwentyA\cong\ltwentyJ R_V$ (\lean{actualC2ArcEquiv}). Theorem~\ref{l20:thm:ad11}
closes the loop inside Lean: it proves that the PVA isomorphism is the unique translation-compatible algebra map
extending the canonical $R_V\to\ltwentygr V$, which is the paper's own characterisation of its canonical map. See also
the warning in Section~\ref{l20:sec:headlines}.
\end{dictionary}

\begin{physics}
The canonical map assigns to each classical field monomial $\prod x_{i_r,n_r}$ the classical limit of the
normally ordered product $\prod(\partial^{n_r}x_{i_r}/n_r!)$ of currents. Surjectivity says the currents
strongly generate; injectivity (the theorem) says the only classical relations are the derivatives of
$x\ltwentyQ^m=0$.
\end{physics}

\begin{status}
Constructions; their well-definedness (\lean{sourceIdeal_le_kernel}, \lean{baseIdeal_le_kernel}) and the PVA
law-fields of \lean{arcLiPoissonHom} are proved for every $m\ge1$; standard axioms.
\end{status}

\subsection{Conformal pieces and the $C_2$-algebra}

\begin{definition}[Conformal weight on arcs; conformal pieces]\label{l20:def:conformal}
The conformal weight of $x_{i,n}$ is $n+1$. The weight-$N$ piece of $\ltwentyB$ is the image of the
weight-$N$ homogeneous polynomials. The weight-$N$ piece of $\ltwentygr W$ consists of the elements all of whose
components of bidegree $(p,M)$ with $M\neq N$ vanish.
\end{definition}
\noindent Files \leanpath{VOA/\allowbreak Commutative/\allowbreak BoundaryAuxiliaryGrading.lean},
\leanpath{VOA/\allowbreak Commutative/\allowbreak BoundaryAuxiliaryHilbert.lean}.
\begin{leancode}
def conformalWeight (v : Fin 3 × ℕ) : ℕ := v.2 + 1

def sourcePiece (m N : ℕ) : Submodule K (Source (R := K) m) :=
  (MvPolynomial.weightedHomogeneousSubmodule K conformalWeight N).map
    (Ideal.Quotient.mkₐ K (sourceIdeal (R := K) m)).toLinearMap
\end{leancode}
\noindent File \leanpath{VOA/\allowbreak Examples/\allowbreak AffineSl2BoundaryClassicalFreeness.lean}.
\begin{leancode}
def arcConformalPiece (m N : ℕ) :
    Submodule ℂ (InfiniteJet.PresentedJet (baseIdeal (R := ℂ) m)) :=
  (MvPolynomial.weightedHomogeneousSubmodule ℂ conformalWeight N).map
    (InfiniteJet.presentedJetMk (baseIdeal (R := ℂ) m)).toLinearMap
\end{leancode}
\noindent File \leanpath{VOA/\allowbreak Filtered/\allowbreak BoundaryArcConformalEquiv.lean}.
\begin{leancode}
def liConformalPiece (N : ℕ) : Submodule ℂ (liAssociatedGraded W) where
  carrier := {x | ∀ p M : ℕ, M ≠ N →
    DirectSum.component ℂ (ℕ × ℕ) (fun d => liWeightLayer W d.1 d.2) (p,M) x = 0}
\end{leancode}

\begin{dictionary}
\lean{weightedHomogeneousSubmodule K w N} is Mathlib's space of polynomials all of whose monomials
$\prod x_v^{d_v}$ have $\sum_vd_vw(v)=N$; with $w=$ \lean{conformalWeight}, $x_{i,n}$ has weight $n+1$, the
conformal weight of $\partial^nx_i/n!$ (currents have weight one). \lean{.map} is the image submodule.
\lean{sourcePiece} and \lean{arcConformalPiece} are the same construction for the two names of $\ltwentyB$
(\lean{sourceIdeal_eq_jetIdeal}). \lean{DirectSum.component … (p,M) x} is the $(p,M)$-component of $x$ in the
direct sum. (\lean{liConformalPiece} lives in namespace \lean{VOA.BoundaryArcConformalSqueeze}, with the section
variable \lean{W}; the omitted fields prove it is a subspace.) Since the relation ideal is homogeneous
(\lean{sourceIdeal_conformalHomogeneous}), $\ltwentyB=\bigoplus_N$\lean{sourcePiece m N}, each piece
finite-dimensional (\lean{finiteDimensional_sourcePiece}). The Lean convention agrees with the book's.
\end{dictionary}

\begin{physics}
$L_0$-eigenvalue: a monomial in field modes has energy equal to the sum of the mode weights. The grading
by $N$ is the one in which characters are computed.
\end{physics}

\begin{status}
Definitions.
\end{status}

The $C_2$-algebra \lean{A₂ W} $=W/C_2(W)$ is that of Chapter~\ref{ch:L12}, Definition~\ref{l12:def:c2}, with its
commutative ring structure, its $\C$-algebra structure (Chapter~\ref{ch:L16}) and the class map
\lean{mk₂ W}; nothing new is defined here.

%======================================================================
\section{The headline theorems}\label{l20:sec:headlines}
%======================================================================

The ten original declarations are in \leanpath{VOA/\allowbreak Examples/\allowbreak AffineSl2BoundaryMainTheorems.lean};
the paper-form Theorem~1.1, \lean{arakawaDai_theorem_1_1}, and its supporting definitions are in
\leanpath{VOA/\allowbreak Examples/\allowbreak AffineSl2BoundaryTheorem11.lean} (Section~\ref{l20:sec:ad11}, added 2026-10-28).
Their only hypothesis is \lean{hm : 1 ≤ m}; besides $m$ they take only the evident arguments of their statements (a weight
$N$ for \lean{coeff_product_le_character} and \lean{arcPoissonEquiv_conformal_iff}, an element $x$ for the latter, and a
coefficient algebra $A$ for \lean{actualC2ArcEquiv}). Their proofs are
one-line assemblies of the work in Sections~\ref{l20:sec:m0}--\ref{l20:sec:m8}: each \lean{...OfCharacter}
helper takes the character equality as an explicit argument, and the headline file discharges it with
\lean{character_eq_product}.

\subsection{Theorem 1.2: the character}

\begin{theorem}[The vacuum character of $L_{k_m}(\ltwentysl)$]\label{l20:thm:character}
For every $m\ge1$, $p=2m+1$:
\begin{equation*}
\sum_{N\ge0}\dim L_{k_m}(\ltwentysl)_N\,q^N=\frac{(q^p;q^p)_\infty^3}{(q)_\infty^3}=S_m(q)
=\frac{1}{(q)_\infty^2}\,T_m(q),
\end{equation*}
with $S_m$, $T_m$ as in Definition~\ref{l20:def:chain}; the last expression is the fermionic chain of
\eqref{l20:eq:fermionic}.
\end{theorem}
\noindent File \leanpath{VOA/\allowbreak Examples/\allowbreak AffineSl2BoundaryMainTheorems.lean}.
\begin{leancode}
theorem coeff_product_le_character (m : ℕ) (hm : 1 ≤ m) (N : ℕ) :
    PowerSeries.coeff N (product m) ≤ PowerSeries.coeff N (character m) := by

theorem character_eq_product (m : ℕ) (hm : 1 ≤ m) :
    character m = product m := by

theorem character_eq_chain (m : ℕ) (hm : 1 ≤ m) :
    character m = QSeries.hilbertBoundSeries m :=
  character_eq_chain_of_product m hm (character_eq_product m hm)

theorem character_eq_fermionic_chain (m : ℕ) (hm : 1 ≤ m) :
    character m = (QSeries.reciprocalEuler : PowerSeries ℚ)^2 * QSeries.fermionicTrace m := by
\end{leancode}

\begin{dictionary}
Namespace \lean{VOA.AffineSl2BoundaryCharacter}. The objects: \lean{character}, \lean{product}
(Definition~\ref{l20:def:character}); \lean{hilbertBoundSeries}, \lean{fermionicTrace}
(Definition~\ref{l20:def:chain}); \lean{reciprocalEuler} (Definition~\ref{l20:def:euler}); all are elements of
$\Q[[q]]$ and \lean{=} is equality of power series (all coefficients). The only hypothesis is $1\le m$, which is
the paper's $m\ge1$. \lean{coeff_product_le_character} is the lower bound of \eqref{l20:eq:squeeze} (Section~\ref{l20:sec:m6}):
a supporting lemma, half of the dimension squeeze, \emph{not} a theorem of the paper; it is listed among the ten headline
declarations of the Lean file only because it is the hard half. For Theorem~1.2 of the paper cite
\lean{character_eq_product} and \lean{character_eq_fermionic_chain};
\lean{character_eq_product} combines it with the upper bound \lean{coeff_character_le_product} (Section~\ref{l20:sec:m5})
by antisymmetry of $\le$ coefficientwise; \lean{character_eq_chain} rewrites with the $q$-series identity
$S_m=\text{product}$ (Section~\ref{l20:sec:m4}); \lean{character_eq_fermionic_chain} is the literal display of the paper's
Theorem~1.2: $\frac{1}{(q)^2_\infty}T_m$.

\textbf{Same, stronger, or weaker?} Same as the paper, coefficient by coefficient, for the actual simple VOA.
The docstring of \lean{character_eq_chain} calls \lean{hilbertBoundSeries} ``the literal descending-chain
expression''; strictly, the literal display of Theorem~1.2 is \lean{character_eq_fermionic_chain}
(\lean{hilbertBoundSeries} is the paper's intermediate $S_m$, with the residual sums not yet evaluated). Both
are proved, so nothing is lost.

\textbf{How could it have been wrong?} (a) If \lean{character} were defined via the product, the theorem
would be circular: it is defined from \lean{Module.finrank} of the actual weight spaces. (b) If \lean{vacuum}
were a quotient by a chosen ideal: it is the maximal-ideal quotient. (c) If the sum in \lean{fermionicTrace}
were infinite in some coefficient, \lean{finsum} would silently return $0$: local finiteness is proved, and our
independent numerical check (Example~\ref{l20:ex:numbers}) confirms the coefficients for $m=1,2$ up to $q^{14}$.
\end{dictionary}

\textbf{Proof on paper.} Sections~\ref{l20:sec:m5} and \ref{l20:sec:m6} prove
$[q^N]\,\mathrm{character}\le[q^N]\,\mathrm{product}\le[q^N]\,\mathrm{character}$ for all $N$; Section~\ref{l20:sec:m4}
proves $S_m=T_m/(q)_\infty^2=\mathrm{product}$. \qed

\begin{physics}
For the $(A_1,D_{2m+1})$ Argyres--Douglas theory the unflavoured Schur index (without $q^{-c/24}$) is
$\operatorname{PE}\bigl[3\frac{q-q^p}{(1-q)(1-q^p)}\bigr]=(q^p;q^p)^3_\infty/(q)^3_\infty$; physicists derived it from
the IR formula of Cordova--Shao (BPS particles of the Coulomb branch) and checked it against the chiral algebra.
The theorem is that statement on the VOA side, proved. The fermionic form is the quasi-particle expansion.
\end{physics}

\begin{status}
\textbf{Unconditional} (only $1\le m$). \texttt{\#print axioms} for all four:
\lean{[propext, Classical.choice, Quot.sound]}; no \lean{VOA.pbw} (maintained audit
\texttt{scripts/audit/ArakawaDaiAxiomAudit.lean}, PASS on 2026-10-13 on the working tree of that date,
aggregate build green, 13150 jobs).
\end{status}

\begin{example}[The numbers for $m=1,2$]\label{l20:ex:numbers}
This is an independent numerical check by the chapter author (Python, exact integer arithmetic), not a Lean
result. Coefficients of $q^0,\dots,q^{14}$:
\begin{center}\small
\begin{tabular}{l@{\ }l}
$1/(q)_\infty^3=\dim V^k_N$: & $1,3,9,22,51,108,221,429,810,1479,2640,4599,7868,13209,21843$\\
$m=1$ product $=T_1/(q)^2_\infty$: & $1,3,9,19,42,81,155,276,486,821,1368,2214,3541,5544,8586$\\
$m=2$ product $=T_2/(q)^2_\infty$: & $1,3,9,22,51,105,212,402,744,1326,2316,3936,6581,10779,17406$
\end{tabular}
\end{center}
For $m=1$ the difference $\dim V^k_N-\dim L_{k,N}$ is $0,0,0,3,9,27,66,153,324,658,\dots$: the first null
states are the triplet $\{v,f_0v,f_0^2v\}$ at weight $3$, and for $N\le8$ the difference is $3\dim V^k_{N-3}$
(the free descendants of the triplet), while at $N=9$ it is $3\cdot221-5\cdot1=658$: the first correction comes
from the second singular vector, at weight $p\cdot1\cdot3=9$ in a quintuplet. This is the Weyl correction of
Section~\ref{l20:sec:m6}: $(4s-1)$ and $(4s+1)$-dimensional $\ltwentysl$-multiplets at weights $ps(2s\mp1)$.
For $m=2$ the first difference is at weight $5$, as it must be.
\end{example}

\subsection{Theorem 1.1: classical freeness}

\begin{theorem}[Classical freeness of $L_{k_m}(\ltwentysl)$]\label{l20:thm:freeness}
For every $m\ge1$ the canonical map $\pi\colon\ltwentyJ\ltwentyA\to\ltwentygr L_{k_m}(\ltwentysl)$ is bijective, and is an
isomorphism of Poisson vertex algebras which, in both directions, maps the conformal weight $N$ piece onto the
conformal weight $N$ piece.
\end{theorem}
\noindent File \leanpath{VOA/\allowbreak Examples/\allowbreak AffineSl2BoundaryMainTheorems.lean}.
\begin{leancode}
theorem canonicalMap_bijective (m : ℕ) (hm : 1 ≤ m) :
    Function.Bijective (canonicalMap m hm) :=
  canonicalMap_bijective_of_character m hm (character_eq_product m hm)

theorem boundary_classically_free (m : ℕ) (hm : 1 ≤ m) :
    Function.Bijective (uniformCandidate m hm).arcLiPoissonHom.toAlgHom :=
  ⟨arcLiPoissonHom_injective_of_character m hm (character_eq_product m hm),
    (uniformCandidate m hm).arcLiPoissonHom_surjective⟩

def arcPoissonEquiv (m : ℕ) (hm : 1 ≤ m) :
    PoissonVertexAlgebra.Equiv (arcPVA m)
      (liPoissonVertexAlgebra (AffineSl2BoundaryCharacter.vacuum m)) :=
  arcPoissonEquivOfCharacter m hm (character_eq_product m hm)

theorem arcPoissonEquiv_conformal_iff (m : ℕ) (hm : 1 ≤ m) (N : ℕ)
    (x : InfiniteJet.PresentedJet (baseIdeal (R := ℂ) m)) :
    (arcPoissonEquiv m hm).toAlgEquiv x ∈
      BoundaryArcConformalSqueeze.liConformalPiece (AffineSl2BoundaryCharacter.vacuum m) N ↔
      x ∈ arcConformalPiece m N :=
  arcPoissonEquivOfCharacter_conformal_iff m hm (character_eq_product m hm) N x
\end{leancode}
\noindent File \leanpath{VOA/\allowbreak Examples/\allowbreak AffineSl2BoundaryClassicalFreeness.lean}.
\begin{leancode}
def arcPoissonEquivOfCharacter (m : ℕ) (hm : 1 ≤ m)
    (hcharacter : character m = product m) :
    PoissonVertexAlgebra.Equiv (arcPVA m) (liPoissonVertexAlgebra (AffineSl2BoundaryCharacter.vacuum m)) :=
  PoissonVertexAlgebra.Equiv.ofBijective (uniformCandidate m hm).arcLiPoissonHom
    ⟨arcLiPoissonHom_injective_of_character m hm hcharacter,
      (uniformCandidate m hm).arcLiPoissonHom_surjective⟩
\end{leancode}

\begin{dictionary}
Namespace \lean{VOA.AffineSl2BoundaryClassicalFreeness}. \lean{Function.Bijective f} is injective and
surjective. \lean{canonicalMap m hm} and \lean{(uniformCandidate m hm).arcLiPoissonHom.toAlgHom} are the
same map $\pi$ on the two names \lean{Source m} and \lean{PresentedJet (baseIdeal m)} of $\ltwentyB$
(Definition~\ref{l20:def:canonical}). \lean{arcPoissonEquiv m hm} is a term of type
\lean{PoissonVertexAlgebra.Equiv (arcPVA m) (liPoissonVertexAlgebra (vacuum m))}: an algebra isomorphism
that commutes with $\partial$ and all modes, whose inverse does too (Definition~\ref{l20:def:pva}). The type
alone already says ``these two PVAs are isomorphic''; the definition (shown: \lean{Equiv.ofBijective} of
\lean{arcLiPoissonHom}) says the isomorphism is the canonical map. \lean{arcPoissonEquiv_conformal_iff} says:
for every $N$ and every $x$, $\pi(x)$ lies in the weight-$N$ piece of $\ltwentygr V$ iff $x$ lies in the weight-$N$
piece of $\ltwentyB$. Since $\pi$ is bijective, this is ``$\pi$ and $\pi^{-1}$ both preserve conformal weight''
(the conformally graded part of the paper's statement).

\textbf{Hypotheses.} Only \lean{hm : 1 ≤ m}. The candidate is chosen internally; the map does not depend on
the choice (Definition~\ref{l20:def:candidate}).

\textbf{Same, stronger, or weaker?} The paper states the isomorphism $\ltwentyJ R_V\to\ltwentygr V$. These four
declarations state it for $\ltwentyJ\ltwentyA$, the arc algebra of the \emph{presented} algebra; together with
\lean{baseC2AlgEquiv} ($\ltwentyA\cong R_V$, below) and the functoriality of arcs this gives the paper's statement. The
single theorem \lean{arakawaDai_theorem_1_1} (Theorem~\ref{l20:thm:ad11}) packages exactly this: it equips the presented
arc algebra with a universal arc of $R_V$, proves the universal property, and proves that \lean{arcPoissonEquiv} is the
canonical map out of $\ltwentyJ R_V$ so realised. The library still has no $\ltwentyJ$ functor applied to the abstract
algebra \lean{A₂ W}; $\ltwentyJ R_V$ is the presented arc algebra plus universal arc, unique up to unique isomorphism.
\lean{canonicalMap_bijective} (source \lean{Source m})
and \lean{boundary_classically_free} (source \lean{PresentedJet (baseIdeal m)}) state the same fact on two isomorphic
presentations of $\ltwentyJ\ltwentyA$ (\lean{sourceIdeal_eq_jetIdeal}). See the warning below for exactly what is and is
not packaged.

\textbf{Wording.} The release audit (2026-10-27) recommended a two-part wording because, at the time, no single theorem had
the paper's shape. That recommendation is superseded by \lean{arakawaDai_theorem_1_1}: the announcement now cites one
theorem, worded as in Theorem~\ref{l20:thm:ad11}, with the caveat that $\ltwentyJ R_V$ is realised as a presented arc
algebra with a universal arc.
\end{dictionary}

\textbf{Proof on paper.} Surjectivity: $\ltwentygr V$ is generated as a differential algebra by the classes of the
currents (Li; Section~\ref{l20:sec:m0}). Injectivity: by the squeeze (Section~\ref{l20:sec:m7}), each conformal
piece of $\ltwentyB$ maps onto $\bigoplus_p\ltwentygr_{p,N}V$ whose total dimension is $\dim V_N$, and
$\dim(\ltwentyB)_N\le[q^N]\mathrm{product}=\dim V_N$; a surjection from a space of dimension at most that of
the target is injective; and a polynomial in the kernel has all its homogeneous components in the kernel. PVA
compatibility: both structures are determined by $\partial$ and the brackets of the generators, which match
(Section~\ref{l20:sec:m8}). \qed

\begin{physics}
The Argyres--Douglas chiral algebra is ``classically free'': its singular support is the arc space of the
(non-reduced) Higgs branch scheme $\operatorname{Spec}\C[e,h,f]/(x\ltwentyQ^m)$. In the language of the
4d/2d correspondence, every Schur operator is accounted for by a classical field configuration on the Higgs
branch; there are no ``extra'' quantum relations. This is a strong form of the Higgs branch / associated
variety correspondence (heuristic, physics literature: Beem--Rastelli).
\end{physics}

\begin{status}
\textbf{Unconditional} (only $1\le m$); \texttt{\#print axioms} of all four:
\lean{[propext, Classical.choice, Quot.sound]}, no \lean{VOA.pbw} (audit as above).
\end{status}

\begin{theorem}[The $C_2$-algebra and the arc universal property]\label{l20:thm:c2arc}
For $m\ge1$: (1) the map $\C[e,h,f]\to R_{L_{k_m}}$, $x_i\mapsto[\overline{x_i}]$, induces an isomorphism
$\ltwentyA=\C[e,h,f]/(e\ltwentyQ^m,h\ltwentyQ^m,f\ltwentyQ^m)\cong R_{L_{k_m}(\ltwentysl)}$; (2) for every commutative
$\C$-algebra $A$, algebra maps $\ltwentyJ\ltwentyA\to A$ correspond bijectively to algebra maps
$R_{L_{k_m}}\to A[[t]]$, via the specific bijection obtained from Theorem~\ref{l20:thm:arc} and (1); (3) this bijection is natural in $A$. So
$\ltwentyJ\ltwentyA$ represents the arc functor $A\mapsto\operatorname{Hom}(R_{L_{k_m}},A[[t]])$, i.e.\ it is the arc algebra
$\ltwentyJ R_{L_{k_m}}$. (Naturality is not part of the type of \lean{actualC2ArcEquiv}; it is the separate theorem
\lean{actualC2ArcEquiv_natural}; see the warning below.)
\end{theorem}
\noindent File \leanpath{VOA/\allowbreak Examples/\allowbreak AffineSl2BoundaryMainTheorems.lean}.
\begin{leancode}
def baseC2AlgEquiv (m : ℕ) (hm : 1 ≤ m) :
    (BasePoly ℂ ⧸ baseIdeal (R := ℂ) m) ≃ₐ[ℂ] A₂ (AffineSl2BoundaryCharacter.vacuum m) :=
  baseC2AlgEquivOfCharacter m hm (character_eq_product m hm)

def actualC2ArcEquiv {A : Type*} [CommRing A] [Algebra ℂ A]
    (m : ℕ) (hm : 1 ≤ m) :
    (InfiniteJet.PresentedJet (baseIdeal (R := ℂ) m) →ₐ[ℂ] A) ≃
      (A₂ (AffineSl2BoundaryCharacter.vacuum m) →ₐ[ℂ] PowerSeries A) :=
  actualC2ArcEquivOfCharacter m hm (character_eq_product m hm)

theorem actualC2ArcEquiv_natural {A B : Type*} [CommRing A] [Algebra ℂ A]
    [CommRing B] [Algebra ℂ B] (m : ℕ) (hm : 1 ≤ m)
    (f : InfiniteJet.PresentedJet (baseIdeal (R := ℂ) m) →ₐ[ℂ] A) (g : A →ₐ[ℂ] B) :
    actualC2ArcEquiv m hm (g.comp f) = (PowerSeries.mapAlgHom g).comp (actualC2ArcEquiv m hm f) :=
\end{leancode}
\noindent File \leanpath{VOA/\allowbreak Filtered/\allowbreak BoundarySimpleArcPoisson.lean}.
\begin{leancode}
def Candidate.baseC2AlgEquiv
    (hinj : Function.Injective C.arcLiPoissonHom.toAlgHom) :
    (BoundaryAuxiliary.BasePoly ℂ ⧸ BoundaryAuxiliary.baseIdeal (R := ℂ) m) ≃ₐ[ℂ]
      A₂ (BoundaryArcLiMap.simpleVOA C) :=
  baseC2AlgEquiv_of_arc_injective (BoundaryArcLiMap.simpleVOA C) (fun _ => 1)
    (fun i => ⟨BoundaryArcLiMap.gen C i,BoundaryArcLiMap.gen_mem C i⟩)
    (BoundaryAuxiliary.baseIdeal m) C.baseC2Relations (BoundaryArcLiMap.baseIdeal_le_kernel C)
    hinj C.baseC2Eval_surjective

def Candidate.actualC2ArcEquiv {A : Type*} [CommRing A] [Algebra ℂ A]
    (hinj : Function.Injective C.arcLiPoissonHom.toAlgHom) :
    (InfiniteJet.PresentedJet (BoundaryAuxiliary.baseIdeal (R := ℂ) m) →ₐ[ℂ] A) ≃
      (A₂ (BoundaryArcLiMap.simpleVOA C) →ₐ[ℂ] PowerSeries A) :=
  InfiniteJet.arcEquivOfBaseAlgebraEquiv (BoundaryAuxiliary.baseIdeal m) (C.baseC2AlgEquiv hinj)

theorem Candidate.actualC2ArcEquiv_natural {A B : Type*} [CommRing A] [Algebra ℂ A]
    [CommRing B] [Algebra ℂ B] (hinj : Function.Injective C.arcLiPoissonHom.toAlgHom)
    (f : InfiniteJet.PresentedJet (BoundaryAuxiliary.baseIdeal (R := ℂ) m) →ₐ[ℂ] A)
    (g : A →ₐ[ℂ] B) :
    C.actualC2ArcEquiv hinj (g.comp f) =
      (PowerSeries.mapAlgHom g).comp (C.actualC2ArcEquiv hinj f) :=

theorem Candidate.baseIdeal_eq_simpleSymbol_kernel
    (hinj : Function.Injective C.arcLiPoissonHom.toAlgHom) :
    BoundaryAuxiliary.baseIdeal (R := ℂ) m = RingHom.ker (simpleSymbol C) := by
\end{leancode}
\noindent File \leanpath{VOA/\allowbreak Filtered/\allowbreak ArcBaseExactness.lean}.
\begin{leancode}
def baseC2Map : (MvPolynomial σ ℂ ⧸ I) →ₐ[ℂ] A₂ W :=
  Ideal.Quotient.liftₐ I (MvPolynomial.aeval (fun i => mk₂ W (gen i).1)) (fun _ hf => hI hf)
\end{leancode}

\begin{dictionary}
\lean{≃ₐ[ℂ]} is an isomorphism of $\C$-algebras; \lean{A₂ (vacuum m)} is $R_V$ (Chapter~\ref{ch:L12}).
\lean{baseC2AlgEquiv} unfolds to \lean{Candidate.baseC2AlgEquiv}, which is \lean{AlgEquiv.ofBijective} of
\lean{baseC2Map}: the map $\ltwentyA\to R_V$ induced by $x_i\mapsto$ \lean{mk₂ W (gen i)}, the class of the $i$-th
current. So the isomorphism is the canonical one. This was originally visible only from the definitions, not from the
type; it is now a theorem: \lean{baseC2AlgEquiv_apply} identifies \lean{baseC2AlgEquiv} with the explicitly defined
\lean{canonicalC2Map}, and \lean{canonicalC2Map_mk_X} gives its value on generators (Section~\ref{l20:sec:ad11}). (The paper itself displays only the
surjection $\ltwentyA\twoheadrightarrow R_V$; the isomorphism is a corollary of Theorem~1.1.) \lean{simpleSymbol C} is the same map from \lean{SymSl2} $=\C[e,f,h]$ (Chapter~\ref{ch:L09});
\lean{baseIdeal_eq_simpleSymbol_kernel} states ``$\ker=I_m$'' literally, but it is stated for a candidate under the
hypothesis \lean{hinj}, which the headlines discharge (\lean{arcLiPoissonHom_injective_of_character}).

\lean{actualC2ArcEquiv} has type ``maps $\ltwentyJ\ltwentyA\to A$'' $\simeq$ ``maps $R_V\to A[[t]]$''; by definition it is
\lean{arcEquivOfBaseAlgebraEquiv} (Theorem~\ref{l20:thm:arc}) composed with \lean{baseC2AlgEquiv}.
\end{dictionary}

\begin{warning}[title={What the types say, and what only the definitions say}]
\begin{enumerate}
\item \textbf{The type of \lean{actualC2ArcEquiv} alone is weak; naturality is a separate theorem.} A term of type
  \lean{(X →ₐ[ℂ] A) ≃ (Y →ₐ[ℂ] PowerSeries A)} is just a bijection of sets for one fixed $A$; such bijections
  exist between many unrelated pairs of sets of the same cardinality. ``$X$ represents the arc functor of $Y$''
  needs the bijections to be \emph{natural in $A$}. That is now proved (2026-10-28): \lean{actualC2ArcEquiv_natural},
  printed above, says $\Phi_B(g\circ f)=g_*\circ\Phi_A(f)$ for every algebra map $g:A\to B$ (and likewise
  \lean{Candidate.actualC2ArcEquiv_natural} for the candidate form). So \lean{actualC2ArcEquiv} together with
  \lean{actualC2ArcEquiv_natural} \emph{is} a representability statement. The mathematically strong content is in
  \lean{baseC2AlgEquiv} (from which (2) and (3) follow formally). The bundled textbook form is
  \lean{c2UniversalArc_universal} (``every arc $R_V\to A[[t]]$ is uniquely the base change of the universal arc'';
  Theorem~\ref{l20:thm:ad11}, item~3).
\item \textbf{\lean{baseC2AlgEquiv} as a type says only $\ltwentyA\cong R_V$} (abstractly); that the isomorphism sends
  $x_i$ to the current classes was originally only by definition (\lean{baseC2Map}). An abstract (ungraded) isomorphism alone would
  give \emph{some} presentation of $R_V$ by three generators, but would say neither that the generators are the current
  classes nor anything about gradings or Hilbert series. \emph{Update (2026-10-28):} canonicity is now a theorem,
  \lean{baseC2AlgEquiv_apply} together with \lean{canonicalC2Map_mk_X} (Section~\ref{l20:sec:ad11}).
\item \textbf{The PVA source is the presented arc algebra.} \lean{arcPVA m} lives on
  \lean{PresentedJet (baseIdeal m)}; there is no PVA structure defined intrinsically on
  ``$\ltwentyJ$ of the abstract Poisson algebra \lean{A₂ (vacuum m)}''. The paper's $\ltwentyJ R_V$ is recovered by
  (1) and (2); and \lean{arcPVA_unique} shows the PVA structure on the presented arc algebra is the canonical
  one. This is the standard way to make $\ltwentyJ R$ concrete and is not a defect, but the reader should not
  expect a theorem whose source is literally ``\lean{J∞ (A₂ V)}''. Theorem~\ref{l20:thm:ad11} makes the identification
  explicit: the presented arc algebra, with the arc \lean{c2UniversalArc}, \emph{is} the arc algebra of $R_V$ by its
  universal property.
\end{enumerate}
\end{warning}

\textbf{Proof on paper.} The composite $\ltwentyA\hookrightarrow\ltwentyJ\ltwentyA\xrightarrow{\pi}\ltwentygr V$ (constant arcs,
then $\pi$) equals $\ltwentyA\to R_V=\ltwentygr_0V\hookrightarrow\ltwentygr V$. The first arrow is injective (it has the
retraction ``set all $x_{i,n}$, $n>0$, to zero''), $\pi$ is injective by Theorem~\ref{l20:thm:freeness}, so
$\ltwentyA\to R_V$ is injective; it is surjective because $R_V$ is generated by the current classes (strong
generation of $V$ by the currents). Then (2) is Theorem~\ref{l20:thm:arc} composed with (1). \qed

\begin{physics}
(1) is the exact chiral ring / Higgs branch \emph{scheme}: $R_V$ is not reduced for any $m\ge1$, since
$\ltwentyQ\neq0$ in $R_V$ (it has degree $2<2m+1$) although $\ltwentyQ^{m+1}=\ltwentyQ^m(h^2+4ef)=0$; for $m=1$ it is
$\C[e,h,f]/(\ltwentyQ e,\ltwentyQ h,\ltwentyQ f)$, whose reduction is $\C[\mathcal N]$ but which keeps the class of $\ltwentyQ$
itself in degree 2 (the image of the Sugawara vector).
\end{physics}

\begin{status}
\textbf{Unconditional} (only $1\le m$); axioms \lean{[propext, Classical.choice, Quot.sound]}, no \lean{VOA.pbw}.
For $m=1$ this settles, on paper, the exactness statement \lean{Sl2Pres.Sl2ExactSymbol} that Chapter~\ref{ch:L16}
records as open; \emph{no Lean theorem yet connects the two} (they are stated for differently named copies of
$L_{-4/3}(\ltwentysl)$), so the L16 status is unchanged until such a bridge is written.
\end{status}

\subsection{Theorem 1.1 in the paper's form: one theorem}\label{l20:sec:ad11}

The declarations above prove everything in Arakawa--Dai's Theorem~1.1, but spread over two presentations: the isomorphism
$\ltwentyA\cong R_V$ on one side, and the PVA isomorphism $\ltwentyJ\ltwentyA\cong\ltwentygr V$ on the other. A reader had to
assemble the paper's sentence ``the canonical map $\ltwentyJ R_V\to\ltwentygr V$ is an isomorphism'' by hand, and two facts
needed for that assembly were visible only from definitions: that \lean{baseC2AlgEquiv} is the canonical map, and that
$\ltwentyJ\ltwentyA$, transported along it, really is the arc algebra of $R_V$. The file
\leanpath{VOA/\allowbreak Examples/\allowbreak AffineSl2BoundaryTheorem11.lean} (added 2026-10-28) states all of this as theorems
and packages it into one statement, \lean{arakawaDai_theorem_1_1}. \emph{Nothing is re-proved there}: every mathematical input
is one of the theorems of the earlier sections; the file adds three definitions (the canonical map, the universal arc of $R_V$,
the zero-jet embedding) and the short arguments that connect them.

Throughout, $m\ge1$, $V=L_{k_m}(\ltwentysl)$ is \lean{AffineSl2BoundaryCharacter.vacuum m} (Definition~\ref{l20:def:simple}),
$R_V$ is \lean{A₂ V} (Chapter~\ref{ch:L12}), $\ltwentyA=\C[e,h,f]/I_m$ is \lean{BasePoly ℂ ⧸ baseIdeal m}
(Definition~\ref{l20:def:base}), and $\ltwentyB=\ltwentyJ\ltwentyA$ is \lean{InfiniteJet.PresentedJet (baseIdeal m)}
(Definition~\ref{l20:def:jet}), with generators $x_{i,n}$ (\lean{InfiniteJet.generator (baseIdeal m) i n}).

\begin{definition}[The canonical map, the universal arc of $R_V$, and the zero-jet embedding]\label{l20:def:canonc2}
Let $\bar e,\bar h,\bar f\in V_1$ be the images in $V$ of the three current states $e_{(-1)}\ltwentyvac$, $h_{(-1)}\ltwentyvac$,
$f_{(-1)}\ltwentyvac$, and write $(x_0,x_1,x_2)=(e,h,f)$.
\begin{enumerate}
\item The \emph{canonical map} $\kappa\colon\ltwentyA\to R_V$ is the algebra map with $\kappa(x_i)=[\bar x_i]$, the $C_2$-class
  of the $i$-th current. (It is well defined because the boundary relations $x_i\ltwentyQ^m$ vanish in $R_V$,
  Section~\ref{l20:sec:m0}.)
\item The \emph{universal arc of $R_V$} is the algebra map
  $\gamma\colon R_V\to\ltwentyB[[t]]$ obtained by composing $\kappa^{-1}\colon R_V\to\ltwentyA$ (Theorem~\ref{l20:thm:canonc2}
  below shows $\kappa$ is bijective) with the tautological arc $\ltwentyA\to\ltwentyB[[t]]$,
  $x_i\mapsto\sum_{n\ge0}x_{i,n}t^n$ of Theorem~\ref{l20:thm:arc}.
\item The \emph{zero-jet embedding} $\iota_0\colon R_V\to\ltwentyB$ is $\kappa^{-1}$ followed by the constant-arc inclusion
  $\ltwentyA\to\ltwentyB$, $x_i\mapsto x_{i,0}$; equivalently, the constant coefficient of $\gamma$.
\end{enumerate}
\end{definition}
\noindent File \leanpath{VOA/\allowbreak Examples/\allowbreak AffineSl2BoundaryTheorem11.lean}.
\begin{leancode}
def canonicalC2Map (m : ℕ) (hm : 1 ≤ m) :
    (BasePoly ℂ ⧸ baseIdeal (R := ℂ) m) →ₐ[ℂ] A₂ (AffineSl2BoundaryCharacter.vacuum m) :=
  baseC2Map (BoundaryArcLiMap.simpleVOA (uniformCandidate m hm)) (fun _ => 1)
    (fun i => ⟨BoundaryArcLiMap.gen (uniformCandidate m hm) i,
      BoundaryArcLiMap.gen_mem (uniformCandidate m hm) i⟩)
    (baseIdeal m) (uniformCandidate m hm).baseC2Relations

def c2UniversalArc (m : ℕ) (hm : 1 ≤ m) :
    A₂ (AffineSl2BoundaryCharacter.vacuum m) →ₐ[ℂ]
      PowerSeries (InfiniteJet.PresentedJet (baseIdeal (R := ℂ) m)) :=
  (InfiniteJet.universalArc (baseIdeal (R := ℂ) m)).comp (baseC2AlgEquiv m hm).symm.toAlgHom

def c2ToArc (m : ℕ) (hm : 1 ≤ m) :
    A₂ (AffineSl2BoundaryCharacter.vacuum m) →ₐ[ℂ]
      InfiniteJet.PresentedJet (baseIdeal (R := ℂ) m) :=
  (InfiniteJet.baseToJet (baseIdeal (R := ℂ) m)).comp (baseC2AlgEquiv m hm).symm.toAlgHom

theorem c2ToArc_eq_constantCoeff (m : ℕ) (hm : 1 ≤ m)
    (a : A₂ (AffineSl2BoundaryCharacter.vacuum m)) :
    c2ToArc m hm a = PowerSeries.constantCoeff (c2UniversalArc m hm a) := rfl
\end{leancode}

\begin{dictionary}
Namespace \lean{VOA.AffineSl2BoundaryClassicalFreeness}. \lean{baseC2Map W weight gen I hI} is the map of
\leanpath{VOA/\allowbreak Filtered/\allowbreak ArcBaseExactness.lean} quoted under Theorem~\ref{l20:thm:c2arc}: it factors
$x_i\mapsto$ \lean{mk₂ W (gen i)} through the quotient by \lean{I}, given the proof \lean{hI} that \lean{I} lies in the kernel. Here
\lean{W} is the simple quotient \lean{simpleVOA (uniformCandidate m hm)}, which unfolds to the same VOA as \lean{vacuum m}
(dictionary of Definition~\ref{l20:def:canonical}); the weights are all $1$ (currents have weight one); \lean{gen i} is the
image of the $i$-th current state, packaged with its membership in $V_1$ (\lean{gen_mem}); and the kernel condition is
\lean{baseC2Relations}, the relations of Section~\ref{l20:sec:m0}. So \lean{canonicalC2Map m hm} is $\kappa$, and it is now a
named definition with an explicit body rather than something read off from the definition of \lean{baseC2AlgEquiv}.

\lean{InfiniteJet.universalArc I} is the tautological arc $x_i\mapsto\sum_nx_{i,n}t^n$ and \lean{InfiniteJet.baseToJet I} the
constant-arc inclusion $x_i\mapsto x_{i,0}$ (dictionary of Theorem~\ref{l20:thm:arc}). \lean{(baseC2AlgEquiv m hm).symm.toAlgHom}
is $\kappa^{-1}$ as an algebra map; \lean{.comp} is composition (right to left). \lean{PowerSeries.constantCoeff} takes the
coefficient of $t^0$. \lean{c2ToArc_eq_constantCoeff} holds by \lean{rfl}: the constant coefficient of
$\sum_nx_{i,n}t^n$ is $x_{i,0}$ by definition, so the two constructions of $\iota_0$ agree definitionally.

\textbf{Could these be the wrong objects?} $\gamma$ and $\iota_0$ are built from \lean{baseC2AlgEquiv}, whose \emph{type}
alone would allow any isomorphism $\ltwentyA\cong R_V$. Theorem~\ref{l20:thm:canonc2} closes this: \lean{baseC2AlgEquiv}
is $\kappa$, by a theorem. So $\gamma$ is the arc $[\bar x_i]\mapsto\sum_nx_{i,n}t^n$ and $\iota_0$ is $[\bar x_i]\mapsto x_{i,0}$.
\end{dictionary}

\begin{physics}
$\kappa$ identifies the Higgs-branch chiral ring generators with the three currents. The universal arc $\gamma$ is the
``generic classical field configuration'' on the Higgs branch scheme: each chiral-ring generator $x_i$ becomes the field
$x_i(t)=\sum_nx_{i,n}t^n$ on the formal disc with independent Taylor coefficients, subject only to the constraints. The
zero-jet embedding remembers only the value of each field at the origin, $x_i(0)=x_{i,0}$.
\end{physics}

\begin{status}
Definitions; \lean{c2ToArc_eq_constantCoeff} by \lean{rfl}. \texttt{\#print axioms} of \lean{canonicalC2Map},
\lean{c2UniversalArc}, \lean{c2ToArc}, \lean{c2ToArc_eq_constantCoeff}: \lean{[propext, Classical.choice, Quot.sound]}
(run 2026-10-28; see the status of Theorem~\ref{l20:thm:ad11}).
\end{status}

\begin{theorem}[The presentation isomorphism is the canonical map]\label{l20:thm:canonc2}
For $m\ge1$: \lean{baseC2AlgEquiv} agrees with $\kappa$ on every element; $\kappa(x_i)$ is the $C_2$-class of the image in
$V$ of the current state $(x_i)_{(-1)}\ltwentyvac$; and $\kappa$ is bijective.
\end{theorem}
\noindent File \leanpath{VOA/\allowbreak Examples/\allowbreak AffineSl2BoundaryTheorem11.lean}.
\begin{leancode}
theorem baseC2AlgEquiv_apply (m : ℕ) (hm : 1 ≤ m) (x : BasePoly ℂ ⧸ baseIdeal (R := ℂ) m) :
    baseC2AlgEquiv m hm x = canonicalC2Map m hm x := rfl

theorem canonicalC2Map_mk_X (m : ℕ) (hm : 1 ≤ m) (i : Fin 3) :
    canonicalC2Map m hm (Ideal.Quotient.mk (baseIdeal (R := ℂ) m) (MvPolynomial.X i)) =
      mk₂ (AffineSl2BoundaryCharacter.vacuum m)
        ((QuotVOA.maxGradedIdeal (AffineVOA.affineSl2VOA_closed (level m)
          (level_noncritical m))).mkQ
          (AffineVOA.J (level m) (AffineSl2UniversalC2PBW.symbolToCurrentColor i) (-1)
            (AffineVOA.vacuumA (level m)))) := by

theorem canonicalC2Map_bijective (m : ℕ) (hm : 1 ≤ m) :
    Function.Bijective (canonicalC2Map m hm) :=
  (baseC2AlgEquiv m hm).bijective
\end{leancode}

\begin{dictionary}
\textbf{Decoding the generator formula.} \lean{Ideal.Quotient.mk (baseIdeal m) (MvPolynomial.X i)} is the class of the
variable $x_i$ in $\ltwentyA$. \lean{AffineVOA.vacuumA k} is the vacuum vector $\ltwentyvac$ of the universal affine VOA
$V^k(\ltwentysl)$, and \lean{AffineVOA.J k a n} is the mode $J^a_{(n)}$ of the $a$-th current acting on it
(Chapter~\ref{ch:L08}, Sections~\ref{l08:sec:currents} and~\ref{l08:sec:grading}); so
\lean{AffineVOA.J (level m) a (-1) (AffineVOA.vacuumA (level m))} is the weight-one state $J^a_{(-1)}\ltwentyvac$ (this is
\lean{currentState k a}, Definition~\ref{l20:def:canonical}). The library's current colours are $(e,f,h)=(0,1,2)$
(Chapter~\ref{ch:L08}), while the paper and \lean{BasePoly} use $(e,h,f)=(0,1,2)$; \lean{symbolToCurrentColor} is the
transposition of $1$ and $2$ that converts the second order into the first (\leanpath{AffineSl2UniversalC2PBW.lean}:
\lean{Equiv.swap 1 2}). \lean{(QuotVOA.maxGradedIdeal …).mkQ} is the projection $V^k\to V^k/\text{(maximal proper graded
ideal)}=L_k$ (Chapter~\ref{ch:L04}), and \lean{mk₂ (vacuum m)} is the class map $V\to R_V=V/C_2(V)$. So the right-hand
side is $[\bar x_i]$, with $x_0=e$, $x_1=h$, $x_2=f$: exactly the paper's ``$e,h,f\mapsto$ the classes of the currents''.
The file also states the three instances separately (\lean{baseC2AlgEquiv_mk_e}, \lean{baseC2AlgEquiv_mk_h},
\lean{baseC2AlgEquiv_mk_f}) and the same formula with the packaged generator (\lean{canonicalC2Map_mk_X_gen}).

\textbf{What changed.} \lean{baseC2AlgEquiv_apply} is a \emph{theorem} (proved by \lean{rfl}, i.e.\ by unfolding): the
abstract isomorphism \lean{baseC2AlgEquiv m hm}, applied to any $x$, gives $\kappa(x)$. Before, this was true but visible only by
unfolding definitions; the release audit (finding F4) asked for exactly this lemma. Together with \lean{canonicalC2Map_mk_X} it
pins the isomorphism completely, since an algebra map out of $\ltwentyA$ is determined by the images of $x_0,x_1,x_2$.
\lean{canonicalC2Map_bijective} restates the bijectivity of \lean{baseC2AlgEquiv} for $\kappa$. \textbf{Hypotheses:} only
\lean{hm : 1 ≤ m}. \textbf{Attribution:} the paper asserts only that $\kappa$ is surjective; bijectivity is a consequence of
Theorem~1.1 (Theorem~\ref{l20:thm:c2arc}).
\end{dictionary}

\textbf{Proof on paper.} \lean{baseC2AlgEquiv} is defined as $\kappa$ made into an isomorphism using its bijectivity, so its
underlying function is $\kappa$. On generators, $\kappa(x_i)=[\bar x_i]$ by the definition of the map induced by
$x_i\mapsto$\lean{mk₂ W (gen i)}, with \lean{gen i} the projected current state. Bijectivity is
Theorem~\ref{l20:thm:c2arc}(1). \qed

\begin{physics}
The Higgs branch chiral ring is $\C[e,h,f]/(x\ltwentyQ^m)$ \emph{with the moment-map operators as generators}: the
isomorphism is not an abstract coincidence of rings but the statement that the three conserved currents generate the
chiral ring subject to exactly these relations.
\end{physics}

\begin{status}
Proved, unconditional (only $1\le m$). \texttt{\#print axioms} of \lean{baseC2AlgEquiv_apply}, \lean{canonicalC2Map_mk_X},
\lean{canonicalC2Map_bijective}: \lean{[propext, Classical.choice, Quot.sound]} (2026-10-28).
\end{status}

\begin{theorem}[The universal property of the arc of $R_V$]\label{l20:thm:univarc}
For $m\ge1$, every commutative $\C$-algebra $A$ and every algebra map $u\colon R_V\to A[[t]]$ there is a unique algebra map
$f\colon\ltwentyB\to A$ such that applying $f$ coefficientwise to $\gamma$ gives $u$, i.e.\ $f_*\circ\gamma=u$. Hence the pair
$(\ltwentyB,\gamma)$ is the arc algebra $\ltwentyJ R_V$ of $R_V$ with its universal arc, unique up to unique isomorphism.
\end{theorem}
\noindent File \leanpath{VOA/\allowbreak Examples/\allowbreak AffineSl2BoundaryTheorem11.lean}.
\begin{leancode}
theorem c2UniversalArc_universal (m : ℕ) (hm : 1 ≤ m) {A : Type*} [CommRing A] [Algebra ℂ A]
    (u : A₂ (AffineSl2BoundaryCharacter.vacuum m) →ₐ[ℂ] PowerSeries A) :
    ∃! f : InfiniteJet.PresentedJet (baseIdeal (R := ℂ) m) →ₐ[ℂ] A,
      (PowerSeries.mapAlgHom f).comp (c2UniversalArc m hm) = u := by
\end{leancode}

\begin{dictionary}
\lean{{A : Type*} [CommRing A] [Algebra ℂ A]}: $A$ is any commutative $\C$-algebra, in any universe. \lean{∃! f, P f} is
``there is exactly one $f$ with $P f$''. \lean{PowerSeries.mapAlgHom f : PowerSeries B →ₐ[ℂ] PowerSeries A} applies $f$ to
every coefficient; so the equation says: for every $r\in R_V$, if $\gamma(r)=\sum_nc_nt^n$ then $u(r)=\sum_nf(c_n)t^n$. This
is the textbook universal property of the arc algebra: $\operatorname{Hom}_{\C\text{-alg}}(\ltwentyJ R,A)\cong
\operatorname{Hom}_{\C\text{-alg}}(R,A[[t]])$ via $f\mapsto f_*\circ\gamma$, stated as existence and uniqueness for each $u$.

\textbf{Why this needs no naturality clause.} Unlike \lean{actualC2ArcEquiv} (whose type is a bijection of sets for each
$A$, with naturality a separate theorem, warning after Theorem~\ref{l20:thm:c2arc}), this statement names the bijection: it is ``compose $f$ coefficientwise with the
fixed arc $\gamma$''. A bijection of this form is automatically natural in $A$, since $(g\circ f)_*\circ\gamma=g_*\circ(f_*\circ\gamma)$.
So this is a genuine representability statement: $\ltwentyB$ represents the functor $A\mapsto\operatorname{Hom}(R_V,A[[t]])$,
with universal element $\gamma$.

\textbf{Unique up to unique isomorphism.} If $(J',\gamma')$ is another pair with the same property, the universal properties
give $f\colon\ltwentyB\to J'$ with $f_*\gamma=\gamma'$ and $g\colon J'\to\ltwentyB$ with $g_*\gamma'=\gamma$; then
$(g\circ f)_*\gamma=\gamma=\mathrm{id}_*\gamma$, so $g\circ f=\mathrm{id}$ by uniqueness, and likewise $f\circ g=\mathrm{id}$. This
is the precise sense in which \lean{PresentedJet (baseIdeal m)} \emph{is} $\ltwentyJ R_V$, although the library has no functor
$\ltwentyJ$ on abstract algebras. (This uniqueness argument is on paper; it is not a separate Lean theorem.)

\textbf{Hypotheses:} only \lean{hm}. \textbf{Universe:} the standalone theorem allows any universe for $A$; the copy inside
\lean{arakawaDai_theorem_1_1} quantifies over \lean{A : Type} (universe $0$), which suffices for every commutative
$\C$-algebra one meets in practice and is all the paper's statement needs.
\end{dictionary}

\textbf{Proof on paper.} By definition $\gamma=\mathrm{arc}\circ\kappa^{-1}$, where $\mathrm{arc}\colon\ltwentyA\to\ltwentyB[[t]]$ is
the tautological arc. For a candidate $f$, the equation $f_*\circ\gamma=u$ is equivalent (precompose with $\kappa$, or with
$\kappa^{-1}$) to $f_*\circ\mathrm{arc}=u\circ\kappa$. By the universal property of the presented arc algebra
(\lean{existsUnique_arcLift}, Theorem~\ref{l20:thm:arc}) the latter has exactly one solution. \qed

\begin{physics}
Points of $\ltwentyJ R_V$ with values in $A$ are exactly the $A$-families of classical field configurations
$x_i(t)$ on the formal disc satisfying the Higgs-branch constraints. The theorem says $\ltwentyB$ is \emph{the} space of
such field configurations, not a model that merely resembles it.
\end{physics}

\begin{status}
Proved, unconditional (only $1\le m$). \texttt{\#print axioms} \lean{c2UniversalArc_universal}:
\lean{[propext, Classical.choice, Quot.sound]} (2026-10-28).
\end{status}

\begin{theorem}[The PVA isomorphism is the paper's canonical map]\label{l20:thm:pvacanon}
Let $c\colon R_V\to\ltwentygr V$ be the canonical inclusion of $R_V=F^0V/F^1V$ as the Li-degree-zero part of $\ltwentygr V$,
and $\pi$ the PVA isomorphism \lean{arcPoissonEquiv m hm} of Theorem~\ref{l20:thm:freeness}. Then:
\begin{enumerate}
\item two algebra maps out of a presented arc algebra that intertwine its translation with the same operator and agree on the
  order-zero generators $x_{i,0}$ are equal;
\item $\pi\circ\iota_0=c$;
\item $\pi$ is the unique algebra map $F\colon\ltwentyB\to\ltwentygr V$ with $F\circ\partial=\partial\circ F$ and $F\circ\iota_0=c$.
\end{enumerate}
\end{theorem}
\noindent File \leanpath{VOA/\allowbreak Examples/\allowbreak AffineSl2BoundaryTheorem11.lean}.
\begin{leancode}
theorem presentedJet_algHom_ext {σ B : Type*} [CommRing B] [Algebra ℂ B]
    (I : Ideal (MvPolynomial σ ℂ)) (D : B → B)
    (F G : InfiniteJet.PresentedJet I →ₐ[ℂ] B)
    (hF : ∀ x, F (InfiniteJet.translation I x) = D (F x))
    (hG : ∀ x, G (InfiniteJet.translation I x) = D (G x))
    (h0 : ∀ i, F (InfiniteJet.generator I i 0) = G (InfiniteJet.generator I i 0)) : F = G := by

theorem arcPoissonEquiv_c2ToArc (m : ℕ) (hm : 1 ≤ m)
    (a : A₂ (AffineSl2BoundaryCharacter.vacuum m)) :
    (arcPoissonEquiv m hm).toAlgEquiv (c2ToArc m hm a) =
      c2ToLiAlgebra (AffineSl2BoundaryCharacter.vacuum m) a := by

theorem eq_arcPoissonEquiv_of_c2ToArc (m : ℕ) (hm : 1 ≤ m)
    (F : InfiniteJet.PresentedJet (baseIdeal (R := ℂ) m) →ₐ[ℂ]
      liAssociatedGraded (AffineSl2BoundaryCharacter.vacuum m))
    (hT : ∀ x, F ((arcPVA m).translation x) =
      (liPoissonVertexAlgebra (AffineSl2BoundaryCharacter.vacuum m)).translation (F x))
    (hR : ∀ a, F (c2ToArc m hm a) = c2ToLiAlgebra (AffineSl2BoundaryCharacter.vacuum m) a) :
    F = (arcPoissonEquiv m hm).toAlgEquiv.toAlgHom := by
\end{leancode}

\begin{dictionary}
\textbf{The objects.} \lean{InfiniteJet.translation I} is the derivation $\partial x_{i,n}=(n+1)x_{i,n+1}$ descended to the
presented arc algebra (dictionary of Definition~\ref{l20:def:arcpva}); \lean{InfiniteJet.generator I i n} is the class of
$x_{i,n}$. In \lean{presentedJet_algHom_ext} the operator \lean{D : B → B} is an arbitrary function (no linearity is assumed of
it). \lean{(arcPVA m).translation} is the translation of the arc PVA; it equals \lean{InfiniteJet.translation (baseIdeal m)}
by the theorem \lean{arcPVA_translation} (\leanpath{VOA/\allowbreak Filtered/\allowbreak BoundaryArcPoisson.lean}).
\lean{(liPoissonVertexAlgebra V).translation} is the derivation of $\ltwentygr V$ induced by $L_{-1}$
(Definition~\ref{l20:def:lipva}). \lean{c2ToLiAlgebra V : A₂ V →ₐ[ℂ] liAssociatedGraded V} is $c$, the algebra inclusion of
$R_V=F^0V/F^1V$ as the degree-zero part of $\ltwentygr V$ (\leanpath{VOA/\allowbreak Filtered/\allowbreak C2LiAlgebra.lean};
its formula on polynomials in generators, \lean{c2ToLiAlgebra_aeval}, is quoted in Section~\ref{l20:sec:m8}).
\lean{(arcPoissonEquiv m hm).toAlgEquiv} is the underlying algebra isomorphism of the PVA isomorphism, and \lean{.toAlgHom}
forgets invertibility.

\textbf{What each statement says.} \lean{presentedJet_algHom_ext} is item~1, for every presented arc algebra (any generator
type \lean{σ}, any ideal \lean{I}). \lean{arcPoissonEquiv_c2ToArc} is item~2: $\pi(\iota_0(a))=c(a)$ for every $a\in R_V$.
\lean{eq_arcPoissonEquiv_of_c2ToArc} is item~3: any algebra map \lean{F} satisfying the two conditions \lean{hT} (intertwines
the translations) and \lean{hR} (restricts to $c$ on $R_V$) \emph{equals} $\pi$.

\textbf{Why this is the paper's canonical map.} Arakawa--Dai (following Arakawa's earlier work) define the canonical map
$\ltwentyJ R_V\to\ltwentygr V$ as the unique differential-algebra map extending $R_V=\ltwentygr_0V\hookrightarrow\ltwentygr V$;
it exists because $\ltwentyJ R_V$ is the free differential algebra on $R_V$ in the appropriate sense. Items~2 and~3 say exactly
that $\pi$ is this map. \textbf{Could item~3 be vacuous?} A uniqueness statement over an empty set of candidates says nothing.
It is not empty here: $\pi$ itself satisfies \lean{hT} (it is a PVA isomorphism, so it preserves translation; the field
\lean{map_translation}) and \lean{hR} (item~2). \textbf{Hypotheses:} only \lean{hm}.
\end{dictionary}

\textbf{Proof on paper.} (1) In characteristic zero $x_{i,n+1}=(n+1)^{-1}\partial x_{i,n}$ (the theorem \lean{generator_succ},
\leanpath{VOA/\allowbreak Commutative/\allowbreak InfiniteJetPoissonUniqueness.lean}). If $F$ and $G$ intertwine $\partial$ with
$D$ and agree on $x_{i,0}$, then by induction on $n$,
$F(x_{i,n+1})=(n+1)^{-1}D(F(x_{i,n}))=(n+1)^{-1}D(G(x_{i,n}))=G(x_{i,n+1})$, using only that $F$ and $G$ are $\C$-linear
(no property of $D$ is needed beyond being a function). So $F$
and $G$ agree on all generators $x_{i,n}$, hence everywhere, since these generate $\ltwentyB$ as an algebra. (2) Both
$\pi\circ\iota_0\circ\kappa$ and $c\circ\kappa$ are algebra maps out of $\ltwentyA$, so it suffices to compare them on $x_i$:
$\iota_0(\kappa(x_i))=x_{i,0}$, and $\pi(x_{i,0})$ is, by the definition of the canonical map
(Definition~\ref{l20:def:canonical}, the case $n=0$), the degree-zero Li class of $\bar x_i$; while $c([\bar x_i])$ is that
same class by definition of $c$. (3) Apply (1) to $F$ and $\pi$ with $D$ the translation of $\ltwentygr V$: both intertwine
(\lean{hT} for $F$, \lean{map_translation} for $\pi$), and on generators
$F(x_{i,0})=F(\iota_0(\kappa(x_i)))=c(\kappa(x_i))=\pi(x_{i,0})$ by \lean{hR} and~(2). \qed

\begin{physics}
The classical limit of a chiral algebra is determined by its chiral ring together with the derivative: once one knows how the
constant modes of the fields map to the chiral ring, the Taylor coefficients follow by differentiating. Item~3 is this
statement for the classical-freeness isomorphism: there is only one way to identify classical field configurations on the Higgs
branch with the associated graded of the chiral algebra compatibly with $L_{-1}$ and with the chiral ring.
\end{physics}

\begin{status}
Proved, unconditional (only $1\le m$; \lean{presentedJet_algHom_ext} has no hypothesis beyond its displayed ones).
\texttt{\#print axioms} of all three: \lean{[propext, Classical.choice, Quot.sound]} (2026-10-28).
\end{status}

\begin{theorem}[Arakawa--Dai, Theorem~1.1, in the paper's form]\label{l20:thm:ad11}
Let $m\ge1$, $k=-2+2/(2m+1)$, $V=L_k(\ltwentysl)$ (the quotient of $V^k(\ltwentysl)$ by its maximal proper graded ideal),
$R_V=V/C_2(V)$, $\ltwentyQ=h^2+4ef$. Realise $\ltwentyJ R_V$ as $\ltwentyB=\ltwentyJ\bigl(\C[e,h,f]/(e\ltwentyQ^m,h\ltwentyQ^m,
f\ltwentyQ^m)\bigr)$ with the arc $\gamma\colon R_V\to\ltwentyB[[t]]$ of Definition~\ref{l20:def:canonc2}. Then:
\begin{enumerate}
\item the canonical map $\kappa\colon\C[e,h,f]/(e\ltwentyQ^m,h\ltwentyQ^m,f\ltwentyQ^m)\to R_V$ is bijective;
\item $\kappa$ sends $e,h,f$ to the $C_2$-classes of the currents $e_{(-1)}\ltwentyvac$, $h_{(-1)}\ltwentyvac$, $f_{(-1)}\ltwentyvac$;
\item $(\ltwentyB,\gamma)$ has the universal property of the arc algebra of $R_V$: for every commutative $\C$-algebra $A$ and
  every algebra map $u\colon R_V\to A[[t]]$ there is a unique $f\colon\ltwentyB\to A$ with $f_*\circ\gamma=u$;
\item the Poisson vertex algebra isomorphism $\pi\colon\ltwentyB\to\ltwentygr V$ restricts on $R_V$ (via the zero-jet
  embedding $\iota_0$) to the canonical inclusion $c\colon R_V\to\ltwentygr V$;
\item $\pi$ is the unique algebra map $\ltwentyB\to\ltwentygr V$ that commutes with the translations and restricts to $c$; that is,
  $\pi$ is the canonical map $\ltwentyJ R_V\to\ltwentygr V$;
\item for every $N$, $\pi(x)$ has conformal weight $N$ if and only if $x$ has conformal weight $N$.
\end{enumerate}
Consequently, the canonical map $\ltwentyJ R_V\to\ltwentygr V$ is an isomorphism of conformally graded Poisson vertex algebras.
\end{theorem}
\noindent File \leanpath{VOA/\allowbreak Examples/\allowbreak AffineSl2BoundaryTheorem11.lean}.
\begin{leancode}
theorem arakawaDai_theorem_1_1 (m : ℕ) (hm : 1 ≤ m) :
    Function.Bijective (canonicalC2Map m hm) ∧
    (∀ i : Fin 3,
      canonicalC2Map m hm (Ideal.Quotient.mk (baseIdeal (R := ℂ) m) (MvPolynomial.X i)) =
        mk₂ (AffineSl2BoundaryCharacter.vacuum m)
          ((QuotVOA.maxGradedIdeal (AffineVOA.affineSl2VOA_closed (level m)
            (level_noncritical m))).mkQ
            (AffineVOA.J (level m) (AffineSl2UniversalC2PBW.symbolToCurrentColor i) (-1)
              (AffineVOA.vacuumA (level m))))) ∧
    (∀ (A : Type) [CommRing A] [Algebra ℂ A]
        (u : A₂ (AffineSl2BoundaryCharacter.vacuum m) →ₐ[ℂ] PowerSeries A),
      ∃! f : InfiniteJet.PresentedJet (baseIdeal (R := ℂ) m) →ₐ[ℂ] A,
        (PowerSeries.mapAlgHom f).comp (c2UniversalArc m hm) = u) ∧
    (∀ a, (arcPoissonEquiv m hm).toAlgEquiv (c2ToArc m hm a) =
      c2ToLiAlgebra (AffineSl2BoundaryCharacter.vacuum m) a) ∧
    (∀ F : InfiniteJet.PresentedJet (baseIdeal (R := ℂ) m) →ₐ[ℂ]
        liAssociatedGraded (AffineSl2BoundaryCharacter.vacuum m),
      (∀ x, F ((arcPVA m).translation x) =
        (liPoissonVertexAlgebra (AffineSl2BoundaryCharacter.vacuum m)).translation (F x)) →
      (∀ a, F (c2ToArc m hm a) = c2ToLiAlgebra (AffineSl2BoundaryCharacter.vacuum m) a) →
      F = (arcPoissonEquiv m hm).toAlgEquiv.toAlgHom) ∧
    (∀ (N : ℕ) (x : InfiniteJet.PresentedJet (baseIdeal (R := ℂ) m)),
      (arcPoissonEquiv m hm).toAlgEquiv x ∈
        BoundaryArcConformalSqueeze.liConformalPiece (AffineSl2BoundaryCharacter.vacuum m) N ↔
        x ∈ arcConformalPiece m N) :=
  ⟨canonicalC2Map_bijective m hm, canonicalC2Map_mk_X m hm,
    fun _ _ _ u => c2UniversalArc_universal m hm u, arcPoissonEquiv_c2ToArc m hm,
    eq_arcPoissonEquiv_of_c2ToArc m hm, arcPoissonEquiv_conformal_iff m hm⟩
\end{leancode}

\begin{dictionary}
\textbf{The six conjuncts} (joined by \lean{∧}), in order, are items~1--6, and the proof term (shown, since it is informative)
is the tuple of the six earlier theorems: \lean{canonicalC2Map_bijective} and \lean{canonicalC2Map_mk_X}
(Theorem~\ref{l20:thm:canonc2}), \lean{c2UniversalArc_universal} (Theorem~\ref{l20:thm:univarc}), \lean{arcPoissonEquiv_c2ToArc}
and \lean{eq_arcPoissonEquiv_of_c2ToArc} (Theorem~\ref{l20:thm:pvacanon}), and \lean{arcPoissonEquiv_conformal_iff}
(Theorem~\ref{l20:thm:freeness}). Every symbol has been decoded in those dictionaries: \lean{canonicalC2Map} is $\kappa$
(Definition~\ref{l20:def:canonc2}); the generator formula is decoded under Theorem~\ref{l20:thm:canonc2};
\lean{c2UniversalArc} is $\gamma$ and \lean{c2ToArc} is $\iota_0$; \lean{arcPoissonEquiv m hm} is the PVA isomorphism $\pi$
(\lean{PoissonVertexAlgebra.Equiv}: an algebra isomorphism preserving translation and every non-negative $\lambda$-bracket
coefficient, with its inverse doing the same; Definition~\ref{l20:def:pva}); \lean{arcPVA m} is the canonical arc PVA of
$\ltwentyA$ (Definition~\ref{l20:def:arcpva}, unique by \lean{arcPVA_unique}); \lean{liConformalPiece} and
\lean{arcConformalPiece} are the conformal-weight-$N$ pieces of Definition~\ref{l20:def:conformal}, $x_{i,n}$ having weight
$n+1$.

\textbf{Hypotheses.} Only \lean{hm : 1 ≤ m}. In the third conjunct $A$ ranges over \lean{Type} (universe $0$); the
standalone \lean{c2UniversalArc_universal} has every universe.

\textbf{Same, stronger, or weaker than the paper?} The paper's Theorem~1.1 states that $L_{-2+2/(2m+1)}(\ltwentysl)$ is
classically free, equivalently that $\ltwentyJ R_V\to\ltwentygr V$ is an isomorphism of conformally graded PVAs. Item~5 identifies $\pi$ with the
paper's canonical map (characterised as the unique translation-compatible extension of $R_V\hookrightarrow\ltwentygr V$), and
$\pi$ is a PVA isomorphism by its type, preserving the conformal grading by item~6. Items~1--2 are \emph{stronger} than the
paper's displayed claims: the paper asserts only the surjection $\kappa$ (citing earlier work). The one point where the Lean
statement is formulated differently is the meaning of ``$\ltwentyJ R_V$'': in Lean it is the presented arc algebra together with
the universal arc $\gamma$, and item~3 is what makes it $\ltwentyJ R_V$ (up to unique isomorphism, dictionary of
Theorem~\ref{l20:thm:univarc}). The PVA structure on it is \lean{arcPVA m}; it is the canonical one on $\ltwentyJ R_V$ because
\lean{arcPVA_unique} pins it by the brackets of the $x_{i,0}=\iota_0(\kappa(x_i))$, which are the Poisson brackets of $R_V$
transported along $\kappa$.

\textbf{How it could have been wrong or vacuous, and why it is not.} (a) An abstract isomorphism $\ltwentyA\cong R_V$ would say
nothing about generators: excluded by item~2. (b) An arbitrary algebra called ``$\ltwentyJ R_V$'' would make items~4--6
meaningless: excluded by item~3, which characterises $(\ltwentyB,\gamma)$ up to unique isomorphism. (c) A uniqueness clause over
an empty set: item~5's candidate set contains $\pi$ (items~4 and \lean{map_translation}). (d) A trivial VOA: $V$ is the actual
simple quotient, nonzero and simple (\lean{simple_vacuum_ne_zero}, \lean{simple_isSimple}; Definition~\ref{l20:def:simple}).
(e) The conformal clause on a trivial grading: both sides are the honest weight gradings, and $\ltwentyB$ is the direct sum of
its finite-dimensional weight pieces (Definition~\ref{l20:def:conformal}).
\end{dictionary}

\textbf{Proof on paper.} Collect the six earlier theorems, as the proof term does. The mathematical content lies in
Theorem~\ref{l20:thm:freeness} (the Hilbert-series squeeze and the PVA compatibility) and Theorem~\ref{l20:thm:c2arc}
(injectivity of $\kappa$ from injectivity of $\pi$); Theorems~\ref{l20:thm:canonc2}--\ref{l20:thm:pvacanon} are the
bookkeeping that identifies the objects with the paper's. \qed

\begin{physics}
For the Argyres--Douglas theories $(A_1,D_{2m+1})$, whose chiral algebra is $L_{k_m}(\ltwentysl)$, the theorem says in one
sentence: the chiral algebra, viewed classically, is exactly the algebra of field configurations on the (non-reduced) Higgs
branch scheme $\operatorname{Spec}\C[e,h,f]/(x\ltwentyQ^m)$, with the moment maps $e,h,f$ given by the $\ltwentysl$ currents
and the Poisson bracket given by the current algebra. The Schur index therefore equals the Hilbert series of the arc space of
the Higgs branch scheme (Theorem~\ref{l20:thm:character}). (Heuristic dictionary; the statement is pure algebra.)
\end{physics}

\begin{status}
\textbf{Proved, unconditional} (only $1\le m$). \texttt{\#print axioms} \lean{arakawaDai_theorem_1_1}:
\lean{[propext, Classical.choice, Quot.sound]}; the same for the eleven supporting declarations of this section
(Section~\ref{l20:sec:verification}, 2026-10-28). \emph{Machine-record caveat:} the compiled module of
\leanpath{VOA/\allowbreak Examples/\allowbreak AffineSl2BoundaryTheorem11.lean} was not present in the main checkout's
\texttt{.lake} at the time of writing; the axiom check was run against the compiled modules of the \texttt{release-cheap}
worktree, whose copy of this file is byte-identical, and the aggregate \texttt{lake build} was not re-run. The project record
(\texttt{Notes/RELEASE\_PLAN.md}, ``Caveats removed in Lean'', 2026-10-28) reports the targeted umbrella builds green and the
axioms standard.
\end{status}

\begin{warning}[title={The remaining caveat: $\ltwentyJ R_V$ is realised, not constructed by a functor}]
The library has no functor $\ltwentyJ$ on abstract commutative algebras, so there is no Lean term ``\lean{J∞ (A₂ V)}''. In
\lean{arakawaDai_theorem_1_1}, $\ltwentyJ R_V$ is the presented arc algebra \lean{PresentedJet (baseIdeal m)} with the universal
arc \lean{c2UniversalArc m hm}, transported along the canonical isomorphism $\kappa$; it is characterised up to unique
isomorphism by the proved universal property (item~3). This is how the arc algebra of a finitely presented algebra is
constructed in practice, and nothing in the mathematics depends on the choice of presentation; but a reader who expects the
source of the map to be literally ``$\ltwentyJ$ applied to \lean{A₂ V}'' will not find it. Two further attributions to keep:
the paper itself asserts only the surjection $\ltwentyA\twoheadrightarrow R_V$ (bijectivity is a consequence of Theorem~1.1);
and the older declaration \lean{actualC2ArcEquiv} is a representability statement only together with its naturality
theorem \lean{actualC2ArcEquiv_natural} (the bundled form is \lean{c2UniversalArc_universal}).
\end{warning}


%======================================================================
\section{The boundary relations hold, and the currents generate}\label{l20:sec:m0}
%======================================================================

This section proves that $e\ltwentyQ^m,h\ltwentyQ^m,f\ltwentyQ^m$ vanish in $R_V$ (and hence their arc coefficients vanish in
$\ltwentygr V$), and that $\pi$ is surjective. In the paper this is an input from the literature; here it is proved.

\subsection{A null vector at weight $p$ from the determinant}

\begin{theorem}[First admissible null vector]\label{l20:thm:firstnull}
Let $k=p'/q-2$ be admissible ($p'\ge2$, $q\ge1$ coprime). Then the maximal ideal of $V^k(\ltwentysl)$ contains a
non-zero vector of conformal weight $q(p'-1)$ and $h_0$-eigenvalue $2(p'-1)$. Such a vector is annihilated by
every positive current mode and by $e_0$. At the boundary ($p'=2$, $q=2m+1$) it has weight $2m+1$ and
$h_0$-charge $2$.
\end{theorem}
\noindent File \leanpath{VOA/\allowbreak Examples/\allowbreak AffineSl2AdmissibleNullVector.lean}.
\begin{leancode}
theorem admissible_first_null (p q : ℕ) (hp : 2 ≤ p) (hq : 0 < q)
    (hcop : p.Coprime q) :
    ∃ x : WakModule ((p : ℂ) / q - 2), x ≠ 0 ∧
      x ∈ jointSpace _ (admissible_level_noncritical p q hp hq)
        (q * (p - 1) : ℕ) (2 * ((p : ℤ) - 1)) ∧
      x ∈ QuotVOA.maxGradedIdeal
        (universal _ (admissible_level_noncritical p q hp hq)) := by
\end{leancode}
\noindent File \leanpath{VOA/\allowbreak Examples/\allowbreak AffineSl2AdmissiblePrimitive.lean}.
\begin{leancode}
theorem admissible_first_null_positive (p q : ℕ) (hp : 2 ≤ p) (hq : 0 < q)
    (hcop : p.Coprime q)
    (x : WakModule ((p : ℂ) / q - 2))
    (hx : x ∈ jointSpace _ (admissible_level_noncritical p q hp hq)
      (q * (p - 1) : ℕ) (2 * ((p : ℤ) - 1)))
    (hi : x ∈ QuotVOA.maxGradedIdeal
      (universal _ (admissible_level_noncritical p q hp hq))) : SingP _ x := by
\end{leancode}

\begin{dictionary}
Here \lean{p} is the numerator of $k+2$ (Chapter~\ref{ch:L09}'s convention), not the chapter's $p$. Annihilation by
$e_0$ is the companion theorem \lean{admissible_first_null_raising} (same file, same hypotheses, conclusion
\lean{J _ colE 0 x = 0}). The
boundary specialisation \lean{boundary_primitive_exists} (Definition~\ref{l20:def:candidate}) is the case
\lean{p = 2}, \lean{q = denominator n}. ``In the maximal ideal'' means ``null'': orthogonal to everything under
the invariant form (Chapter~\ref{ch:L04}, Section~\ref{l04:sec:gram}).
\end{dictionary}

\textbf{Proof on paper.} The weight-$N$, charge-$c$ block $G_{N,c}(k)$ of the Shapovalov (invariant) form on
$V^k(\ltwentysl)$, written in the PBW basis, is a matrix of polynomials in $k$; its determinant factorises
(Kac--Kazhdan, via the reflected-Verma cancellation of Section~\ref{l20:sec:m6}) into linear factors in $k$,
and coprimality of $p'$ and $q$ shows that the first factor vanishing at $k=p'/q-2$ occurs at weight
$q(p'-1)$, charge $2(p'-1)$, and none below. A vanishing determinant gives a non-zero radical vector, i.e.\ a
vector of the maximal ideal (Chapter~\ref{ch:L04}). Applying a positive mode $a_r$ ($r>0$) gives an element of
the maximal ideal of lower weight, i.e.\ in the range of weights below $q(p'-1)$, where the determinant is non-zero, so it is $0$.
Applying $e_0$ gives charge $2p'$ at the same weight, where again the determinant is non-zero. \qed

\begin{physics}
At the boundary level the first null state of the vacuum module sits at $L_0=p=2m+1$ in the adjoint
($\ltwentysl$-spin one) representation; in the AD theory it is the first relation in the chiral ring. Its
existence is ``where the Kac determinant first vanishes''.
\end{physics}

\begin{status}
Proved for every admissible level; standard axioms. This is a genuine improvement on Chapter~\ref{ch:L09},
where the boundary candidate was available only for $m=1,2$.
\end{status}

\subsection{The symbol of the null vector, and the relations}

\begin{theorem}[The $C_2$ symbol is a non-zero multiple of $e\ltwentyQ^m$]\label{l20:thm:symbol}
Let $v$ be a candidate (Definition~\ref{l20:def:candidate}). Then $v$ lies in the maximal ideal, its class in
$R_{V^k}=\C[e,h,f]$ is $c\cdot e\ltwentyQ^m$ for some $c\neq0$, and consequently
$(e\ltwentyQ^m,h\ltwentyQ^m,f\ltwentyQ^m)\subseteq\ker\bigl(\C[e,h,f]\to R_{L_k}\bigr)$.
\end{theorem}
\noindent Files \leanpath{VOA/\allowbreak Examples/\allowbreak AffineSl2BoundaryRelationBridge.lean},
\leanpath{VOA/\allowbreak Examples/\allowbreak AffineSl2BoundarySymbolNonzero.lean}.
\begin{leancode}
theorem Candidate.mem_maxGradedIdeal {k : ℂ} {n : ℕ} (C : Candidate k n) :
    C.vector ∈ QuotVOA.maxGradedIdeal (affineSl2VOA_closed k C.noncritical) := by

theorem charge_two_highest_symbol_ne_zero (k : ℂ) (hk : k+2 ≠ 0)
    (x : WakModule k) (hne : x ≠ 0)
    (hE : J k colE 0 x = 0) (hH : J k colH 0 x = (2 : ℂ) • x)
    (hpos : SingP k x) :
    mk₂ (affineSl2VOA_closed k hk) x ≠ 0 := by

theorem Candidate.exists_highestSymbol {k : ℂ} {n : ℕ} (C : Candidate k n) :
    ∃ c : ℂ, c ≠ 0 ∧ mk₂ (affineSl2VOA_closed k C.noncritical) C.vector =
      c • symbolAlg k C.noncritical
        (eSym * casSym ^ n) :=

theorem Candidate.boundaryIdeal_le_kernel {k : ℂ} {n : ℕ} (C : Candidate k n) :
    AffineSl2BoundaryC2.boundaryIdeal n ≤
      RingHom.ker (simpleSymbol C) :=
\end{leancode}

\begin{dictionary}
\lean{mk₂ W x} is the class of $x$ in $R_W$. \lean{symbolAlg k hk : SymSl2 →ₐ[ℂ] A₂ (V^k)} is the evaluation
$\C[e,h,f]\to R_{V^k}$, an isomorphism by the PBW theorem (Chapter~\ref{ch:L09}); \lean{eSym}, \lean{casSym}
are $e$, $\ltwentyQ$. \lean{simpleSymbol C} is the composite $\C[e,h,f]\to R_{V^k}\to R_{L_k}$. The hypotheses of
\lean{charge_two_highest_symbol_ne_zero} are exactly ``non-zero affine highest-weight vector of $\ltwentysl$-weight
two''; note that the theorem holds at \emph{every} non-critical level, it is not specific to the boundary.
\end{dictionary}

\textbf{Proof on paper.} (1) A vector killed by all positive modes and of positive weight generates a proper
ideal (its descendants never reach the vacuum line), hence lies in the maximal ideal. (2) \emph{Non-vanishing of
the symbol} is the delicate point. Embed $V^k$ in the Wakimoto free-field module (faithful at non-critical level,
Chapter~\ref{ch:L08}); $e(z)$ becomes $\beta(z)$, and the conditions $e_rv=0$ ($r\ge0$) say the image is free of
$\gamma$-modes, $h_0v=2v$ and $h_rv=0$ ($r>0$) become a triangular system of differential equations for the
coefficients of the $\beta$-oscillators. The system forces a non-zero coefficient on a monomial with no
derivatives (``depth zero''), and the Wakimoto map is compatible with Li filtrations, so such a monomial
survives in $R$: $\mathrm{symbol}(v)\neq0$. (3) The symbol is homogeneous of degree $2m+1$, killed by
$\{e,-\}$ (because $e_0v=0$) and of $h$-weight $2$. In $\C[e,h,f]=\C[\ltwentyQ]\otimes\mathcal H$ (harmonics), the
$\ltwentysl$-highest vectors of weight $2$ and degree $2m+1$ form the line $\C e\ltwentyQ^m$; hence the symbol is $c\,e\ltwentyQ^m$,
$c\neq0$. (4) The image of the maximal ideal in $R_{V^k}$ is a Poisson ideal containing $e\ltwentyQ^m$; bracketing
with $f$ (and using $\{f,\ltwentyQ\}=0$) gives $h\ltwentyQ^m$ and $f\ltwentyQ^m$ up to non-zero constants. \qed

\begin{physics}
The classical limit of the null state is the highest weight component of $e\cdot(\text{Casimir})^m$; the
three relations are its $\ltwentysl$-orbit. Physically: the Higgs branch chiral ring has the relation
``moment map times $(\text{Casimir})^m=0$''.
\end{physics}

\begin{status}
All four proved for every candidate, standard axioms. In particular, combined with
\lean{boundary_primitive_exists}, the relations hold for every $m\ge1$
(\lean{uniform_boundaryIdeal_le_kernel}). This supersedes the conditional status in Chapter~\ref{ch:L09},
Section~\ref{l09:sec:boundary} for the \emph{containment}; exactness is Theorem~\ref{l20:thm:c2arc}.
\end{status}

\begin{theorem}[Arc relations vanish; currents generate]\label{l20:thm:m0arc}
For every candidate $C$: the three base relations vanish on the degree-zero Li classes of the currents; all
coefficients $(x_i\ltwentyQ^m)_{[n]}$ vanish under $\pi$; and $\pi\colon\C[x_{i,n}]\to\ltwentygr L_k$ is surjective.
\end{theorem}
\noindent File \leanpath{VOA/\allowbreak Filtered/\allowbreak BoundaryArcLiMap.lean}.
\begin{leancode}
theorem baseIdeal_le_kernel : BoundaryAuxiliary.baseIdeal (R := ℂ) m ≤
    RingHom.ker (MvPolynomial.aeval (fun i =>
      liDegreeZeroClass (simpleVOA C) 1 ⟨gen C i,gen_mem C i⟩)).toRingHom := by

theorem sourceIdeal_le_kernel : BoundaryAuxiliary.sourceIdeal (R := ℂ) m ≤
    RingHom.ker (polynomialMap C).toRingHom := by

theorem polynomialMap_surjective : Function.Surjective (polynomialMap C) := by
\end{leancode}

\begin{dictionary}
\lean{RingHom.ker φ} is the kernel; \lean{I ≤ ker φ} says $\varphi$ kills $I$. The first statement is the passage from
$R_V$ to the degree-zero part of $\ltwentygr V$ (they are the same space, $F^0/F^1$, but the library proves the
algebra structures agree, \lean{c2ToLiAlgebra}). The second is the arc version. The third is Li's generation
theorem for affine quotients: every Li layer is spanned by current words of the right depth.
\end{dictionary}

\textbf{Proof on paper.} $\pi$ is a homomorphism of differential algebras (it sends $x_{i,n}$ to
$\partial^n\bar x_i/n!$), so it sends the generic arc $x_i(t)=\sum x_{i,n}t^n$ to $e^{t\partial}\bar x_i$, and
$f(x(t))\mapsto e^{t\partial}f(\bar x)$. If $f(\bar x)=0$ in $\ltwentygr V$, every coefficient vanishes. Surjectivity:
by Li's theorem (proved in the library for every graded quotient of $V^k(\ltwentysl)$), $F^pV_N$ is spanned by current
monomials $a^1_{(-n_1-1)}\cdots a^r_{(-n_r-1)}\ltwentyvac$ with $\sum n_i\ge p$; their classes in the exact layer
$(p,N)$ are images of $\prod x_{a^i,n_i}$. \qed

\begin{physics}
Strong generation by the currents: every state is a normally ordered polynomial in currents and their
derivatives. Classically, the arc of the Higgs branch surjects onto the singular support.
\end{physics}

\begin{status}
Proved for every candidate; standard axioms.
\end{status}

%======================================================================
\section{Degeneration to the auxiliary arc algebra}\label{l20:sec:m1}
%======================================================================

The upper bound $\dim(\ltwentyB)_N\le[q^N]\,\mathrm{product}$ is the heart of the paper (its Sections~2--5). The first
step replaces $\ltwentyB$ by a ``monomial-like'' algebra that is easier to count.

\begin{definition}[Auxiliary relations and weight]\label{l20:def:aux}
Give $e,f$ auxiliary weight $1$ and $h$ weight $0$ (for all jet orders). The auxiliary relations are
$x_i(ef)^m$; $\ltwentyCm=\ltwentyJ\bigl(\C[e,h,f]/(x_i(ef)^m)\bigr)$. The \emph{initial ideal} of an ideal $I$ of jet
polynomials is the ideal generated by the top auxiliary-weight components of the elements of $I$.
\end{definition}
\noindent File \leanpath{VOA/\allowbreak Commutative/\allowbreak BoundaryAuxiliaryInitial.lean}.
\begin{leancode}
def auxiliaryRelation (m : ℕ) (i : Fin 3) : BasePoly R :=
  MvPolynomial.X i * (MvPolynomial.X 0 * MvPolynomial.X 2) ^ m

def auxiliaryIdeal (m : ℕ) : Ideal (JetPoly R) :=
  VOA.InfiniteJet.coefficientIdeal (Set.range (auxiliaryRelation (R := R) m))

def currentWeight (i : Fin 3) : ℕ := if i = 1 then 0 else 1
def auxWeight (v : Fin 3 × ℕ) : ℕ := currentWeight v.1

def initialIdeal (I : Ideal (JetPoly R)) : Ideal (JetPoly R) :=
  Ideal.span {p | ∃ g ∈ I, ∃ d : ℕ,
    (∀ k, d < k → MvPolynomial.weightedHomogeneousComponent auxWeight k g = 0) ∧
      MvPolynomial.weightedHomogeneousComponent auxWeight d g = p}

theorem relationCoefficient_top (m n : ℕ) (i : Fin 3) :
    MvPolynomial.weightedHomogeneousComponent auxWeight (currentWeight i + 2 * m)
      (VOA.InfiniteJet.relationCoefficient (baseRelation (R := R) m i) n) =
        (4 : JetPoly R) ^ m *
          VOA.InfiniteJet.relationCoefficient (auxiliaryRelation (R := R) m i) n := by

theorem auxiliaryIdeal_le_initialIdeal (m : ℕ) :
    auxiliaryIdeal (R := K) m ≤ initialIdeal (sourceIdeal (R := K) m) := by
\end{leancode}
\noindent File \leanpath{VOA/\allowbreak Commutative/\allowbreak BoundaryAuxiliaryHilbert.lean}.
\begin{leancode}
def auxiliaryPiece (m N : ℕ) : Submodule K (Auxiliary (R := K) m) :=
  (MvPolynomial.weightedHomogeneousSubmodule K conformalWeight N).map
    (Ideal.Quotient.mkₐ K (auxiliaryIdeal (R := K) m)).toLinearMap

theorem sourcePiece_finrank_le_auxiliaryPiece (m N : ℕ) :
    Module.finrank K (sourcePiece (K := K) m N) ≤
      Module.finrank K (auxiliaryPiece (K := K) m N) := by
\end{leancode}

\begin{dictionary}
\lean{weightedHomogeneousComponent auxWeight d g} is the part of $g$ of auxiliary weight $d$. In
\lean{initialIdeal}, $d$ is any weight above which $g$ has no components, so the spanned elements are the top
components. \lean{relationCoefficient_top} is the paper's equation (2.15): expanding
$x_i\ltwentyQ^m=x_i(h^2+4ef)^m$ binomially, the top auxiliary-weight part of each arc coefficient is $4^m$ times the
corresponding coefficient of $x_i(ef)^m$ (the term $(4ef)^m$ has the most $e,f$). Hence the auxiliary ideal lies
in the initial ideal of the source ideal (\lean{auxiliaryIdeal_le_initialIdeal}; the field \lean{K} has
characteristic $0$ so $4^m$ is invertible; ``equality of initial ideals is neither needed nor assumed''). The
final theorem is Proposition~2.3 of the paper: $\dim(\ltwentyB)_N\le\dim(\ltwentyCm)_N$.
\end{dictionary}

\textbf{Proof on paper.} Fix the conformal weight $N$; there only finitely many jet variables occur, so we are in
a finite-dimensional space $P_N$ of polynomials. Filter $P_N$ by auxiliary weight. For any ideal $I$
homogeneous for conformal weight, $\dim(P/I)_N=\dim(P/\operatorname{in}I)_N$ (the initial forms of a basis of
$I_N$ adapted to the filtration form a basis of $(\operatorname{in}I)_N$). If $J\subseteq\operatorname{in}I$ then
$\dim(P/J)_N\ge\dim(P/\operatorname{in}I)_N$. With $J=$ auxiliary ideal this gives the claim. (The library proves
the inequality directly, by injecting each top-weight slice of $P/I$ into the corresponding slice of $P/J$.)
\qed

\begin{physics}
A ``classical limit of the classical limit'': scaling $e,f$ to be large compared with $h$, the Casimir
$h^2+4ef$ is dominated by $4ef$. The degenerate theory has a larger (or equal) state space, so counting its
states gives an upper bound.
\end{physics}

\begin{status}
Proved; standard axioms. No degeneration map or Hilbert inequality is assumed.
\end{status}

%======================================================================
\section{Restricted dual, Gordon divisibility and the partition filtration}\label{l20:sec:m23}
%======================================================================

This is the combinatorial core: a bound on $\dim(\ltwentyCm)_N$ by an explicit fermionic sum. We explain it on
paper in the language of Feigin--Stoyanovsky ``functional realisations''; the library follows the paper's
Sections~3--4 closely, with every step proved (the campaign ledger lists these two steps in detail).

\subsection{Sectors and particle polynomials}

$\ltwentyCm$ is spanned by monomials $\prod x_{i,n}$. Fix the numbers $a,b,c$ of factors of colour $e,f,h$ and the
total derivative degree $D=\sum n$; the span of these monomials is a finite-dimensional \emph{sector}, and the
weight-$N$ piece is spanned by the sectors with $a+b+c+D=N$.

\begin{definition}[Sector, particle polynomial]\label{l20:def:sector}
Label the $a+b+c$ factors by a finite set $\iota$ with a colouring $\iota\to\{e,h,f\}$. For a monomial with jet
orders $r\colon\iota\to\N$, $\sum r=D$, write $\mathrm{mode}(r)=\prod_{s\in\iota}x_{\mathrm{col}(s),r(s)}$. To a linear
functional $L$ on the sector associate the \emph{particle polynomial}
$\Phi_L(z)=\sum_{r}L(\mathrm{mode}(r))\,z^r\in\C[z_s:s\in\iota]$, homogeneous of degree $D$.
\end{definition}
\noindent File \leanpath{VOA/\allowbreak Commutative/\allowbreak GordonRestrictedDual.lean}.
\begin{leancode}
def quotientMode (m : ℕ) (color : ι → Fin 3) (r : ι →₀ ℕ) : Auxiliary (R := K) m :=
  Ideal.Quotient.mk (auxiliaryIdeal (R := K) m) (∏ i, X (color i, r i))

def sector (m D : ℕ) (color : ι → Fin 3) : Submodule K (Auxiliary (R := K) m) :=
  Submodule.span K (quotientMode m color '' (↑(degreeFiber (ι := ι) D) : Set (ι →₀ ℕ)))

theorem particleMap_injective (m D : ℕ) (color : ι → Fin 3) :
    Function.Injective (particleMap (K := K) m D color) := by

theorem finrank_functionalModel (m D : ℕ) (color : ι → Fin 3) :
    Module.finrank K (functionalModel (K := K) m D color) =
      Module.finrank K (sector (K := K) m D color) := by
\end{leancode}

\begin{dictionary}
\lean{ι →₀ ℕ} is a finitely supported function (an exponent vector); \lean{degreeFiber D} is the set of those of
total degree $D$; \lean{''} is the image of a set. So \lean{sector m D color} is the span of all coloured mode
products of derivative degree $D$ in the actual quotient $\ltwentyCm$. \lean{particleMap} sends a functional to its
particle polynomial; it is injective (a functional is determined by its values on spanning monomials), and
\lean{functionalModel} is its image, of the same dimension as the sector (dual of a finite-dimensional space).
\end{dictionary}

\textbf{Proof on paper.} A functional on a span is determined by its values on the spanning set, which are the
coefficients of $\Phi_L$. \qed

\begin{physics}
$\Phi_L(z)$ is a ``correlation function'' $\langle L,\prod_s x_{\mathrm{col}(s)}(z_s)\rangle$ of classical fields at
points $z_s$: the generating function of all mode products. Relations among modes become vanishing conditions
on these functions.
\end{physics}

\begin{status}
Proved; standard axioms.
\end{status}

\subsection{Vanishing conditions and Gordon divisibility}

Because $\ltwentyCm$ is a quotient by the arc ideal of $x(ef)^m$, every particle polynomial $\Phi$ in the functional
model is \emph{symmetric in the variables of each colour} (the variables commute), and vanishes on the
\emph{relation diagonals}: setting $m$ $e$-variables, $m$ $f$-variables and one further variable (of any
colour) equal to a common $z$ kills $\Phi$, \emph{for arbitrary values of the remaining ``spectator''
variables} (\lean{functionalModel_relationDiagonal}, in \leanpath{GordonRestrictedDual.lean}, resting on
\leanpath{GordonParticleDiagonals.lean}). Indeed the
relation $x(z)e(z)^mf(z)^m=0$ holds as an identity of fields. From these conditions one derives divisibility.

\paragraph{Set-up for the divisor.} In a sector with $a$ $e$'s and $b$ $f$'s, choose $R\le\min(a,b)$ of the
$e$-variables and $R$ of the $f$-variables and match them into $R$ \emph{$e$--$f$ pairs}; the other $s=a-R$
$e$-variables, $t=b-R$ $f$-variables and the $c$ $h$-variables are \emph{residual}. The \emph{pair evaluation}
sets the two variables of each pair equal. A map $f\colon\{\text{pairs}\}\to\tau$ groups the pairs into
\emph{clusters} (blocks); its \emph{shape} is the partition $\lambda$ of $R$ formed by the block sizes, and $r_i$ is
the number of blocks of size $i$. The \emph{cluster evaluation} sets all pair variables of one block equal to a
single cluster point $w_u$ ($u\in\tau$). The \emph{residual filtration} step indexed by $\lambda$ consists of the
pair-evaluated functions killed by every cluster evaluation of lexicographically larger shape. Finally define
\begin{equation}\label{l20:eq:Delta}
\Delta_\lambda=\prod_{u<v}(w_u-w_v)^{2\min(|u|,|v|)}\cdot\prod_{|u|=m}\ \prod_{y\ \text{residual}}(w_u-y)
\cdot\prod_{x\ \text{residual }e}\ \prod_{y\ \text{residual }f}(x-y),
\end{equation}
the product over pairs of distinct blocks, over blocks of the maximal size $m$ against all residual variables,
and over residual $e$--$f$ pairs.

\begin{theorem}[Full Gordon divisor]\label{l20:thm:gordon}
Let $\Phi$ be a particle polynomial of the functional model (for the sector with pairs as above) such that
(i) $\Phi$ vanishes when any residual $e$-variable is set equal to any residual $f$-variable after pair
evaluation (``no further pair can be formed''), and (ii) its pair evaluation lies in the residual filtration step
of shape $\lambda=\operatorname{shape}(f)$. Then $\Delta_\lambda$ divides the cluster evaluation of $\Phi$, and
\begin{equation*}
\deg\Delta_\lambda=\sum_{u\neq v}\min(|u|,|v|)+r_m(s+t+c)+st
=\sum_{i,j}\min(i,j)\,r_ir_j-\sum_i i\,r_i+r_m(s+t+c)+st .
\end{equation*}
\end{theorem}
\noindent File \leanpath{VOA/\allowbreak Commutative/\allowbreak GordonFullDivisor.lean}.
\begin{leancode}
def edgeExponent (m : ℕ) (f : ρ → τ) : Edge (α := α) (β := β) (γ := γ) m f → ℕ
  | .inl ij => 2 * min (blockSize f ij.val.1) (blockSize f ij.val.2)
  | .inr _ => 1

def fullDivisor (m : ℕ) (f : ρ → τ) : MvPolynomial (Target τ α β γ) K :=
  ∏ e : Edge (α := α) (β := β) (γ := γ) m f,
    collisionFactor (edgePair m f e).1 (edgePair m f e).2 ^ edgeExponent m f e

theorem fullDegree_eq (m : ℕ) (f : ρ → τ) :
    fullDegree (α := α) (β := β) (γ := γ) m f = clusterDegree f +
      Fintype.card {u : τ // blockSize f u = m} *
        (Fintype.card α + Fintype.card β + Fintype.card γ) +
      Fintype.card α * Fintype.card β := by

theorem fullDivisor_dvd [CharZero K] (m D : ℕ) (f : ρ → τ)
    (F : MvPolynomial (Source ρ α β γ) K)
    (hF : F ∈ functionalModel m D (sourceColor (ρ := ρ) (α := α) (β := β) (γ := γ)))
    (hnext : ∀ x : α, ∀ y : β, nextPairEval x y F = 0)
    (U : Submodule K (MvPolynomial (Target ρ α β γ) K))
    (hG : pairEval F ∈ residualFiltration (τ := τ) U (shape f)) :
    fullDivisor (K := K) (α := α) (β := β) (γ := γ) m f ∣ clusterEval f F := by
\end{leancode}

\begin{dictionary}
Index types: \lean{ρ} labels the pairs, \lean{τ} the blocks (linearly ordered, so ``$u<v$'' makes sense),
\lean{α}, \lean{β}, \lean{γ} the residual $e$, $f$, $h$ variables. \lean{Source ρ α β γ} indexes the original variables
(pair-$e$, residual $e$, pair-$f$, residual $f$, $h$), \lean{Target} the variables after identification.
\lean{blockSize f u} is the number of pairs in block $u$; \lean{shape f} the partition of block sizes.
\lean{Edge m f} is the disjoint union of the three families in \eqref{l20:eq:Delta} (increasing pairs of blocks;
blocks of size $m$ times residual variables; residual $e$ times residual $f$), \lean{collisionFactor i j} $=x_i-x_j$, and
\lean{edgeExponent} is $2\min$ on block pairs and $1$ otherwise; so \lean{fullDivisor m f} is $\Delta_\lambda$.
\lean{clusterDegree f} $=\sum_{u<v}2\min(|u|,|v|)=\sum_{i,j}\min(i,j)r_ir_j-\sum_i ir_i$ (the diagonal $u=v$
contributes $\sum_u|u|=\sum_i ir_i$). In \lean{fullDivisor_dvd}: \lean{hF} says $\Phi$ is in the functional model of the
sector, \lean{hnext} is condition (i) (\lean{nextPairEval x y} = pair evaluation followed by $x=y$), \lean{hG} is
condition (ii) (\lean{residualFiltration U λ} is $U$ intersected with the kernels of the cluster evaluations of
lexicographically larger shape), and \lean{clusterEval f} is the cluster evaluation.

\textbf{Relation to the sector bound.} The exponent of Theorem~\ref{l20:thm:sectorbound} is
$\deg\Delta_\lambda+(a+b+c)$, since $\bigl(\sum_{i,j}\min(i,j)r_ir_j+\sum_iir_i+(r_m+1)(s+t+c)+st\bigr)-\deg\Delta_\lambda
=2R+s+t+c=a+b+c$; the extra $a+b+c$ is exactly the shift from derivative degree $D$ to conformal weight
$D+a+b+c$ in the coefficient index there. So $\deg\Delta_\lambda$ is the derivative-degree cost of the clusters.

The decisive one-variable input is the \emph{symmetric occupancy lemma} (paper Lemma~4.2,
\leanpath{SymmetricOccupancy.lean}): for a symmetric polynomial, specialising a subset of $s$ variables to $z+\delta$
and the rest to $z$, the coefficient of $\delta^d$ is a polynomial in $s$ of degree at most $d$; so vanishing at
many occupancy values forces high-order vanishing (\leanpath{GordonTwoBlock.lean}).
\end{dictionary}

\textbf{Proof on paper (outline).} Collapse each cluster to a single point. When two collapsed points (or a
collapsed point and a residual particle) approach each other, the relation diagonals --- applied with the
other particles as spectators --- give vanishing at a whole range of ``occupancy'' specialisations (how many of
the particles of one block are moved), and the occupancy lemma converts this into vanishing to a definite order
along the collision divisor. Different collision divisors are coprime, so the product $\Delta_\lambda$ of all these
factors divides the evaluated polynomial (\lean{prod_pow_dvd_of_pairwise_prime_varying}). The library computes
the total degree of $\Delta_\lambda$ exactly (\lean{fullDegree_eq}); added to $a+b+c$ it is the exponent appearing in
Theorem~\ref{l20:thm:sectorbound} below (see the dictionary). We do not reproduce the case analysis of the individual collision orders; it is the paper's
Section~4, formalised in \leanpath{GordonPairedRelations.lean}, \leanpath{GordonCluster.lean},
\leanpath{GordonCommonFactors.lean} and \leanpath{GordonFullDivisor.lean}. \qed

\begin{physics}
This is the classical analogue of \emph{exclusion statistics}: the $(m+1)$-st pair cannot sit at the same
point as $m$ others ($e^{m+1}f^m=0$-type relations), so wave functions vanish on collisions to a definite order,
as in Laughlin/Haldane fractional statistics. The coefficients $\min(i,j)$ of the resulting quadratic form play
the role of a ``statistical interaction'' between clusters of sizes $i$ and $j$ (a heuristic reading).
\end{physics}

\begin{status}
Proved (all divisor families and their exact degrees); standard axioms.
\end{status}

\subsection{The pair/partition filtration and the sector bound}

\begin{theorem}[Sector bound]\label{l20:thm:sectorbound}
For a sector with $a$ $e$'s, $b$ $f$'s, $c$ $h$'s and derivative degree $D$,
\begin{multline*}
\dim\mathrm{sector}\ \le\ \sum_{R\le\min(a,b)}\ \sum_{\substack{r_1,\dots,r_m\ge0\\ \sum ir_i=R}}
[q^{D+a+b+c}]\\ \frac{q^{\sum_{i,j}\min(i,j)r_ir_j+\sum_iir_i+(r_m+1)(s+t+c)+st}}{(q)_{r_1}\cdots(q)_{r_m}(q)_s(q)_t(q)_c},
\qquad s=a-R,\ t=b-R.
\end{multline*}
\end{theorem}
\noindent Files \leanpath{VOA/\allowbreak Commutative/\allowbreak QMultiplicityProfile.lean},
\leanpath{VOA/\allowbreak Commutative/\allowbreak GordonPartitionProfile.lean},
\leanpath{VOA/\allowbreak Commutative/\allowbreak GordonSectorCharacter.lean}.
\begin{leancode}
def profileQuadratic {m : ℕ} (r : Fin m → ℕ) : ℕ :=
  ∑ i, ∑ j, min (i.val + 1) (j.val + 1) * r i * r j

def profileLinear {m : ℕ} (r : Fin m → ℕ) : ℕ :=
  ∑ i, (i.val + 1) * r i

def profileEnergy {m : ℕ} (r : Fin m → ℕ) : ℕ := profileQuadratic r + profileLinear r

def multiplicityResidualTerm {m : ℕ} (r : Fin (m + 1) → ℕ) (s t c : ℕ) : PowerSeries K :=
  PowerSeries.X ^ (profileEnergy r + (r (Fin.last m) + 1) * (s + t + c) + s * t) *
    (∏ i, reciprocalQFactorial (r i)) * reciprocalQFactorial s *
      reciprocalQFactorial t * reciprocalQFactorial c

def Profile (m R : ℕ) := {r : Fin m → ℕ // ∑ i, (i.val + 1) * r i = R}

theorem sector_finrank_le_profile_sum {K : Type*} [Field K] [CharZero K]
    (m D a b c : ℕ) :
    (Module.finrank K (sector (K := K) (m + 1) D (originalColor a b c)) : ℚ) ≤
      ∑ R : PairCount a b, ∑ r : Profile (m + 1) R.val,
        PowerSeries.coeff (D + a + b + c)
          (multiplicityResidualTerm (K := ℚ) r.val (a - R.val) (b - R.val) c) := by
\end{leancode}

\begin{dictionary}
\lean{Fin m → ℕ} is a vector $(r_1,\dots,r_m)$ written with $0$-based indices; \lean{i.val + 1} restores the
$1$-based cluster size $i$. \lean{Fin.last m} is the last index (largest cluster size). \lean{PairCount a b} is
$\{0,\dots,\min(a,b)\}$; \lean{Profile m R} is the set of $(r_i)$ with $\sum ir_i=R$ (cluster profile of $R$ pairs).
\lean{originalColor a b c} colours $a+b+c$ slots $e^af^bh^c$. \textbf{Index shift:} the theorem is stated for
\lean{sector … (m + 1) …}, i.e.\ paper index $m+1$, with profiles of length $m+1$; this is a harmless
reparametrisation that avoids the hypothesis $m\ge1$. The extra $a+b+c$ in the coefficient index converts derivative
degree to conformal weight (each current has weight one).
\end{dictionary}

\textbf{Proof on paper.} Filter the functional model by \emph{how many pairs can be formed} and then by the
\emph{partition} $\lambda$ of the pairs into clusters, in descending lexicographic order: the step indexed by $\lambda$
is the kernel of all evaluations at partitions larger than $\lambda$. Each subquotient injects, by evaluation at
$\lambda$ followed by division by $\Delta_\lambda$ (Theorem~\ref{l20:thm:gordon}), into the space of polynomials in the
cluster and residual points that are symmetric within each block type; the Hilbert series of such polynomials
is $\prod_i1/(q)_{r_i}\cdot1/(q)_s(q)_t(q)_c$ (fundamental theorem of symmetric polynomials, applied blockwise:
\leanpath{ColoredSymmetricHilbert.lean}, \leanpath{TwoBlockSymmetricInvariant.lean}), shifted by the degree of
$\Delta_\lambda$. Summing the subquotient bounds over the finite filtration (a telescope) gives the bound. \qed

\begin{physics}
Each term is a gas of free ``quasi-particles'': $r_i$ bound states of $i$ $e$--$f$ pairs, and $s,t,c$ free
$e$, $f$, $h$ particles, with mutual exclusion energies in the exponent. This is how fermionic character
formulas are derived in the physics of integrable models (Bethe ansatz counting) and in Feigin--Stoyanovsky's
treatment of $\widehat{\ltwentysl}$ principal subspaces.
\end{physics}

\begin{status}
Proved; standard axioms. Only an inequality is proved (that is all that is needed); no claim that the
filtration bound is sharp is made at this stage. Sharpness follows a posteriori from the squeeze.
\end{status}

\begin{theorem}[Global auxiliary bound]\label{l20:thm:auxbound}
$\dim(\ltwentyCm)_N\le[q^N]\,\mathcal S_m$, where $\mathcal S_m=\sum_{r_1,\dots,r_m\ge0}q^{\sum\min(i,j)r_ir_j+\sum ir_i}
\prod_{i<m}(q)_{r_i}^{-1}\,\mathrm{res}(r_m)$ is the paper's source series (5.8).
\end{theorem}
\noindent Files \leanpath{VOA/\allowbreak Commutative/\allowbreak QMultiplicityProfile.lean},
\leanpath{VOA/\allowbreak Commutative/\allowbreak GordonConformalCharacter.lean}.
\begin{leancode}
def multiplicityTerm {m : ℕ} (r : Fin (m + 1) → ℕ) : PowerSeries K :=
  PowerSeries.X ^ profileEnergy r * profileInitialDenominator r * residualSum (r (Fin.last m))

def multiplicitySource (m : ℕ) : PowerSeries K :=
  VOA.LocallyFinitePowerSeriesSum.sum (fun r : Fin (m + 1) → ℕ => multiplicityTerm r)

theorem multiplicitySource_eq_hilbertBoundSeries (m : ℕ) :
    multiplicitySource (K := K) m = hilbertBoundSeries (m + 1) := by

theorem auxiliaryPiece_finrank_le_source {K : Type*} [Field K] [CharZero K] (m N : ℕ) :
    (Module.finrank K (auxiliaryPiece (K := K) (m + 1) N) : ℚ) ≤
      PowerSeries.coeff N (multiplicitySource (K := ℚ) m) := by
\end{leancode}

\begin{dictionary}
\lean{profileInitialDenominator r} is $\prod_{i<\text{last}}(q)_{r_i}^{-1}$; the last cluster size carries the
residual sum \lean{residualSum (r last)} (Definition~\ref{l20:def:chain}), which is the result of summing the three
residual species $s,t,c$ in \lean{multiplicityResidualTerm} (proved in the same file). \lean{multiplicitySource m}
is the paper's (5.8) at index $m+1$; \lean{multiplicitySource_eq_hilbertBoundSeries} is the paper's passage
(5.8)$\to$(5.12): grouping profiles by total multiplicity and peeling the first coordinate turns the quadratic form
into the cumulative-count chain $\sum_jn_j(n_j+1)$ with $n_j=r_j+\dots+r_m$. The bound sums
Theorem~\ref{l20:thm:sectorbound} over sectors, injecting the finite set of sector labels into the profile/residual
index set and using that all terms have non-negative coefficients (\lean{masterTerm_nonneg}), so adding unused terms
only increases the sum.
\end{dictionary}

\textbf{Proof on paper.} Sum the sector bounds over $a+b+c+D=N$; reindex; extend to all profiles by
positivity. The identity $\sum_{i,j}\min(i,j)r_ir_j+\sum ir_i=\sum_jn_j(n_j+1)$ with $n_j=\sum_{i\ge j}r_i$ is the
standard Andrews--Gordon change of variables. \qed

\begin{physics}
$\mathcal S_m$ is the full quasi-particle partition function of the degenerate theory $\ltwentyCm$.
\end{physics}

\begin{status}
Proved; standard axioms.
\end{status}

%======================================================================
\section{The \texorpdfstring{$q$}{q}-series identities}\label{l20:sec:m4}
%======================================================================

\begin{theorem}[Chain, trace and product]\label{l20:thm:qseries}
For $m\ge1$, $p=2m+1$:
\begin{equation*}
S_m=\frac{T_m}{(q)^2_\infty},\qquad T_m=\frac{(q^p;q^p)_\infty^3}{(q)_\infty},\qquad
S_m=\frac{(q^p;q^p)^3_\infty}{(q)^3_\infty}.
\end{equation*}
\end{theorem}
\noindent Files \leanpath{VOA/\allowbreak Commutative/\allowbreak QHilbertProduct.lean},
\leanpath{VOA/\allowbreak Commutative/\allowbreak QChainTrace.lean}.
\begin{leancode}
theorem hilbertBoundSeries_eq_trace (m : ℕ) :
    hilbertBoundSeries (K := K) m = fermionicTrace m * reciprocalEuler ^ 2 := by

theorem fermionicTrace_eq (m : ℕ) (hm : 1 ≤ m) :
    fermionicTrace (K := K) m =
      (PowerSeries.expand (2 * m + 1) (by omega) (euler : PowerSeries K)) ^ 3 *
        reciprocalEuler := by

theorem hilbertBoundSeries_product (m : ℕ) (hm : 1 ≤ m) :
    hilbertBoundSeries (K := K) m =
      (PowerSeries.expand (2 * m + 1) (by omega) (euler : PowerSeries K)) ^ 3 *
        reciprocalEuler ^ 3 := by
\end{leancode}
\noindent Files \leanpath{VOA/\allowbreak Commutative/\allowbreak QBinomialResidual.lean},
\leanpath{VOA/\allowbreak Commutative/\allowbreak QBinomialAlternating.lean},
\leanpath{VOA/\allowbreak Commutative/\allowbreak QJacobiCube.lean},
\leanpath{VOA/\allowbreak Commutative/\allowbreak QDurfeeChain.lean}.
\begin{leancode}
theorem residualSum_eq (n : ℕ) :
    (residualSum n : PowerSeries R) =
      qFactorial PowerSeries.X n ^ 2 * reciprocalEuler ^ 2 *
        reciprocalQFactorial (2 * n + 1) := by

theorem finite_alternating (n : ℕ) :
    (qFactorial (PowerSeries.X : PowerSeries K) n) ^ 2 =
      ∑ r ∈ Finset.range (n + 1),
        (-1 : PowerSeries K) ^ r * ((2 * r + 1 : ℕ) : PowerSeries K) *
          PowerSeries.X ^ ((r + 1).choose 2) *
            gaussian PowerSeries.X (2 * n + 1) (n - r) := by

noncomputable def alternatingTerm (r : ℕ) : PowerSeries R :=
  (-1 : PowerSeries R) ^ r * ((2 * r + 1 : ℕ) : PowerSeries R) *
    PowerSeries.X ^ ((r + 1).choose 2)

noncomputable def alternatingTheta : PowerSeries R :=
  VOA.LocallyFinitePowerSeriesSum.sum alternatingTerm

theorem jacobi_cube : (alternatingTheta : PowerSeries K) = euler ^ 3 := by

theorem chainSeries_eq (ℓ r : ℕ) (hℓ : 1 ≤ ℓ) :
    chainSeries (R := R) ℓ r =
      PowerSeries.X ^ (ℓ * (r * (r + 1))) * reciprocalEuler := by
\end{leancode}

\begin{dictionary}
\lean{gaussian q n k} is the Gaussian binomial $\binom nk_q$, defined by the $q$-Pascal recursion
(\leanpath{QBinomial.lean}). \lean{(r + 1).choose 2} $=r(r+1)/2$. \lean{alternatingTheta} is
$\sum_{r\ge0}(-1)^r(2r+1)q^{r(r+1)/2}$, and \lean{jacobi_cube} is Jacobi's identity
$(q)_\infty^3=\sum_r(-1)^r(2r+1)q^{r(r+1)/2}$. \lean{chainSeries ℓ r} is the descending chain of length $\ell$
started at $r$ with the shifted kernel (\leanpath{QDurfeeChain.lean}); \lean{chainSeries_eq} is its closed form
$q^{\ell r(r+1)}/(q)_\infty$, proved by collapsing each bounded inner sum with a finite shifted identity
(\lean{chainKernel_eq}) and inducting on the chain length (a Durfee-square type argument). Everything is an identity in $K[[q]]$ for any field $K$ (only
characteristic-free algebra is used; \lean{finite_alternating} uses Hahn series over $\Z$-exponents to divide by
$q$-powers).
\end{dictionary}

\textbf{Proof on paper.} (1) \emph{Residual sums.} Summing the two free species with Euler's identity
$\sum_jq^{Aj}/(q)_j=1/(q^A;q)_\infty$ and the third with the $q$-binomial theorem gives
$\mathrm{res}(n)=(q)_n^2/((q)_{2n+1}(q)^2_\infty)$ (\lean{residualSum_eq}); inserting this at the end of the chain gives
$S_m=T_m/(q)^2_\infty$. (2) \emph{Terminal factor.} The finite identity
$(q)_n^2=\sum_{r=0}^n(-1)^r(2r+1)q^{\binom{r+1}2}\binom{2n+1}{n-r}_q$ (\lean{finite_alternating}) rewrites
$K_0(n)=(q)^2_n/(q)_{2n+1}$ as $\sum_r(-1)^r(2r+1)q^{\binom{r+1}2}/((q)_{n-r}(q)_{n+r+1})$. (3) \emph{Chain collapse.}
After exchanging the finite sums, $T_m=\sum_r(-1)^r(2r+1)q^{\binom{r+1}2}\cdot\mathrm{chain}_m(r)$ and
$\mathrm{chain}_m(r)=q^{mr(r+1)}/(q)_\infty$ (\lean{chainSeries_eq}). (4) \emph{Jacobi.} Since
$\binom{r+1}2+mr(r+1)=p\binom{r+1}2$, $T_m=\frac1{(q)_\infty}\sum_r(-1)^r(2r+1)q^{p\binom{r+1}2}
=\frac{(q^p;q^p)^3_\infty}{(q)_\infty}$ by Jacobi's identity with $q\mapsto q^p$. \qed

\begin{physics}
A ``fermionic = bosonic'' identity: the quasi-particle sum (fermionic) equals the product (bosonic,
free-field) form of the same partition function, as in the Rogers--Ramanujan and Andrews--Gordon identities that
compute characters of Virasoro minimal models. Jacobi's triple product is the vanishing of all but the
Weyl-orbit contributions in the $\widehat{\ltwentysl}$ denominator identity.
\end{physics}

\begin{status}
Proved for every $m\ge1$ over any field; standard axioms. These are pure $q$-series facts, independent of any
VOA.
\end{status}

%======================================================================
\section{The Hilbert-series upper bound}\label{l20:sec:m5}
%======================================================================

\begin{theorem}[Upper bound]\label{l20:thm:upper}
For $m\ge1$ and all $N$: $\dim L_{k_m}(\ltwentysl)_N\le\dim(\ltwentyB)_N\le[q^N]\,\mathrm{product}(m)$.
\end{theorem}
\noindent File \leanpath{VOA/\allowbreak Commutative/\allowbreak GordonConformalCharacter.lean}.
\begin{leancode}
theorem sourcePiece_finrank_le_product {K : Type*} [Field K] [CharZero K] (m N : ℕ) :
    (Module.finrank K (sourcePiece (K := K) (m + 1) N) : ℚ) ≤
      PowerSeries.coeff N ((PowerSeries.expand (2 * (m + 1) + 1) (by omega) euler) ^ 3 *
        (reciprocalEuler : PowerSeries ℚ) ^ 3) := by
\end{leancode}
\noindent File \leanpath{VOA/\allowbreak Examples/\allowbreak AffineSl2BoundaryVacuumCharacter.lean}.
\begin{leancode}
theorem vacuum_finrank_le_sourcePiece (m : ℕ) (hm : 1 ≤ m) (N : ℕ) :
    Module.finrank ℂ ((vacuum m).grading N) ≤
      Module.finrank ℂ (sourcePiece (K := ℂ) m N) := by

theorem coeff_character_le_product (m : ℕ) (hm : 1 ≤ m) (N : ℕ) :
    PowerSeries.coeff N (character m) ≤ PowerSeries.coeff N (product m) := by
\end{leancode}
\noindent File \leanpath{VOA/\allowbreak Filtered/\allowbreak LiAssociatedGraded.lean}.
\begin{leancode}
theorem sum_liWeightLayer_finrank_eq_grading (W : VertexOperatorAlgebra 𝒱) (N : ℕ) :
    (∑ p ∈ Finset.range (N + 1), Module.finrank ℂ (liWeightLayer W p N)) =
      Module.finrank ℂ (W.grading N) := by
\end{leancode}

\begin{dictionary}
The first statement is at paper index $m+1$ (see the index shift above) over any field of characteristic zero.
\emph{Caution, two declarations share this short name:} a second \lean{sourcePiece_finrank_le_product}, in
\leanpath{AffineSl2BoundaryVacuumCharacter.lean} (namespace \lean{VOA.AffineSl2BoundaryCharacter}), restates it at index
$m$ with \lean{hm : 1 ≤ m} over $\C$ with right-hand side \lean{PowerSeries.coeff N (product m)}; that second one is the
one used in \lean{coeff_character_le_product}. \lean{vacuum_finrank_le_sourcePiece} is the surjection argument: $\pi$ maps the weight-$N$ piece of
$\ltwentyB$ onto $\bigoplus_{p\le N}\ltwentygr_{p,N}V$, whose dimension is $\dim V_N$ by the telescope
\lean{sum_liWeightLayer_finrank_eq_grading} (each $V_N$ is filtered by $F^pV_N$, $F^{N+1}V_N=0$). The telescope holds
for every VOA, with no generation hypothesis; this is the form quoted in Chapter~\ref{ch:L11}.
\end{dictionary}

\textbf{Proof on paper.} The degeneration of Section~\ref{l20:sec:m1} + Theorem~\ref{l20:thm:auxbound} + the $q$-series identities of Section~\ref{l20:sec:m4}: $\dim(\ltwentyB)_N\le\dim(\ltwentyCm)_N\le[q^N]S_m=[q^N]\,
\mathrm{product}$. And $\pi$ restricted to weight $N$ is onto $\bigoplus_p\ltwentygr_{p,N}V$ (Theorem~\ref{l20:thm:m0arc}),
of dimension $\dim V_N$. \qed

\begin{physics}
Every quantum state has a classical shadow, and classical states are bounded by the quasi-particle count:
``Schur index $\le$ arc-space Hilbert series of the Higgs branch scheme'' (a general inequality, true for any VOA
generated by its weight-one space and any relations known to hold).
\end{physics}

\begin{status}
Proved for every $m\ge1$, unconditionally; standard axioms. No VOA input enters the bound on $\dim(\ltwentyB)_N$:
it is pure commutative algebra and combinatorics.
\end{status}

%======================================================================
\section{The lower bound for the character}\label{l20:sec:m6}
%======================================================================

This is where the library goes beyond the paper (which quotes Kac--Wakimoto). The lower bound
$[q^N]\,\mathrm{product}\le\dim L_{k_m,N}$ is proved by counting null vectors with the determinant.

\subsection{The idea}

Write $V^k_N=\bigoplus_cV^k_{N,c}$ by $h_0$-charge $2c$. On each finite block the invariant form has a Gram matrix
$G_{N,c}(k)$ in the PBW basis, whose entries are \emph{polynomials in the level} $k$. The simple quotient's block
$L_{k,N,c}$ has dimension $\operatorname{rank}G_{N,c}(k)$ (the maximal ideal is the radical of the form). Hence
\begin{equation}\label{l20:eq:rank}
\dim V^k_{N,c}-\dim L_{k,N,c}=\dim\ker G_{N,c}(k)\ \le\ \operatorname{ord}_{t=k}\det G_{N,c}(t),
\end{equation}
the inequality being a general fact about polynomial matrices (Lemma~\ref{l20:lem:rootorder}). Summing over $c$:
\begin{equation*}
\dim L_{k,N}\ \ge\ \dim V^k_N-\sum_c\operatorname{ord}_{t=k}\det G_{N,c}(t).
\end{equation*}
So it suffices to compute the total order of vanishing of the determinants at $k=k_m$ and show that
$\dim V^k_N-\text{(total order)}=[q^N]\,\mathrm{product}$. Since $\dim V^k_N=[q^N](q)_\infty^{-3}$ (PBW), and by
Jacobi $\mathrm{product}=(q)_\infty^{-3}\sum_{r\ge0}(-1)^r(2r+1)q^{p\binom{r+1}2}$, the claim is that the total order
is the \emph{Weyl correction}
\begin{equation}\label{l20:eq:weyl}
\sum_{s\ge1}\Bigl[(4s-1)\,P_3\bigl(N-ps(2s-1)\bigr)-(4s+1)\,P_3\bigl(N-ps(2s+1)\bigr)\Bigr],\qquad
P_3(n)=[q^n](q)^{-3}_\infty\ (P_3(n)=0\text{ for }n<0),
\end{equation}
which is $-\sum_{r\ge1}(-1)^r(2r+1)P_3(N-p\binom{r+1}2)$ with $r=2s-1$ and $r=2s$ grouped in pairs.

\subsection{The universal character and charge decomposition}

\begin{theorem}[PBW character and charge decomposition]\label{l20:thm:univchar}
For every non-critical $k$: $\dim V^k_N=[q^N](q)^{-3}_\infty$; $V^k_N$ and $L_{k,N}$ are the direct sums of their
charge-$2c$ pieces, $c\in[-N,N]$ (in Lean: the dimension identities below, which is all that is used); the charge-$2c$ piece of $V^k_N$ has the PBW basis indexed by \lean{JointIndex N (2*c)}.
\end{theorem}
\noindent Files \leanpath{VOA/\allowbreak Examples/\allowbreak AffineSl2UniversalCharacter.lean},
\leanpath{VOA/\allowbreak Examples/\allowbreak AffineSl2ChargeDecomposition.lean},
\leanpath{VOA/\allowbreak Examples/\allowbreak AffineSl2AdmissibleCharacterRange.lean}.
\begin{leancode}
theorem universal_grading_finrank_eq_reciprocalEuler_cube
    (k : ℂ) (hk : k+2≠0) (N : ℕ) :
    (Module.finrank ℂ ((universal k hk).grading N) : ℚ) =
      PowerSeries.coeff N ((reciprocalEuler : PowerSeries ℚ)^3) := by

def chargeWindow (N : ℕ) : Finset ℤ := Finset.Icc (-(N : ℤ)) N

theorem universal_grading_finrank_eq_sum_joint (k : ℂ) (hk : k+2≠0) (N : ℕ) :
    Module.finrank ℂ ((universal k hk).grading N) =
      ∑ c ∈ chargeWindow N,Fintype.card (JointIndex N (2*c)) := by

theorem simple_grading_finrank_eq_sum_joint (k : ℂ) (hk : k+2≠0) (N : ℕ) :
    Module.finrank ℂ ((simple k hk).grading N) =
      ∑ c ∈ chargeWindow N,Module.finrank ℂ (simpleJointPiece k hk N (2*c)) := by

def simpleJointPiece (k : ℂ) (hk : k + 2 ≠ 0) (n h : ℤ) :=
  (jointSpace k hk n h).map (QuotVOA.maxGradedIdeal (universal k hk)).mkQ
\end{leancode}

\begin{dictionary}
\lean{JointIndex n h} is the set of PBW occupation data (Chapter~\ref{ch:L08}) of energy $n$ and $h_0$-charge $h$;
\lean{Fintype.card} its cardinality. \lean{Finset.Icc a b} is $\{a,\dots,b\}$. Charges of current monomials are even
(\lean{currentCharge_even}), so the sum over $c\in[-N,N]$ of charge $2c$ covers everything. \lean{simpleJointPiece}
is the image in $L_k$ of the charge-$2c$, energy-$n$ piece.
\end{dictionary}

\textbf{Proof on paper.} PBW: $V^k$ has basis the ordered monomials in $e_{-n},h_{-n},f_{-n}$ ($n\ge1$), three colours of
bosonic oscillators of weight $n$, whose generating function is $\prod_n(1-q^n)^{-3}$. $L_0$ and $h_0$ commute and are
diagonal on PBW monomials; the quotient map respects the decomposition. \qed

\begin{physics}
Three free bosons: the universal theory has no null states at generic level, so its partition function is
$1/\eta^3$ up to $q^{c/24}$.
\end{physics}

\begin{status}
Proved for every non-critical level; standard axioms.
\end{status}

\subsection{Gram rank and determinant root order}

\begin{lemma}[Kernel dimension is at most the root order]\label{l20:lem:rootorder}
Let $M(t)$ be a square matrix of polynomials over a field with $\det M\neq0$, and $z$ a point. Then
$\dim\ker M(z)\le\operatorname{ord}_{t=z}\det M(t)$.
\end{lemma}
\noindent File \leanpath{VOA/\allowbreak Commutative/\allowbreak PolynomialMatrixRootOrder.lean}.
\begin{leancode}
theorem finrank_ker_le_rootMultiplicity (M : Matrix ι ι K[X]) (z : K)
    (hdet : M.det ≠ 0) :
    Module.finrank K (evaluate M z).mulVecLin.ker ≤ M.det.rootMultiplicity z := by
\end{leancode}

\begin{dictionary}
\lean{K[X]} is the polynomial ring; \lean{evaluate M z} is the entrywise evaluation (\leanpath{VOA/\allowbreak Core/\allowbreak PolynomialMatrixSimpleRoot.lean});
\lean{mulVecLin.ker} is the kernel of $v\mapsto M(z)v$; \lean{rootMultiplicity z} is the order of vanishing at $z$. The
hypothesis $\det M\neq0$ (as a polynomial) is needed: otherwise the order is not defined (Mathlib sets it to $0$).
\end{dictionary}

\textbf{Proof on paper.} Choose a constant invertible matrix $P$ whose first $d=\dim\ker M(z)$ columns are a basis of
$\ker M(z)$. The first $d$ columns of $M(t)P$ vanish at $t=z$, hence are divisible by $(t-z)$, so
$(t-z)^d\mid\det(M(t)P)=\det M(t)\cdot\det P$. \qed

\begin{physics}
A null state at level $k$ is a direction in which the Shapovalov form degenerates; each such direction costs at least
one power of $(t-k)$ in the Kac determinant. (Equality would need a Jantzen-filtration argument; only the inequality is
needed.)
\end{physics}

\begin{status}
Proved for all fields; standard axioms.
\end{status}

\begin{theorem}[Simple joint dimension and the root-order bound]\label{l20:thm:rootbound}
For every non-critical $k$ and every $N$: $\dim L_{k,N,c}=\operatorname{rank}G_{N,c}(k)$ and
\begin{equation*}
\dim V^k_N\le\dim L_{k,N}+\sum_{c\in[-N,N]}\operatorname{ord}_{t=k}\det G_{N,c}(t).
\end{equation*}
\end{theorem}
\noindent Files \leanpath{VOA/\allowbreak Examples/\allowbreak AffineSl2VacuumFirstBlock.lean},
\leanpath{VOA/\allowbreak Examples/\allowbreak AffineSl2SimpleCharacterRank.lean}.
\begin{leancode}
def vacuumChevalleyPolynomial (n : ℕ) (c : ℤ) :
    Matrix (AffineSl2GeneralRigidity.JointIndex (n : ℤ) (2*c))
      (AffineSl2GeneralRigidity.JointIndex (n : ℤ) (2*c)) ℂ[X] := fun i j =>
  MvPolynomial.eval₂Hom Polynomial.C ![X,Polynomial.C 0]
    (wordPolynomial (adjointWord ((vacuumLiftWord i.val).map lowerLetter)++
      (vacuumLiftWord j.val).map lowerLetter))

theorem vacuumChevalleyPolynomial_eval (k : ℂ) (n : ℕ) (c : ℤ)
    (i j : AffineSl2GeneralRigidity.JointIndex (n : ℤ) (2*c)) :
    (vacuumChevalleyPolynomial n c i j).eval k=
      contravariant k 0 (indexedWord k 0 (vacuumLiftWord i.val))
        (indexedWord k 0 (vacuumLiftWord j.val)) := by

def actualJointGram (k : ℂ) (n : ℕ) (c : ℤ) :
    Matrix (JointIndex n (2 * c)) (JointIndex n (2 * c)) ℂ := by
  classical
  exact (Polynomial.evalRingHom k).mapMatrix (vacuumChevalleyPolynomial n c)

theorem jointSynthesis_mem_maxIdeal_iff (k : ℂ) (hk : k + 2 ≠ 0)
    (n : ℕ) (c : ℤ) (r : JointIndex n (2 * c) → ℂ) :
    jointSynthesis k hk n c r ∈ QuotVOA.maxGradedIdeal (universal k hk) ↔
      (actualJointGram k n c).vecMulLinear r = 0 := by

theorem simpleJointPiece_finrank_eq_gramRank (k : ℂ) (hk : k + 2 ≠ 0)
    (n : ℕ) (c : ℤ) :
    Module.finrank ℂ (simpleJointPiece k hk n (2 * c)) = (actualJointGram k n c).rank := by
\end{leancode}
\noindent File \leanpath{VOA/\allowbreak Examples/\allowbreak AffineSl2BoundaryCharacterLowerBound.lean}.
\begin{leancode}
def vacuumRootOrderTotal (k : ℂ) (N : ℕ) : ℕ :=
  ∑ c ∈ chargeWindow N, (vacuumChevalleyPolynomial N c).det.rootMultiplicity k

theorem universal_grading_finrank_le_simple_add_rootOrder (k : ℂ) (hk : k+2 ≠ 0)
    (N : ℕ) :
    Module.finrank ℂ ((universal k hk).grading N) ≤
      Module.finrank ℂ ((simple k hk).grading N) + vacuumRootOrderTotal k N := by
\end{leancode}

\begin{dictionary}
\lean{vacuumChevalleyPolynomial n c} is $G_{n,c}(t)$: entry $(i,j)$ is the polynomial in the level variable \lean{X}
whose value at $k$ is the contravariant (Shapovalov) pairing \lean{contravariant k 0} of the $i$-th and $j$-th PBW
words in the Verma module of highest weight $0$ (\lean{vacuumChevalleyPolynomial_eval}); \lean{eval₂Hom} substitutes
level $=$ \lean{X} and highest weight $=0$ in a two-variable polynomial \lean{wordPolynomial}. ``Chevalley'' refers
to the adjoint convention $e^\dagger=f$. \lean{actualJointGram k n c} is $G_{n,c}(k)$: by its body, the entrywise evaluation at $k$ of the
same polynomial matrix whose determinant's root multiplicity is counted below (so rank and root order refer to one
matrix). \lean{jointSynthesis k hk n c r} is the vector $\sum_ir_i\cdot(\text{PBW word }i)$ in $V^k$;
\lean{jointSynthesis_mem_maxIdeal_iff} says it is in the maximal ideal iff $r$ is in the (row) kernel of $G$:
``maximal ideal $=$ radical of the form'' made concrete. \lean{vacuumRootOrderTotal k N} is
$\sum_c\operatorname{ord}_k\det G_{N,c}$. The final inequality is \eqref{l20:eq:rank} summed over charges (the
determinant is a non-zero polynomial, \lean{vacuumChevalleyPolynomial_det_ne_zero}, proved below).
\end{dictionary}

\textbf{Proof on paper.} Rank--nullity on each block, Lemma~\ref{l20:lem:rootorder}, and the charge decompositions of
Theorem~\ref{l20:thm:univchar}. \qed

\begin{physics}
The number of physical states at level $k$ is at least the number of states of the free theory minus the number
of zeros of the Kac determinant at $k$.
\end{physics}

\begin{status}
Proved for every non-critical level; standard axioms.
\end{status}

\subsection{The vacuum determinant and its zeros}

\paragraph{Kac--Kazhdan factors for affine $\ltwentysl$.} Let $M(\lambda)$ be the Verma module of
$\widehat{\ltwentysl}$ at level $k$ whose highest-weight vector has $h_0$-eigenvalue $\lambda$ (the
\emph{horizontal highest weight}). Its weight spaces are indexed by energy $n$ and charge; the Shapovalov form
on each weight space has a determinant, a polynomial in $(k,\lambda)$. Write $\alpha$ for the positive root of
$\ltwentysl$ (horizontal simple root) and $\delta$ for the imaginary root. The positive real roots are
$\beta=\alpha+j\delta$ ($j\ge0$) and $\beta=-\alpha+j\delta$ ($j\ge1$). The Kac--Kazhdan formula says that the
determinant is, up to a non-zero constant, a product over pairs $(\beta,r)$, $r\ge1$, of the affine-linear factors
\begin{equation*}
\phi_{\beta,r}(k,\lambda)=\begin{cases}\lambda+1+j(k+2)-r,&\beta=\alpha+j\delta,\\
-\lambda-1+j(k+2)-r,&\beta=-\alpha+j\delta,\end{cases}
\end{equation*}
each raised to a power $\mathrm{mult}_{\beta,r}$, the number of PBW states at the weight lowered by $r\beta$ (the
descendants of the singular vector created when $\phi_{\beta,r}=0$). (Imaginary roots $j\delta$ contribute factors
proportional to $j(k+2)$, independent of $\lambda$ and non-zero at non-critical level; they are absorbed into the
normalisation below.) The pair
$(\alpha,1)$ gives $\phi_{\alpha,1}=\lambda$, which vanishes \emph{identically} at $\lambda=0$: this is the
singular vector $f_0\ltwentyvac$, and the vacuum module is $V^k(\ltwentysl)=M(0)/\langle f_0\ltwentyvac\rangle$.
The submodule $\langle f_0\ltwentyvac\rangle$ is a copy of $M(-2)$ ($-2=s_\alpha\cdot0$, the reflected weight); its
highest-weight vector $f_0\ltwentyvac$ has charge $-2$, so the part of it of charge $2c$ is the part of $M(-2)$
of charge $2(c+1)$ \emph{relative to its own highest weight}. That is the origin of the index $c+1$ below.

Because $\det M(0)$ vanishes identically, one cannot simply divide determinants. The library instead removes from
the Kac--Kazhdan product \emph{exactly} the multiplicity of the identically vanishing factor $\phi_{\alpha,1}=\lambda$
(file docstring of \leanpath{AffineSl2KacVacuumCoefficient.lean}: ``removing exactly its descendant multiplicity leaves a
polynomial in central level''), and proves that what remains equals the Gram determinant of the quotient times the
determinant of $M(-2)$ at the shifted charge, up to a non-zero constant and an explicit normalising polynomial
$\ell_{n,c}(k)$, a product of ``derivative multipliers'' of PBW words (\lean{affineLeadingPolynomial}, non-zero at every
non-critical level by
\lean{affineLeadingPolynomial_eval_ne_zero}; the constant is \lean{vacuumDeterminantScale}, non-zero by
\lean{vacuumDeterminantScale_ne_zero}).

\begin{theorem}[Vacuum/reflected determinant identity]\label{l20:thm:vacdet}
For all $n,c$, as polynomials in the level $t$:
\begin{equation*}
\det G_{n,c}(t)\cdot D^{\mathrm{refl}}_{n,c+1}(t)=\mathrm{const}_{n,c}\cdot\ell_{n,c}(t)\cdot
\prod_{(\beta,r)\neq(\alpha,1)}\phi_{\beta,r}(t,0)^{\mathrm{mult}_{\beta,r}},
\end{equation*}
where $D^{\mathrm{refl}}_{n,c+1}(t)$ is the Shapovalov determinant of $M(-2)$ at energy $n$ and relative charge
$2(c+1)$, $\mathrm{const}_{n,c}\neq0$, $\ell_{n,c}(k)\neq0$ for every non-critical $k$, and the product runs over
the Kac--Kazhdan labels except the horizontal simple root with $r=1$. Consequently, for every non-critical $k$,
\begin{equation*}
\operatorname{ord}_k\det G_{n,c}+\mathrm{reflectedOrder}_k(n,c+1)=\mathrm{vacuumOrder}_k(n,c).
\end{equation*}
\end{theorem}
\noindent Files \leanpath{VOA/\allowbreak Examples/\allowbreak AffineSl2VacuumDeterminantOrder.lean},
\leanpath{VOA/\allowbreak Examples/\allowbreak AffineSl2KacVacuumCoefficient.lean},
\leanpath{VOA/\allowbreak Examples/\allowbreak AffineSl2KacProduct.lean}.
\begin{leancode}
theorem affine_vacuum_polynomial_identity_all (n : ℕ) (c : ℤ) :
    (vacuumChevalleyPolynomial n c).det * affineMinusTwoPolynomial n (c+1) =
      Polynomial.C (vacuumDeterminantScale n c) * affineLeadingPolynomial n c *
        affineVacuumProduct n c := by

theorem vacuumChevalleyPolynomial_det_ne_zero (n : ℕ) (c : ℤ) : (vacuumChevalleyPolynomial n c).det ≠ 0 := by

def vacuumOrder (k : ℂ) (n : ℕ) (c : ℤ) : ℕ :=
  ∑ i ∈ (affineKacLabels n c).erase ((1,0),0),
    if realRootWeight (k+2) i = 0 then affineKacMultiplicity n c i else 0

def reflectedOrder (k : ℂ) (n : ℕ) (c : ℤ) : ℕ :=
  ∑ i ∈ affineKacLabels n c,
    if realRootWeight (k+2) i = -2 then affineKacMultiplicity n c i else 0

theorem vacuum_determinant_order_add_reflected (k : ℂ) (hk : k+2 ≠ 0)
    (n : ℕ) (c : ℤ) :
    (vacuumChevalleyPolynomial n c).det.rootMultiplicity k + reflectedOrder k n (c+1) =
      vacuumOrder k n c := by

def realRootWeight (t : ℂ) (i : (Fin 3×ℕ)×ℕ) : ℂ :=
  (realRootSlope i.1 : ℂ)*t+(realRootOffset i.1 i.2 : ℂ)

def affineKacMultiplicity (n : ℕ) (c : ℤ) (i : (Fin 3×ℕ)×ℕ) : ℕ :=
  Fintype.card (JointIndex (n-(i.2+1)*oscillatorEnergy i.1)
    (c-((i.2+1 : ℕ) : ℤ)*oscillatorCharge i.1))
\end{leancode}

\begin{dictionary}
A Kac label \lean{i = ((col, j), r)} encodes a positive real root $\beta$ (through an oscillator \lean{(col, j)}) and a
multiplicity index \lean{r} (the factor for $\lean{r}+1$ copies of $\beta$, i.e.\ the book's $r$ is \lean{r}${}+1$):
colour \lean{1} with depth \lean{j} is $\beta=\alpha+j\delta$, colour \lean{0} with depth \lean{j} is $\beta=-\alpha+(j+1)\delta$
(\lean{realRootSlope}, \lean{realRootOffset} in \leanpath{AffineSl2KacProduct.lean}). \lean{realRootWeight (k+2) i} is the
affine-linear function $\mathrm{slope}\cdot(k+2)+\mathrm{offset}$, and \lean{realRootWeight (k+2) i} $=\lambda$ is exactly
$\phi_{\beta,r}(k,\lambda)=0$: e.g.\ for $\beta=\alpha+j\delta$ the slope is $-j$, the offset \lean{r}, and
$\lambda=\lean{r}-j(k+2)$ iff $\lambda+1+j(k+2)-(\lean{r}+1)=0$. So the condition ``$=0$'' (resp.\ ``$=-2$'') is the
Kac--Kazhdan condition for the vacuum (resp.\ reflected) Verma module.
\lean{affineKacMultiplicity} is the number of PBW states at the shifted energy and charge (the descendants of the
singular vector). \lean{((1,0),0)} is the label of the horizontal simple root, removed because the vacuum module
$V^k=M(k\Lambda_0)/\langle f_0\ltwentyvac\rangle$ is already the quotient by it. \lean{vacuumOrder}, \lean{reflectedOrder}
count the resonances with multiplicity. The identity is an identity of polynomials in $t$, so it holds even at levels where
the reflected determinant also vanishes; this is why it can be used at $k=k_m$ where many factors collide.
\end{dictionary}

\textbf{Proof on paper.} The vacuum module is the quotient of the Verma module $M(0)$ (horizontal highest weight $0$) by
the submodule generated by $f_0\ltwentyvac$, which is isomorphic to the Verma module $M(-2)$. Since the Shapovalov
determinant of $M(0)$ itself is identically zero, one works with the Kac--Kazhdan product with the factor
$\phi_{\alpha,1}=\lambda$ removed to exactly its multiplicity, and shows it equals (Gram determinant of the
quotient)$\times$(determinant of $M(-2)$ at the shifted charge)$\times$(a change-of-basis constant)$\times\ell_{n,c}$,
using an adapted PBW basis compatible with the submodule
(\leanpath{AffineSl2ParabolicAdaptedBasis.lean}, \leanpath{AffineSl2ParabolicFirstBlock.lean}). The Kac--Kazhdan formula
for both Verma determinants (proved in the library for affine $\ltwentysl$, \leanpath{AffineSl2KacProduct.lean} and its
companions) then gives the identity.
Taking orders of vanishing (orders of products add) gives the second statement. \qed

\begin{physics}
The vacuum module is a Verma module with its trivial horizontal null vector $f_0\ltwentyvac$ removed; its null states are
those of the Verma module minus those of the removed submodule (the ``reflected'' Verma module). This is the first step of
the BGG / Kac--Wakimoto resolution, done at the level of determinants.
\end{physics}

\begin{status}
Proved for every $n,c$ and every non-critical $k$; standard axioms.
\end{status}

\subsection{The boundary resonances and the Weyl correction}

\begin{theorem}[Total root order at the boundary level]\label{l20:thm:resonances}
Let $P$ be odd and $k=2/P-2$ (we write $P$, as the Lean statements do, for an arbitrary odd denominator; at
$P=2m+1$ it is the chapter's $p$ and $k=k_m$). For every $N$, the total root order equals the Weyl correction \eqref{l20:eq:weyl}, and
\begin{equation*}
\mathrm{vacuumRootOrderTotal}(k,N)+[q^N]\frac{(q^P;q^P)^3_\infty}{(q)^3_\infty}=\dim V^k_N.
\end{equation*}
\end{theorem}
\noindent Files \leanpath{VOA/\allowbreak Examples/\allowbreak AffineSl2BoundaryVacuumResonances.lean},
\leanpath{VOA/\allowbreak Examples/\allowbreak AffineSl2ParabolicChargeSum.lean}.
\begin{leancode}
def lowerEnergy (P s : ℕ) : ℕ := P*s*(2*s-1)
def upperEnergy (P s : ℕ) : ℕ := P*s*(2*s+1)

theorem vacuumOrder_boundary (P : ℕ) (hP : 0<P) (hcop : P.Coprime 2)
    (n : ℕ) (c : ℤ) :
    Wak.Charged.vacuumOrder (2/(P:ℂ)-2) n c =
      ∑ s∈Finset.Icc 1 n, (shiftedJointCount n c (lowerEnergy P s) (2*(s:ℤ)-1) +
        shiftedJointCount n c (upperEnergy P s) (-(2*(s:ℤ)+1))) := by

def universalPBWCount (N : ℕ) : ℕ := Nat.card (EnergyIndex N)

def shiftedPBWCount (N energy : ℕ) : ℕ :=
  if energy≤N then universalPBWCount (N-energy) else 0

theorem sum_boundary_order_difference (P : ℕ) (hP : 0<P) (hcop : P.Coprime 2) (N : ℕ) :
    (∑ c ∈ chargeWindow N,
      ((Wak.Charged.vacuumOrder (2/(P:ℂ)-2) N c : ℤ) -
        Wak.Charged.reflectedOrder (2/(P:ℂ)-2) N (c+1))) =
      ∑ s ∈ Finset.Icc 1 N,
        ((4*(s : ℤ)-1)*(shiftedPBWCount N
          (AffineSl2BoundaryVacuumResonances.lowerEnergy P s) : ℤ) -
         (4*(s : ℤ)+1)*(shiftedPBWCount N
          (AffineSl2BoundaryVacuumResonances.upperEnergy P s) : ℤ)) := by
\end{leancode}
\noindent Files \leanpath{VOA/\allowbreak Commutative/\allowbreak QWeylPairing.lean},
\leanpath{VOA/\allowbreak Examples/\allowbreak AffineSl2BoundaryDeterminantCharacter.lean}.
\begin{leancode}
def weylPairTail (p n : ℕ) : K :=
  ∑ s ∈ Finset.Icc 1 n,
    ((if p*s*(2*s-1) ≤ n then
      (4*s-1 : ℕ) * PowerSeries.coeff (n-p*s*(2*s-1)) (reciprocalEuler^3 : PowerSeries K)
      else 0) -
    (if p*s*(2*s+1) ≤ n then
      (4*s+1 : ℕ) * PowerSeries.coeff (n-p*s*(2*s+1)) (reciprocalEuler^3 : PowerSeries K)
      else 0))

theorem weylPairTail_add_product (p : ℕ) (hp : p ≠ 0) (n : ℕ) :
    (weylPairTail p n : K) + PowerSeries.coeff n
        ((PowerSeries.expand p hp (euler : PowerSeries K))^3 * reciprocalEuler^3) =
      PowerSeries.coeff n (reciprocalEuler^3 : PowerSeries K) := by

theorem vacuumRootOrderTotal_eq_weylPairTail (P : ℕ) (hP : 0<P)
    (hcop : P.Coprime 2) (N : ℕ) :
    (vacuumRootOrderTotal (2/(P : ℂ)-2) N : ℚ) = weylPairTail P N := by

theorem vacuumRootOrderTotal_add_product_eq_universal_finrank (P : ℕ) (hP : 0<P)
    (hcop : P.Coprime 2) (hk : 2/(P : ℂ)-2+2 ≠ 0) (N : ℕ) :
    (vacuumRootOrderTotal (2/(P : ℂ)-2) N : ℚ) +
      PowerSeries.coeff N
        ((PowerSeries.expand P (Nat.ne_of_gt hP) (euler : PowerSeries ℚ))^3 * reciprocalEuler^3) =
      (Module.finrank ℂ ((universal (2/(P : ℂ)-2) hk).grading N) : ℚ) := by
\end{leancode}

\begin{dictionary}
\lean{P.Coprime 2} says $P$ is odd; the level is written \lean{2/(P:ℂ)-2}, which equals \lean{level m} when
$P=2m+1$ (the headline file makes the conversion, \lean{hlevel}). \lean{shiftedJointCount n c E Q} counts PBW states of
energy $n-E$ and the correspondingly shifted charge; \lean{EnergyIndex N} is the set of all PBW states of energy $N$, so
\lean{universalPBWCount N} $=P_3(N)$ (\lean{universalPBWCount_eq_coeff}). The \lean{if … then … else 0} guards
implement $P_3(n)=0$ for $n<0$ (natural-number subtraction would otherwise silently give $P_3(0)$; the guards are
exactly what prevents that trap). \lean{weylPairTail p n} is \eqref{l20:eq:weyl}; \lean{weylPairTail_add_product} is the
identity ``Weyl correction $+$ product $=P_3$'', which is Jacobi's identity for $(q^p;q^p)^3_\infty$ read coefficientwise;
the last theorem combines it with the universal character. The range $s\in[1,N]$ is large enough since
$ps(2s-1)>N$ for $s>N$.
\end{dictionary}

\textbf{Proof on paper.} At $k+2=2/P$ with $P$ odd, a Kac--Kazhdan factor vanishes iff an integrality condition holds;
coprimality of $P$ and $2$ forces the relevant root energies to be multiples of $P$ and ties the multiplicity
index to them (compare \leanpath{AffineSl2BoundaryParabolicRoots.lean}: ``the next affine root has relative energy
$P\cdot r$''); the library enumerates the solutions exactly (\lean{vacuum_labels}, \lean{reflected_labels}). Solving, the vacuum
resonances come in two families with energy shifts $Ps(2s-1)$ and $Ps(2s+1)$, and the reflected resonances in the
``reversed'' families (\lean{vacuumOrder_boundary}, \lean{reflectedOrder_boundary}). The differences, summed over all
charges $c$, combine into $(2r+1)$ full copies of the PBW count at the shifted energy, with
$2r+1=4s-1$ and $4s+1$ (the dimensions of the $\ltwentysl$-multiplets of spin $2s-1$ and $2s$); see
\lean{sum_shiftedParabolicCount}. This is \lean{sum_boundary_order_difference}. Combining with
Theorem~\ref{l20:thm:vacdet} gives the total root order, and Jacobi's identity finishes. \qed

\begin{physics}
At the boundary level the singular vectors of the vacuum module sit at the affine Weyl orbit $w\cdot0$: an
$\ltwentysl$-multiplet of spin $r$ at energy $p\binom{r+1}2$, with sign $(-1)^r$. The character is the alternating sum
$\sum_r(-1)^r(2r+1)q^{p\binom{r+1}2}/(q)^3_\infty$ (the unflavoured Kac--Wakimoto / Weyl--Kac formula), and Jacobi's
triple product turns it into $(q^p;q^p)^3_\infty/(q)^3_\infty$. In the lower-bound step this formula is not assumed: only its inequality
$\ge$ is derived, from determinant zeros. For $m=1$, $N=9$: $3P_3(6)-5P_3(0)=663-5=658$ (Example~\ref{l20:ex:numbers}).
\end{physics}

\begin{status}
Proved for every odd $P$; standard axioms. The supporting parabolic-module theory in the umbrella
(\leanpath{AffineSl2Parabolic*.lean}, \leanpath{AffineSl2HorizontalReflection.lean}: exact parabolic quotient
sequence, adapted bases, joint dimensions) is proved and integrated; the campaign ledger notes that the final
bound uses the determinant argument, not an exact resolution.
\end{status}

\begin{theorem}[The lower bound]\label{l20:thm:lower}
For $m\ge1$ and all $N$: $[q^N]\,\mathrm{product}(m)\le\dim L_{k_m}(\ltwentysl)_N$.
\end{theorem}
\noindent This is \lean{coeff_product_le_character}, quoted in Theorem~\ref{l20:thm:character}
(file \leanpath{VOA/\allowbreak Examples/\allowbreak AffineSl2BoundaryMainTheorems.lean}).

\textbf{Proof on paper.} By Theorem~\ref{l20:thm:rootbound},
$\dim V^k_N\le\dim L_{k,N}+\mathrm{order}_N$, and by Theorem~\ref{l20:thm:resonances}
$\mathrm{order}_N+[q^N]\mathrm{product}=\dim V^k_N$. Subtract. (In Lean the only extra work is converting
\lean{2/(denominator m)-2} to \lean{level m}.) \qed

\begin{physics}
Combined with the Hilbert-series upper bound: the null states found by the determinant are \emph{all} the null states, a fact that is usually
proved via the Kazhdan--Lusztig / Kac--Wakimoto machinery; here it falls out of the squeeze.
\end{physics}

\begin{status}
Unconditional; standard axioms.
\end{status}

%======================================================================
\section{The conformal Hilbert squeeze}\label{l20:sec:m7}
%======================================================================

\begin{theorem}[Injectivity from the dimension squeeze]\label{l20:thm:squeeze}
Let $W$ be a VOA, $x_1,x_2,x_3\in W_1$, and suppose the arc relations of $I_m$ vanish under
$\pi\colon x_{i,n}\mapsto\partial^n\bar x_i/n!$ and that $\pi$ is surjective. If
$\dim(\ltwentyB)_N\le\dim W_N$ for all $N$, then $\pi\colon\ltwentyB\to\ltwentygr W$ is injective. For
$W=L_{k_m}(\ltwentysl)$ the dimension hypothesis follows from the character equality.
\end{theorem}
\noindent Files \leanpath{VOA/\allowbreak Filtered/\allowbreak BoundaryArcConformalSqueeze.lean},
\leanpath{VOA/\allowbreak Core/\allowbreak FiniteHilbertSqueeze.lean}.
\begin{leancode}
abbrev Layers (N : ℕ) := ∀ p : Fin (N + 1), liWeightLayer W p.val N

theorem layers_finrank (N : ℕ) :
    Module.finrank ℂ (Layers W N) = Module.finrank ℂ (W.grading N) := by

def pieceMap (N : ℕ) : sourcePiece (K := ℂ) m N →ₗ[ℂ] Layers W N :=
  LinearMap.pi (fun p =>
    (DirectSum.component ℂ (ℕ × ℕ) (fun d => liWeightLayer W d.1 d.2) (p.val, N)).comp
      ((sourceMap W gen hgen m hI).toLinearMap.comp (sourcePiece (K := ℂ) m N).subtype))

theorem pieceMap_injective_of_hilbert_bound
    (hsurj : Function.Surjective (polynomialMap W gen hgen)) (N : ℕ)
    (hbound : Module.finrank ℂ (sourcePiece (K := ℂ) m N) ≤
      Module.finrank ℂ (W.grading N)) :
    Function.Injective (pieceMap W gen hgen m hI N) := by

theorem sourceMap_injective_of_hilbert_bound
    (hsurj : Function.Surjective (polynomialMap W gen hgen))
    (hbound : ∀ N, Module.finrank ℂ (sourcePiece (K := ℂ) m N) ≤
      Module.finrank ℂ (W.grading N)) :
    Function.Injective (sourceMap W gen hgen m hI) := by

theorem injective_of_surjective_of_finrank_le
    (f : A →ₗ[K] B) (hsurj : Function.Surjective f)
    (hle : Module.finrank K A ≤ Module.finrank K B) :
    Function.Injective f := by
\end{leancode}
\noindent Files \leanpath{VOA/\allowbreak Examples/\allowbreak AffineSl2BoundaryVacuumCharacter.lean},
\leanpath{VOA/\allowbreak Examples/\allowbreak AffineSl2BoundaryClassicalFreeness.lean}.
\begin{leancode}
theorem sourcePiece_finrank_le_vacuum_of_character (m : ℕ) (hm : 1 ≤ m)
    (hcharacter : character m = product m) (N : ℕ) :
    Module.finrank ℂ (sourcePiece (K := ℂ) m N) ≤
      Module.finrank ℂ ((vacuum m).grading N) := by

theorem canonicalMap_bijective_of_character (m : ℕ) (hm : 1 ≤ m)
    (hcharacter : character m = product m) :
    Function.Bijective (canonicalMap m hm) := by
\end{leancode}

\begin{dictionary}
In \leanpath{BoundaryArcConformalSqueeze.lean}, \lean{W}, \lean{gen}, \lean{hgen} (weight-one generators), \lean{m} and
\lean{hI} (the source ideal lies in the kernel) are section variables, and \lean{polynomialMap W gen hgen} is that file's
own \lean{abbrev} for \lean{liInfiniteJetPolynomialMap W gen (fun _ => 1) hgen} (Definition~\ref{l20:def:canonical});
it is a different declaration from \lean{BoundaryArcLiMap.polynomialMap C}, which is the special case
$W=$ \lean{simpleVOA C}, \lean{gen} $=$ \lean{gen C}. \lean{Layers W N} is the finite product
$\prod_{p\le N}\ltwentygr_{p,N}W$, of total dimension $\dim W_N$ (\lean{layers_finrank}, from the telescope of
Theorem~\ref{l20:thm:upper}). \lean{pieceMap N} is $\pi$ restricted to the weight-$N$ piece, followed by projection to the
layers of weight $N$. \lean{injective_of_surjective_of_finrank_le} is the linear-algebra fact used (in namespace
\lean{FiniteHilbertSqueeze}, Chapter~\ref{ch:L11}). \lean{..._of_character} are the helpers that take the character
equality as an explicit argument; the headline \lean{canonicalMap_bijective} supplies it.
\end{dictionary}

\textbf{Proof on paper.} Fix $N$. Surjectivity of $\pi$ and homogeneity give that \lean{pieceMap N} is onto $\prod_p\ltwentygr_{p,N}W$,
a space of dimension $\dim W_N\ge\dim(\ltwentyB)_N$; a surjective linear map from a space of dimension at most that of the
target is injective. Globally: if $\pi(f)=0$ for a polynomial $f$, then each conformal component $f_N$ satisfies
$\mathrm{pieceMap}_N([f_N])=0$ (the conformal projection is compatible with the Li components,
\lean{component_conformal}), hence $f_N$ is in the relation ideal, hence so is $f=\sum_Nf_N$. For
$W=L_{k_m}$: $\dim(\ltwentyB)_N\le[q^N]\mathrm{product}=[q^N]\mathrm{character}=\dim W_N$. \qed

\begin{physics}
Equal partition functions plus a surjection of state spaces in every energy forces an isomorphism: the classical
and quantum Hilbert spaces have the same size level by level.
\end{physics}

\begin{status}
Proved; standard axioms. The helpers are conditional on \lean{hcharacter}, discharged in the headline file.
\end{status}

%======================================================================
\section{Poisson vertex structure, inverse, and the \texorpdfstring{$C_2$}{C2} presentation}\label{l20:sec:m8}
%======================================================================

\begin{theorem}[The canonical map is a PVA homomorphism; uniqueness]\label{l20:thm:pvahom}
$\pi$ commutes with $\partial$ and with every $(n)$-product. Moreover two PVA structures on a presented arc algebra
with the standard translation that agree on pairs of order-zero generators are equal, and a translation-compatible
algebra map that is compatible with the modes on order-zero generators is compatible with all modes.
\end{theorem}
\noindent File \leanpath{VOA/\allowbreak Commutative/\allowbreak InfiniteJetPoissonUniqueness.lean}.
\begin{leancode}
theorem compatible_of_base_generators
    (hP : P.translation = translation I)
    (hT : ∀ x, F (P.translation x) = Q.translation (F x))
    (hbase : ∀ i j, Compatible P Q F (generator I i 0) (generator I j 0)) :
    ∀ x y, Compatible P Q F x y := by

theorem poissonVertexAlgebra_unique
    (Q : VOA.PoissonVertexAlgebra (PresentedJet I))
    (hP : P.translation = translation I) (hQ : Q.translation = translation I)
    (hbase : ∀ i j n, P.mode n (generator I i 0) (generator I j 0) =
      Q.mode n (generator I i 0) (generator I j 0)) : P = Q := by
\end{leancode}
\noindent File \leanpath{VOA/\allowbreak Commutative/\allowbreak PoissonVertexHom.lean}.
\begin{leancode}
def Compatible (P : VOA.PoissonVertexAlgebra A) (Q : VOA.PoissonVertexAlgebra B)
    (F : A →ₐ[ℂ] B) (x y : A) : Prop :=
  ∀ n, F (P.mode n x y) = Q.mode n (F x) (F y)
\end{leancode}

\begin{dictionary}
\lean{Compatible P Q F x y}: $F(x_{(n)}y)=F(x)_{(n)}F(y)$ for all $n$. \lean{generator I i 0} is the class of $x_{i,0}$.
The first theorem reduces homomorphism checks to generators, using Leibniz (products), sesquilinearity (derivatives) and
skew-symmetry (swapping slots). For \lean{arcLiPoissonHom} the generator check is the computation
$\overline{x_i}{}_{(0)}\overline{x_j}=\overline{[x_i,x_j]}$ and $\overline{x_i}{}_{(n+1)}\overline{x_j}=0$ in $\ltwentygr L_k$
(the level term has negative Li degree).
\end{dictionary}

\textbf{Proof on paper.} Every element of the arc algebra is a polynomial in the $x_{i,n}=\partial^nx_{i,0}/n!$; the PVA laws
express $(\partial^ra)_{(n)}b$, $a_{(n)}(bc)$, $b_{(n)}a$ in terms of modes of the factors, so compatibility propagates
from generators to everything. \qed

\begin{physics}
A classical field theory is determined by its basic Poisson brackets $\{x_i(z),x_j(w)\}$; everything else follows from the
Leibniz rule and derivatives.
\end{physics}

\begin{status}
Proved; standard axioms.
\end{status}

\begin{theorem}[Injectivity on the PVA side; the inverse; the $C_2$ presentation]\label{l20:thm:m8}
Under the character equality: \lean{arcLiPoissonHom} is injective, hence a PVA equivalence (with compatible inverse);
the equivalence and its inverse preserve conformal pieces; the induced map $\ltwentyA\to R_V$ is injective, hence an
isomorphism.
\end{theorem}
\noindent Files \leanpath{VOA/\allowbreak Examples/\allowbreak AffineSl2BoundaryClassicalFreeness.lean},
\leanpath{VOA/\allowbreak Filtered/\allowbreak ArcBaseExactness.lean}.
\begin{leancode}
theorem arcLiPoissonHom_injective_of_character (m : ℕ) (hm : 1 ≤ m)
    (hcharacter : character m = product m) :
    Function.Injective (uniformCandidate m hm).arcLiPoissonHom.toAlgHom := by

theorem canonicalMap_conformal_iff (m : ℕ) (hm : 1 ≤ m)
    (hcharacter : character m = product m) (N : ℕ) (x : Source (R := ℂ) m) :
    canonicalMap m hm x ∈ BoundaryArcConformalSqueeze.liConformalPiece
        (AffineSl2BoundaryCharacter.vacuum m) N ↔
      x ∈ sourcePiece (K := ℂ) m N :=

theorem baseC2Map_injective_of_arc_injective
    (hLi : I ≤ RingHom.ker (MvPolynomial.aeval
      (fun i => liDegreeZeroClass W (weight i) (gen i))).toRingHom)
    (hinj : Function.Injective (liInfinitePresentedJetMap W
      (fun i => liDegreeZeroClass W (weight i) (gen i)) I hLi)) :
    Function.Injective (baseC2Map W weight gen I hI) := by
\end{leancode}
\noindent File \leanpath{VOA/\allowbreak Filtered/\allowbreak C2LiAlgebra.lean}.
\begin{leancode}
theorem c2ToLiAlgebra_aeval {σ : Type*} (weight : σ → ℕ)
    (gen : ∀ i, W.grading (weight i : ℤ)) (f : MvPolynomial σ ℂ) :
    c2ToLiAlgebra W (MvPolynomial.aeval (fun i => mk₂ W (gen i).1) f) =
      MvPolynomial.aeval (fun i => liDegreeZeroClass W (weight i) (gen i)) f := by
\end{leancode}

\begin{dictionary}
\lean{arcLiPoissonHom_injective_of_character} transfers injectivity from \lean{canonicalMap} (on \lean{Source m}) to
\lean{arcLiPoissonHom} (on \lean{PresentedJet (baseIdeal m)}) through \lean{sourceIdeal_eq_jetIdeal}.
\lean{canonicalMap_conformal_iff} is the conformal statement on the \lean{Source} side; the headline
\lean{arcPoissonEquiv_conformal_iff} is its transport. \lean{c2ToLiAlgebra W : A₂ W →ₐ[ℂ] liAssociatedGraded W} is the
inclusion of $R_W=F^0W/F^1W$ as the degree-zero part of $\ltwentygr W$, as an algebra map; \lean{c2ToLiAlgebra_aeval} says it
sends the $C_2$ class of a polynomial in generators to the same polynomial in degree-zero Li classes. In
\lean{baseC2Map_injective_of_arc_injective} the section variables are a VOA \lean{W}, weights, generators \lean{gen}, an
ideal \lean{I} and \lean{hI : I ≤ ker(C₂ evaluation)}.
\end{dictionary}

\textbf{Proof on paper.} Injectivity transfers along the identification of the two names of $\ltwentyB$. For the presentation:
the square ``$\ltwentyA\to\ltwentyJ\ltwentyA\xrightarrow\pi\ltwentygr V$'' $=$ ``$\ltwentyA\to R_V\xrightarrow{c_2\to\mathrm{Li}}\ltwentygr V$''
commutes (\lean{baseC2Map_li_commutes}); the left composite is injective (constant arcs have a retraction, then $\pi$),
so $\ltwentyA\to R_V$ is injective. Conformal compatibility: $\pi$ maps the weight-$N$ piece \emph{onto} the weight-$N$ piece
(\lean{sourcePiece_map_eq_liConformalPiece}), and injectivity gives the converse inclusion. \qed

\begin{physics}
The classical limit determines the chiral ring: once the full arc space is known to be faithful, its constant part ---
the Higgs branch chiral ring --- has no further relations.
\end{physics}

\begin{status}
Proved; the \lean{_of_character} forms carry \lean{hcharacter}, discharged in the headline file. Standard axioms.
\end{status}

%======================================================================
\section{The Virasoro companion: classical freeness of the minimal models \texorpdfstring{$M(2,2m+3)$}{M(2,2m+3)}}\label{l20:sec:vir}
%======================================================================

This section and the next document three files that prove the Virasoro analogue of the theorems above, and a family of
commutative-algebra files left from an alternative route to the affine theorem that was not completed. (Their Lean
quotations were first copied on 2026-10-27 and are checked against the released source by \texttt{sync.py}; see
Section~\ref{l20:sec:verification}.)

\textbf{Machine status in these two sections.} ``Proved, unconditional'' in a status box below means: the Lean statement has
no hypothesis beyond the displayed variables, and \texttt{\#print axioms} on it gives Lean's three standard axioms. These
axiom checks are \emph{isolated-module} runs (\texttt{lake env lean -DmaxSynthPendingDepth=3} on a scratch file importing one
compiled module; Section~\ref{l20:sec:verification}). The aggregate \texttt{lake build} was \emph{not} re-run for these
sections, and the project counts a result as established only when the aggregate build is green; each status box repeats the
caveat ``(isolated-module check; aggregate build not re-run)''.

\subsection{What is proved, and why it belongs here}

The earlier sections proved (the Arakawa--Dai theorem) that the simple affine VOA $L_{k_m}(\ltwentysl)$ at the boundary
admissible level $k_m = -2+\frac{2}{2m+1}$ is classically free. This section proves the analogous statement for the Virasoro
minimal models of the series $(2,2m+3)$: $M(2,5)$ (Lee--Yang), $M(2,7)$, $M(2,9)$, and so on. For them the $C_2$-algebra
is a single fat point,
\begin{equation*}
R_{L(c_{2,2m+3},0)}\cong\C[T]/(T^{m+1}),\qquad T=[\omega],
\end{equation*}
and classical freeness says that Li's associated graded is the arc algebra of this fat point,
\begin{equation}\label{l20:eq:virfree}
\ltwentyJ\bigl(\C[T]/(T^{m+1})\bigr)\xrightarrow{\ \cong\ }\ltwentygr L(c_{2,2m+3},0).
\end{equation}
On characters, \eqref{l20:eq:virfree} is equivalent to the Andrews--Gordon identities (the Hilbert series of the
left side is the Feigin--Stoyanovsky fermionic sum, the character of the right side is the Rocha-Caridi product);
this is the result of van Ekeren and Heluani. The library proves \eqref{l20:eq:virfree} for every $m\ge0$, together with
the presentation of the full non-reduced $C_2$-algebra and the Poisson vertex structure, about the
\emph{actual} simple quotient of the \emph{actual} universal Virasoro VOA. The structure of the statements is the
same as for $L_{k_m}(\ltwentysl)$ in Section~\ref{l20:sec:headlines}, which is why the files are documented here.
\textbf{The role of character identities.} The injectivity proof uses the Andrews--Gordon identity \eqref{l20:eq:ag} and
the Feigin--Fuchs/Rocha-Caridi lower bound for $\dim L(c_m,0)_N$ (both proved in Lean), but only through the leading-symbol
theorem \lean{virSymbol_twoSeries} (Theorem~\ref{l20:thm:virinputs}). In contrast to the affine case, no comparison of the
Hilbert series of the two sides of \eqref{l20:eq:virfree} is made at the end; the character of $L(c_m,0)$ is a by-product.

\textbf{The Andrews--Gordon identity used here.} For $n\ge0$ and $N\ge0$,
\begin{equation}\label{l20:eq:ag}
\#\Bigl\{\text{partitions of }N\text{ into parts}\ge2\text{ with }f_j+f_{j+1}\le n\text{ for all }j\Bigr\}
=\sum_{k\in\Z}(-1)^k\,p\Bigl(N-\tfrac{k((2n+3)k+2n+1)}{2}\Bigr),
\end{equation}
where $f_j$ is the multiplicity of the part $j$ and $p$ is the partition function ($p(<0)=0$). The right side is the
coefficient of $q^N$ in the Rocha-Caridi character of $L(c_{2,2n+3},0)$; by Gordon's theorem both sides also equal the number
of partitions of $N$ into parts $\not\equiv0,\pm1\pmod{2n+3}$, the product form
$\prod_{j\not\equiv0,\pm1\,(2n+3)}(1-q^j)^{-1}$ of the character. In the library this is
\lean{TwoSeriesCharacter.AG.gordon_sum_eq_rcq} (\leanpath{VOA/\allowbreak Examples/\allowbreak TwoSeriesCharacterAGIdentity.lean}).
(Check at $n=1$, $N=4$: the left side counts $\{4\}$ only, and the right side is $p(4)-p(3)-p(0)=5-3-1=1$.)

\textbf{Gordon-good words.} A creation word $L_{-j_1-2}\cdots L_{-j_b-2}\ltwentyvac$ with $j_1\ge\dots\ge j_b$ (sorted
decreasingly) is \emph{Gordon-good of order $n$} if $j_i\ge j_{i+n}+2$ whenever both entries exist (\lean{GordonGood n},
\leanpath{VOA/\allowbreak Examples/\allowbreak TwoSeriesCharacterSpanning.lean}). Via ``depth $j$ $\leftrightarrow$ part $j+2$'' the
sorted Gordon-good words of weight $N$ correspond exactly to the partitions counted on the left of \eqref{l20:eq:ag}: the
condition $j_i\ge j_{i+n}+2$ says that no $n+1$ parts lie within distance one of each other, i.e.\ $f_j+f_{j+1}\le n$. For
$M(2,2m+3)$ the order is $n=m$; for Lee--Yang ($m=1$) these are the Rogers--Ramanujan words, parts differing by at least $2$.

\textbf{Index convention.} Throughout this section $m\ge0$ and $M(2,2m+3)$ is the minimal model with
\begin{equation*}
c_m=c_{2,2m+3}=1-\frac{3(2m+1)^2}{2m+3},\qquad c_0=0,\quad c_1=-\tfrac{22}{5},\quad c_2=-\tfrac{68}{7}.
\end{equation*}
So $m=0$ is the trivial endpoint $M(2,3)$ ($c=0$, $L(0,0)=\C$, $R=\C[T]/(T)=\C$), and $m=1$ is Lee--Yang
$M(2,5)$. This $m$ is \emph{not} the $m$ of the affine theorem: there $m\ge1$, $p=2m+1$ and the level is
$k_m=-2+2/p$; here the minimal model is $M(2,2m+3)$ and the relation is $T^{m+1}$. The same letter thus labels different
algebras in the two halves of the chapter; the coincidence is the library's. (Chapter~\ref{ch:L11} writes the same series as
$M(2,2n+1)$ with $n=m+1$.) In four-dimensional language $M(2,2m+3)$ is the chiral algebra of the $(A_1,A_{2m})$
Argyres--Douglas theory.

We rely on Chapter~\ref{ch:L06} for the universal Virasoro VOA $V(c,0)$ (the library's \lean{virasoroVOA c},
Definition~\ref{l06:def:virVOA}), its creation words $L_w\ltwentyvac=L_{-j_1-2}\cdots L_{-j_b-2}\ltwentyvac$
(\lean{negWordVir c w}, with entry $j$ standing for $L_{-j-2}$, Definition~\ref{l06:def:grading}), the fact that these
words span (\lean{negSpanVir_eq_top}), and the simple quotient $L(c,0)=V(c,0)/I_{\max}(c)$ (\lean{virasoroSimple c},
\lean{maxIdealVir c}, Definition~\ref{l06:def:simple}). Li's filtration \lean{liFiltration W p}, its layers and
\lean{liAssociatedGraded W} are those of Chapter~\ref{ch:L11} (Section~\ref{l11:sec:li}); presented jet algebras are
those of Chapter~\ref{ch:L11} (Section~\ref{l11:sec:jets}) and Definition~\ref{l20:def:jet}. The generic
generator-word machinery is Section~\ref{l11:sec:genwords} of Chapter~\ref{ch:L11}: Definition~\ref{l11:def:genwordpiece}
(the word filtration $C_p$ for generators $a^1,\dots,a^g$), Theorem~\ref{l11:thm:tracked} (word spanning plus a tracked
estimate give $F^p=C_p$) and Theorem~\ref{l11:thm:genwordsurj} (if $F^p=C_p$, the canonical infinite-jet map is surjective).
Theorem~\ref{l20:thm:virli} below is the case $g=1$, $a^1=\omega$ of that machinery (L11 excludes the Virasoro case from its
Lie-current theorem and refers here for it).

\subsection{Li's filtration of the Virasoro VOA is the word filtration}

\begin{definition}[Depth of a creation word; the word filtration]\label{l20:def:virword}
The \emph{Li depth} of a creation word $L_{-j_1-2}\cdots L_{-j_b-2}\ltwentyvac$ is $j_1+\dots+j_b$. For $p\in\N$ let
$C_p\subseteq V(c,0)$ be the span of all creation words of depth $\ge p$ (in any order of the letters).
\end{definition}
\noindent File \leanpath{VOA/\allowbreak Filtered/\allowbreak VirasoroWordLi.lean}.
\begin{leancode}
abbrev virWordLevel (w : List ℕ) : ℕ := w.sum

def virWordPiece (p : ℕ) : Submodule ℂ (VirUniv c) :=
  Submodule.span ℂ {x | ∃ w : List ℕ, p ≤ virWordLevel w ∧ negWordVir c w = x}
\end{leancode}

\begin{dictionary}
\lean{VirUniv c} is the underlying space of $V(c,0)$. A word is a list \lean{w : List ℕ}; its depth is the sum of its
entries (\lean{w.sum}). Since $L_{-j-2}=\omega_{(-j-1)}$ and $\omega$ has weight $2$, a word of weight $N$ and length $b$
has depth $N-2b$: depth counts the derivatives, length counts the copies of $T(z)$. \lean{virWordPiece c p} is $C_p$.
Words of any order are allowed (no sorting): $C_p$ is a span, so this costs nothing, and it is what the proof needs.
The family is decreasing (\lean{virWordPiece_antitone}) and $C_0=V(c,0)$ (\lean{virWordPiece_zero_eq_top}, because words
span).
\end{dictionary}

\begin{physics}
$L_{-j-2}\ltwentyvac$ is the state of $\partial^jT/j!$. A word of depth $p$ is a normal-ordered product of stress
tensors carrying $p$ derivatives in total. Li's filtration is ``the number of derivatives'', and this definition
makes the slogan literal for the Virasoro VOA.
\end{physics}

\begin{status}
Definition; no hypotheses.
\end{status}

\begin{theorem}[Li's filtration of $V(c,0)$ and of $L(c,0)$ is the word filtration]\label{l20:thm:virli}
For every $c\in\C$ and every $p\in\N$: $F^pV(c,0)=C_p$, and $F^pL(c,0)$ is the image of $C_p$ under the quotient map
$V(c,0)\to L(c,0)$.
\end{theorem}
\noindent File \leanpath{VOA/\allowbreak Filtered/\allowbreak VirasoroWordLi.lean}.
\begin{leancode}
theorem liFiltration_eq_virWordPiece (p : ℕ) :
    liFiltration (virasoroVOA c) p = virWordPiece c p :=

theorem simpleLiFiltration_eq_virWordPiece (p : ℕ) :
    liFiltration (VirSimple.virasoroSimple c) p =
      (virWordPiece c p).map (VirSimple.maxIdealVir c).mkQ := by

theorem virWordNegativeModeClosed (q n : ℕ) (u : VirUniv c) {v : VirUniv c}
    (hv : v ∈ virWordPiece c q) :
    (virasoroVOA c).smul u (-(n : ℤ) - 1) v ∈ virWordPiece c (q + n) := by
\end{leancode}

\begin{dictionary}
Namespace \lean{VOA.VirasoroLi}; \lean{c : ℂ} is a section variable, so both statements hold for \emph{every} central
charge. \lean{liFiltration W p} is Li's $F^pW$: the span of $u^1_{(-n_1-1)}\cdots u^r_{(-n_r-1)}\ltwentyvac$ with
$\sum n_i\ge p$, the $u^i$ arbitrary states (Chapter~\ref{ch:L11}, Definition~\ref{l11:def:gfil}). \lean{.map mkQ} is
the image under the quotient map. \lean{W.smul u n v} is $u_{(n)}v$, so \lean{virWordNegativeModeClosed} says
$u_{(-n-1)}C_q\subseteq C_{q+n}$ for \emph{every} state $u$: this is the key estimate.

\textbf{Same, stronger or weaker?} The same as the textbook statement. Note the asymmetry: Li's definition allows
arbitrary states $u^i$, while $C_p$ uses only $\omega$; the theorem says that nothing is lost.
\end{dictionary}

\textbf{Proof on paper.} $C_p\subseteq F^pV$: the word $L_{-j_1-2}\cdots L_{-j_b-2}\ltwentyvac$ is
$\omega_{(-j_1-1)}\cdots\omega_{(-j_b-1)}\ltwentyvac$, a generator of $F^{\sum j_i}$. For the converse one shows by
induction that negative modes of arbitrary states raise depth correctly. Put $t(q,n)=\max(q-n-1,0)$ for $q\in\N$,
$n\in\Z$; for $n=-s-1<0$ this is $q+s$.
\begin{enumerate}
\item \emph{One mode of $\omega$.} $\omega_{(n)}C_q\subseteq C_{t(q,n)}$ for all $n\in\Z$. For $n<0$, $\omega_{(n)}=L_{n-1}$
  prepends a letter of depth $-n-1$. For $n\ge0$ commute $L_{n-1}$ through the word with $[L_a,L_b]=(a-b)L_{a+b}+
  \frac{c}{12}(a^3-a)\delta_{a+b,0}$: each commutator replaces two letters by one, or by a scalar, and an induction on the
  length of the word checks the depth bound term by term (\lean{Lop_negWordVir_mem_virWordPiece}).
\item \emph{Modes of an arbitrary word.} For $u=\omega_{(-r-1)}u'$ the Borcherds iterate formula writes
  $u_{(N)}=\sum_{i\ge0}(-1)^i\binom{-r-1}{i}\bigl(\omega_{(-r-1-i)}u'_{(N+i)}-(-1)^{r+1}u'_{(-r-1+N-i)}\omega_{(i)}\bigr)$; by
  induction on the length of $u$ and step~1, each term maps $C_q$ into $C_{t(q,N)}$ (one checks the inequality
  $t(t(q,N+i),-r-1-i)\ge t(q,N)$ and its analogue for the second term, \lean{negWordVir_mode_mem_virWordPiece}).
\item Since words span $V(c,0)$, step~2 holds for every $u$; with $N=-n-1$ it is
  $u_{(-n-1)}C_q\subseteq C_{q+n}$. Applying it letter by letter to a generator of $F^p$ shows $F^p\subseteq C_p$.
\end{enumerate}
The statement for $L(c,0)$ follows because Li's filtration of a quotient by a graded ideal is the image of Li's
filtration (Chapter~\ref{ch:L11}). Only the spanning of words is used, never their linear independence. \qed

\begin{physics}
In a theory generated by one field $T$, every local operator is a normal-ordered polynomial in $T$ and its
derivatives, and Li's ``derivative counting'' of an arbitrary composite operator can be read off from such an
expression without loss. Physically obvious; the content is that OPE rearrangements (normal-ordering corrections) never
\emph{decrease} the derivative count, which is what the tracked estimate verifies.
\end{physics}

\begin{status}
\textbf{Proved, unconditional} (isolated-module check; aggregate build not re-run), every $c$; \texttt{\#print axioms} of both: \lean{[propext, Classical.choice, Quot.sound]}.
\end{status}

\subsection{The canonical map and what ``classically free'' means here}

\begin{definition}[The generator, the relation ideal, the canonical map, classical freeness]\label{l20:def:vircanon}
Let $\bar\omega\in L(c,0)$ be the class of the conformal vector and $x=[\bar\omega]\in F^0L_2/F^1L_2\subseteq\ltwentygr L(c,0)$
its Li class of bidegree $(0,2)$. In $\C[T]$ (one variable) let $I_m=(T^{m+1})$. If $I_m$ is killed by $T\mapsto x$, the
\emph{canonical map} is the unique homomorphism of differential algebras
\begin{equation*}
\pi\colon\ltwentyJ\bigl(\C[T]/(T^{m+1})\bigr)\longrightarrow\ltwentygr L(c,0),\qquad T_{(j)}\longmapsto\frac{\partial^jx}{j!}
=\bigl[L_{-j-2}\ltwentyvac\bigr]\in F^jL_{j+2}/F^{j+1}L_{j+2},
\end{equation*}
and $L(c,0)$ is \emph{classically free} (for this presentation) if $\pi$ is bijective. The central charge of $M(2,2m+3)$ is
$c_m$ as above.
\end{definition}
\noindent File \leanpath{VOA/\allowbreak Filtered/\allowbreak InfiniteClassicalFreeness.lean}.
\begin{leancode}
noncomputable def liInfinitePresentedJetMap (W : VertexOperatorAlgebra V)
    (x : σ → liAssociatedGraded W) (I : Ideal (MvPolynomial σ ℂ))
    (hI : I ≤ RingHom.ker (MvPolynomial.aeval x).toRingHom) :
    InfiniteJet.PresentedJet (R := ℂ) I →ₐ[ℂ] liAssociatedGraded W :=
  InfiniteJet.presentedJetToDifferentialAlgebra (liAssociatedGradedTranslation W)
    (liAssociatedGradedTranslation_leibniz W) (liAssociatedGradedTranslation_one W) x I hI

def LiPresentationClassicallyFree (W : VertexOperatorAlgebra V)
    (x : σ → liAssociatedGraded W) (I : Ideal (MvPolynomial σ ℂ))
    (hI : I ≤ RingHom.ker (MvPolynomial.aeval x).toRingHom) : Prop :=
  Function.Bijective (liInfinitePresentedJetMap W x I hI)
\end{leancode}
\noindent File \leanpath{VOA/\allowbreak Examples/\allowbreak VirasoroMinimalClassicalFreeness.lean}.
\begin{leancode}
def omegaBar : VirUniv c ⧸ maxIdealVir c := (maxIdealVir c).mkQ (omegaVir c)

abbrev gen : Fin 1 → VirUniv c ⧸ maxIdealVir c := fun _ => omegaBar c

abbrev wt : Fin 1 → ℕ := fun _ => 2

abbrev x : Fin 1 → liAssociatedGraded (virasoroSimple c) :=
  fun i => liDegreeZeroClass (virasoroSimple c) (wt i) ⟨gen c i, gen_mem c i⟩

abbrev relIdeal (m : ℕ) : Ideal (MvPolynomial (Fin 1) ℂ) :=
  Ideal.span {MvPolynomial.X 0 ^ (m + 1)}

abbrev cm (m : ℕ) : ℂ := (TwoSeries.twoSeriesLevels m).c

theorem cm_eq (m : ℕ) :
    cm m = 1 - 3 * (2 * (m : ℂ) + 1) ^ 2 / (2 * (m : ℂ) + 3) := by

theorem relIdeal_le (m : ℕ) :
    relIdeal m ≤ RingHom.ker (MvPolynomial.aeval (x (cm m))).toRingHom :=
\end{leancode}

\begin{dictionary}
\textbf{Decoding.} The generator set is \lean{Fin 1}, a one-element type: one variable $T$ (\lean{MvPolynomial.X 0}).
\lean{omegaVir c} is $\omega=L_{-2}\ltwentyvac$, \lean{omegaBar c} its class in $L(c,0)$, \lean{wt} its weight $2$, and
\lean{liDegreeZeroClass W N v} is the class of $v\in W_N$ in $F^0W_N/F^1W_N$, inserted into $\ltwentygr W$
(Chapter~\ref{ch:L11}). \lean{InfiniteJet.PresentedJet I} is the presented arc algebra $\ltwentyJ(\C[T]/I)$ of
Definition~\ref{l20:def:jet}, with Hasse--Schmidt jets $T_{(j)}\leftrightarrow\partial^jT/j!$.
\lean{presentedJetToDifferentialAlgebra} is the universal property of $\ltwentyJ$: a map from $\C[T]/I$ into a differential
algebra extends uniquely to a differential-algebra map from the arc algebra. Here the target derivation is the descended
$L_{-1}$ (\lean{liAssociatedGradedTranslation W}). The hypothesis \lean{hI} (``the relations hold in $\ltwentygr W$'') is
what makes the map exist; it is a proof argument, and the map does not depend on which proof is supplied.
\lean{LiPresentationClassicallyFree} is bijectivity of that map. Its docstring says it ``deliberately does not claim'' that the
source is identified with an independently defined arc functor; for this example that identification is
\lean{actualC2ArcEquiv} below.

\textbf{The central charge.} \lean{cm m} is defined through the library's $(2,2n+3)$-series level data
(\lean{twoSeriesLevels}), and \lean{cm_eq} proves that it equals $1-3(2m+1)^2/(2m+3)=c_{2,2m+3}$. The closed form agrees with the
minimal-model charge of Definition~\ref{l06:def:cpq} by arithmetic, $6(2m+1)^2/(2(2m+3))=3(2m+1)^2/(2m+3)$ (that no Lean
theorem states this is recorded in the warning at the end of this section). \lean{cm_zero : cm 0 = 0} and \lean{cm_one : cm 1 = -22 / 5} are proved.

\textbf{How could this be vacuous?} If \lean{relIdeal m} were too big (say $(T)$ for every $m$), the source would be
too small and bijectivity false; if the target were not the actual $L(c,0)$, the statement would be about the wrong
object. Neither happens: the ideal is literally $(T^{m+1})$, and the target is \lean{virasoroSimple (cm m)}, the quotient by
the maximal proper graded ideal.
\end{dictionary}

\begin{physics}
$x$ is the classical stress tensor $T(z)$ viewed as a coordinate on the associated scheme; $\pi$ sends the classical
field $\partial^jT/j!$ to the class of the quantum operator $\partial^jT/j!$. The relation $T^{m+1}=0$ is the classical
shadow of the vacuum null vector at level $2m+2$.
\end{physics}

\begin{status}
Definitions; \lean{cm_eq}, \lean{relIdeal_le}: \lean{[propext, Classical.choice, Quot.sound]}.
\end{status}

\subsection{The two inputs from the singular vector}

The relation $T^{m+1}$ and the injectivity both come from the vacuum singular vector of $M(2,2m+3)$, which sits at weight
$2m+2$ and has leading term $\alpha L_{-2}^{m+1}\ltwentyvac$ with $\alpha\ne0$. The library packages its two consequences as
two propositions, both proved for every $m$ in its $(2,2n+3)$-series files (\leanpath{VOA/\allowbreak Examples/\allowbreak TwoSeries*.lean});
we state exactly what is used and explain the proofs below.

\begin{definition}[Jet relations; the symbol property]\label{l20:def:virinputs}
Let $P_k^{(n)}=\sum_{e_1+\dots+e_n=k}L_{-e_1-2}\cdots L_{-e_n-2}\ltwentyvac$. Let $\mathcal F_pV(c,0)_N$ be the span of sorted
creation words of weight $N$ and length $\le p$ (Chapter~\ref{ch:L19}, Definition~\ref{l19:def:lengthflag}). Then:
\begin{enumerate}
\item $c$ has \emph{jet relations of order $n+1$} if there is $\alpha\ne0$ such that for every $k$ some element of the maximal
  ideal agrees with $\alpha P_k^{(n+1)}$ modulo $\mathcal F_nV(c,0)_{2n+2+k}$.
\item Let $\mathcal R_n=\C[t_0,t_1,\dots]/(\text{coefficients of }(\sum_jt_js^j)^n)$. The \emph{symbol property} for $(c,\mathcal R_n,t)$ is the existence
  of linear maps $\sigma_b\colon V(c,0)\to \mathcal R_n$ ($b\in\N$) that kill the maximal ideal, kill every creation word of length $<b$,
  and send every creation word $L_{-j_1-2}\cdots L_{-j_b-2}\ltwentyvac$ of length $b$ to $t_{j_1}\cdots t_{j_b}$.
\end{enumerate}
\end{definition}
\noindent File \leanpath{VOA/\allowbreak Examples/\allowbreak TwoSeriesCharacterLeading.lean}.
\begin{leancode}
def jetSum : ℕ → ℕ → VirUniv c
  | 0, k => if k = 0 then vacVec c else 0
  | m + 1, k => ∑ p ∈ Finset.antidiagonal k, Lop c (-(p.1 : ℤ) - 2) (jetSum m p.2)

def HasJetRelations (c : ℂ) (n : ℕ) : Prop :=
  ∃ α : ℂ, α ≠ 0 ∧ ∀ k : ℕ, ∃ σ ∈ maxIdealVir c,
    σ - α • jetSum c (n + 1) k ∈ virFil c n ((2 * n + 2 + k : ℕ) : ℤ)
\end{leancode}
\noindent File \leanpath{VOA/\allowbreak Examples/\allowbreak TwoSeriesJetRing.lean}.
\begin{leancode}
abbrev TPoly : Type := MvPolynomial ℕ ℂ

def tSer : PowerSeries TPoly := PowerSeries.mk X

def jetIdealT (n : ℕ) : Ideal TPoly :=
  Ideal.span (Set.range fun j => PowerSeries.coeff j (tSer ^ n))

abbrev Rn (n : ℕ) : Type := TPoly ⧸ jetIdealT n

def tn (n : ℕ) (j : ℕ) : Rn n := Ideal.Quotient.mk (jetIdealT n) (X j)
\end{leancode}
\noindent File \leanpath{VOA/\allowbreak Examples/\allowbreak TwoSeriesDSVirSymbol.lean}.
\begin{leancode}
def vBelow (b : ℕ) : Submodule ℂ (VirUniv c) :=
  Submodule.span ℂ {x | ∃ w : List ℕ, w.length < b ∧ negWordVir c w = x}

def VirSymbol : Prop :=
  ∃ σ : ℕ → (VirUniv c →ₗ[ℂ] R),
    (∀ b, ∀ v ∈ maxIdealVir c, σ b v = 0) ∧
    (∀ b, ∀ v ∈ vBelow c b, σ b v = 0) ∧
    (∀ w : List ℕ, σ w.length (negWordVir c w) = (w.map t).prod)
\end{leancode}

\begin{dictionary}
\lean{Lop c n} is $L_n$ on $V(c,0)$ and \lean{vacVec c} the vacuum (Chapter~\ref{ch:L06}); \lean{jetSum c n k} is
$P^{(n)}_k$, defined recursively over the first letter (\lean{Finset.antidiagonal k} is the set of pairs summing to $k$).
\lean{virFil c n N} is $\mathcal F_nV(c,0)_N$ (sorted words of energy $N$ and length $\le n$; Chapter~\ref{ch:L19}).
\lean{HasJetRelations c n} is item~1, with \lean{σ} the element of the maximal ideal (\lean{σ} here is a vector, not the
symbol map). \lean{TPoly} is $\C[t_0,t_1,\dots]$, \lean{tSer} the series $\sum_jt_js^j$, \lean{jetIdealT n} the ideal of its
$n$-th power's coefficients, \lean{Rn n} the quotient $\mathcal R_n$, and \lean{tn n j} the class of $t_j$; so $\mathcal R_n$ is the arc
algebra $\ltwentyJ(\C[T]/(T^n))$ written with one family of variables. \lean{vBelow c b} is the span of words of length $<b$, in any
order. \lean{VirSymbol c R t} is item~2 (\lean{c}, \lean{R}, \lean{t} are section variables); \lean{w.map t} is the list
$(t_{j_1},\dots,t_{j_b})$ and \lean{.prod} its product.

\textbf{Why item 2 is strong.} The third clause of \lean{VirSymbol} is required for words in \emph{every} order and of \emph{every} weight, with one
linear map per length. Since different orderings of the same letters differ by commutators of strictly smaller length,
this is consistent; but it also forces every linear relation among length-$b$ words in $L(c,0)$, modulo shorter words, to hold
among the corresponding monomials in $\mathcal R_{m+1}$. That is exactly injectivity at the level of leading symbols.
\end{dictionary}

\begin{physics}
Item~1 says that the null state at level $2m+2$ is, at leading order in the number of stress tensors, the classical constraint
$T(z)^{m+1}=0$, and that its descendants $L_{-1}^k$ give the derivatives of that constraint. Item~2 says that the map ``quantum
state $\mapsto$ classical monomial in $T$ and its derivatives'' is well defined on the physical (null-free) Hilbert space at leading
order: the quantum theory has no leading-order relations beyond the classical ones.
\end{physics}

\begin{status}
Definitions (Props). Both are proved for $c=c_m$ in Theorem~\ref{l20:thm:virinputs}; they are not open inputs.
\end{status}

\begin{theorem}[Both inputs hold for every $M(2,2m+3)$]\label{l20:thm:virinputs}
For every $m\in\N$: $c_m$ has jet relations of order $m+1$, and the symbol property holds for $(c_m,\mathcal R_{m+1},t)$.
\end{theorem}
\noindent File \leanpath{VOA/\allowbreak Examples/\allowbreak TwoSeriesCharacterEndpoint.lean}.
\begin{leancode}
theorem twoSeries_hasJetRelations_all (n : ℕ) :
    VOA.TwoSeriesCharacter.HasJetRelations (twoSeriesLevels n).c n := by
\end{leancode}
\noindent File \leanpath{VOA/\allowbreak Examples/\allowbreak TwoSeriesDSVirSymbol.lean}.
\begin{leancode}
theorem virSymbol_twoSeries (m : ℕ) :
    VirSymbol (twoSeriesLevels m).c (Rn (m + 1)) (tn (m + 1)) :=
\end{leancode}

\begin{dictionary}
\lean{(twoSeriesLevels m).c} is \lean{cm m}. No hypotheses. Both are theorems, not named inputs. \textbf{Same or weaker?} \lean{HasJetRelations} allows a possibly \emph{different} element of the maximal ideal for each
$k$, rather than the $L_{-1}$-derivatives of one singular vector; this is weaker as a statement about singular vectors but is exactly
what classical freeness needs. \textbf{Vacuity at $m=0$:} there the relations say $L_{-2-k}\ltwentyvac$ lies in the maximal ideal at
$c=0$, which is true (the simple quotient is $\C$). \textbf{Dependence:} the proof of \lean{virSymbol_twoSeries} uses
\lean{gordonLowerBound_twoSeries}, i.e.\ the Andrews--Gordon identity and the Kac-determinant lower bound (both proved).
\end{dictionary}

\textbf{Proof on paper.} \emph{Jet relations.} The vacuum singular vector of $V(c_m,0)$ at weight $2m+2$ is constructed
explicitly (a Benoit--Saint-Aubin type recursion) and its leading term is computed to be $\alpha L_{-2}^{m+1}\ltwentyvac$ modulo
words of length $\le m$, with $\alpha\ne0$. (For $m=1$ it is $\chi=(L_{-2}^2-\frac35L_{-4})\ltwentyvac$ of Example~\ref{l20:ex:leeyang},
$\alpha=1$.) Apply $L_{-1}^k/k!$. On a word, $L_{-1}$ acts as a derivation on the letters, $[L_{-1},L_{-e-2}]=(e+1)L_{-e-3}$, up to
commutators that shorten the word; so on leading terms $L_{-1}P^{(n)}_k=(k+1)P^{(n)}_{k+1}$, and $L_{-1}^kL_{-2}^{m+1}/k!$ has leading
term $P^{(m+1)}_k$. The image of the singular vector under $L_{-1}^k/k!$ lies in the maximal ideal, which gives item~1 with that
image as the element $\sigma$ of the maximal ideal.

\emph{Symbol maps: the mechanism in three steps.}
\begin{enumerate}
\item \emph{Straightening (spanning).} Fix the weight $N$ and the length $b$. Modulo words of length $<b$ and modulo the maximal
  ideal, every length-$b$ word is a combination of \emph{sorted} length-$b$ words (reordering letters costs commutators, which are
  shorter). A sorted word that is not Gordon-good of order $m$ contains $m+1$ letters within distance one; the jet relation
  $P^{(m+1)}_k$ with the appropriate $k$ rewrites it, modulo shorter words and the maximal ideal, as a combination of sorted words that
  are larger in a fixed monomial order. Since there are finitely many words of given weight, the rewriting stops: \emph{the Gordon-good
  sorted words span $L(c_m,0)_N$ modulo shorter words}, and the same rewriting holds verbatim in $\mathcal R_{m+1}$, where the relations are
  the coefficients of $(\sum_jt_js^j)^{m+1}$.
\item \emph{Counting (independence).} Summing over $b$, the number of Gordon-good sorted words of weight $N$ is the left side of
  \eqref{l20:eq:ag} (with $n=m$), which by the Andrews--Gordon identity equals the Rocha-Caridi number on the right; and the
  Kac determinant (Feigin--Fuchs) shows that the Rocha-Caridi number is $\le\dim L(c_m,0)_N$. Spanning gives the reverse inequality, so
  the Gordon-good words are a basis of $L(c_m,0)_N$, and their images form a basis of each length layer.
\item \emph{Definition of $\sigma_b$.} Define $\sigma_b$ on the basis: a Gordon-good word of length $b$ goes to its monomial
  $t_{j_1}\cdots t_{j_b}$, shorter basis words go to $0$; extend linearly and compose with $V(c_m,0)\to L(c_m,0)$ (so the maximal ideal
  is killed). Because step~1 rewrites a word and its monomial \emph{by the same relations}, $\sigma_b$ sends every length-$b$ word, in
  any order, to its monomial, and kills shorter words: the three clauses of the symbol property.
\end{enumerate}
\emph{The case $m=1$.} The relation is $P^{(2)}_k=\sum_{a+b=k}L_{-a-2}L_{-b-2}\ltwentyvac$, the Gordon-good words are the
Rogers--Ramanujan words (sorted, consecutive letters differing by $\ge2$), and the counts are those of the table of
Example~\ref{l20:ex:leeyang}: at weight $6$, for instance, the sorted length-two words are $L_{-4}L_{-2}$, $L_{-3}^2$ (depths $(2,0)$,
$(1,1)$); only $(2,0)$ is Gordon-good, and $L_{-3}^2\ltwentyvac$ is rewritten through $P^{(2)}_2 = 2L_{-4}L_{-2}+L_{-3}^2$ (leading terms),
matching $T_1^2\equiv-2T_0T_2$ in the table. \qed

\begin{physics}
The singular vector is the null state that the minimal model removes; its leading term $L_{-2}^{m+1}\ltwentyvac$ says
``$T^{m+1}=0$'' classically. The symbol property says that, at leading order in the number of stress tensors, this
null state and its derivatives are the \emph{only} relations: the Gordon-good words are the Feigin--Stoyanovsky
quasi-particle basis of the minimal model.
\end{physics}

\begin{status}
\textbf{Proved, unconditional} (isolated-module check; aggregate build not re-run); \texttt{\#print axioms} of both: \lean{[propext, Classical.choice, Quot.sound]}.
\end{status}

\begin{warning}[title={Recorded overclaim: \texttt{virSymbol\_twoSeries} is not by itself classical freeness}]
\lean{virSymbol_twoSeries} is a PBW-\emph{length} leading-symbol theorem into the presented ring $\mathcal R_{m+1}$. It does not
mention Li's filtration, $V/C_2(V)$ or the canonical map. Earlier project notes and docstrings described it as
``classical freeness''; that description was corrected in the project, and the file's docstring says so (the record is cited in
Section~\ref{l20:sec:verification}). The bridge from the length filtration to Li's filtration is Theorem~\ref{l20:thm:virli} together with the
comparison in the proof of Theorem~\ref{l20:thm:virfree}: weight $=$ depth $+\,2\cdot$length, so at fixed weight and fixed
Li depth the length is determined.
\end{warning}

\subsection{The headline: classical freeness and the \texorpdfstring{$C_2$}{C2}-algebra}

\begin{theorem}[Classical freeness of $M(2,2m+3)$]\label{l20:thm:virfree}
For every $m\in\N$, the canonical map $\pi\colon\ltwentyJ(\C[T]/(T^{m+1}))\to\ltwentygr L(c_m,0)$ of
Definition~\ref{l20:def:vircanon} is bijective; hence it is an isomorphism of differential algebras. In particular this holds
for $M(2,3)$ ($m=0$) and for Lee--Yang $M(2,5)$ ($m=1$).
\end{theorem}
\noindent File \leanpath{VOA/\allowbreak Examples/\allowbreak VirasoroMinimalClassicalFreeness.lean}.
\begin{leancode}
theorem classicallyFree (m : ℕ) :
    LiPresentationClassicallyFree (virasoroSimple (cm m)) (x (cm m)) (relIdeal m)
      (relIdeal_le m) :=
  ⟨jetMap_injective (cm m) m (hasJetRelations m) (TwoSeries.DSSlice.virSymbol_twoSeries m),
    jetMap_surjective (cm m) m (hasJetRelations m)⟩

def canonicalEquiv (m : ℕ) :
    InfiniteJet.PresentedJet (relIdeal m) ≃ₐ[ℂ] liAssociatedGraded (virasoroSimple (cm m)) :=
  AlgEquiv.ofBijective _ (classicallyFree m)

theorem classicallyFree_M23 :
    LiPresentationClassicallyFree (virasoroSimple (cm 0)) (x (cm 0)) (relIdeal 0)
      (relIdeal_le 0) :=
  classicallyFree 0

theorem classicallyFree_leeYang :
    LiPresentationClassicallyFree (virasoroSimple (cm 1)) (x (cm 1)) (relIdeal 1)
      (relIdeal_le 1) :=
  classicallyFree 1

theorem jetMap_injective (m : ℕ) (hJ : TwoSeriesCharacter.HasJetRelations c m)
    (hsym : TwoSeries.DSSlice.VirSymbol c (TwoSeries.StrictModel.Rn (m + 1))
      (TwoSeries.StrictModel.tn (m + 1))) :
    Function.Injective
      (liInfinitePresentedJetMap (virasoroSimple c) (x c) (relIdeal m) (relIdeal_le_li c m hJ)) := by
\end{leancode}

\begin{dictionary}
Namespace \lean{VOA.VirasoroClassicalFreeness}. \lean{classicallyFree m} unfolds (Definition~\ref{l20:def:vircanon}) to
``\lean{liInfinitePresentedJetMap (virasoroSimple (cm m)) (x (cm m)) (relIdeal m) _} is bijective'': the map $\pi$ from
$\ltwentyJ(\C[T]/(T^{m+1}))$ to $\ltwentygr L(c_m,0)$. The one-line proof term shown is informative: injectivity is
\lean{jetMap_injective} fed with the two inputs of Theorem~\ref{l20:thm:virinputs}, and surjectivity is
\lean{jetMap_surjective}, fed with the jet relations only. \lean{jetMap_injective} is stated for an arbitrary
\lean{c}; it is the general implication ``jet relations $+$ symbol property $\Rightarrow$ injective''.
\lean{canonicalEquiv m} packages $\pi$ as an algebra isomorphism (\lean{AlgEquiv.ofBijective}); its definition, not just its
type, says it is $\pi$. \lean{classicallyFree_M23}, \lean{classicallyFree_leeYang} are the instances $m=0,1$.

\textbf{Hypotheses.} None besides \lean{m : ℕ}. In particular $m=0$ is included.

\textbf{Same, stronger or weaker?} The same as \eqref{l20:eq:virfree} for the presented arc algebra; the identification of
the presented arc algebra with the arc algebra of the actual $C_2$-algebra is Theorem~\ref{l20:thm:virc2}. The canonical map is
a map of differential algebras by construction (\lean{jetMap_translation}); compatibility with the $\lambda$-brackets is
Theorem~\ref{l20:thm:virpva}. Unlike the affine \lean{arcPoissonEquiv_conformal_iff}, there is no separate theorem that
$\pi$ preserves conformal weight; it does, because $T_{(j)}$ has weight $j+2$ and $[L_{-j-2}\ltwentyvac]$ lies in weight $j+2$
by the definition of $\pi$, and the injectivity proof works one bidegree at a time.
\end{dictionary}

\textbf{Proof on paper.} Grade $\ltwentyJ\C[T]$ by bidegree: $T_{(j)}$ has Li depth $j$ and weight $j+2$, so a monomial with
$b$ letters and depth $p$ has bidegree $(p,p+2b)$; $\pi$ preserves bidegree, where $\ltwentygr L$ is bigraded by
$(p,N)\mapsto F^pL_N/F^{p+1}L_N$.

\emph{The relations hold.} In $R_L=F^0L/F^1L$ the class $[\omega]^{m+1}$ is the class of $L_{-2}^{m+1}\ltwentyvac$
(\lean{mk₂_replicate}). By the jet relation with $k=0$, $\alpha L_{-2}^{m+1}\ltwentyvac$ is congruent, modulo the maximal
ideal, to a combination of sorted words of weight $2m+2$ and length $\le m$; such words have depth $(2m+2)-2\cdot\text{length}\ge2$,
so they lie in $C_1$, whose image is $F^1L=C_2(L)$ by Theorem~\ref{l20:thm:virli} and Li's $F^1=C_2$. Hence $[\omega]^{m+1}=0$ in
$R_L$ (\lean{mk₂_omegaBar_pow_eq_zero}), so $T^{m+1}\mapsto0$ in $\ltwentygr_0L$, and $\pi$ is defined.

\emph{Surjective.} This is Theorem~\ref{l11:thm:genwordsurj}(4) for the single generator $\omega$, whose hypothesis is
Theorem~\ref{l20:thm:virli}; directly: by Theorem~\ref{l20:thm:virli}, $F^pL_N$ is spanned by images of words of depth $\ge p$; so the layer
$F^pL_N/F^{p+1}L_N$ is spanned by classes of words of depth exactly $p$. The class of
$L_{-j_1-2}\cdots L_{-j_b-2}\ltwentyvac$ in that layer is the product $\prod_i\partial^{j_i}x/j_i!=\pi(T_{(j_1)}\cdots T_{(j_b)})$
(the generic statement ``a word's Li class is the product of the divided derivatives of its letters'' is in
Chapter~\ref{ch:L11}). Hence $\pi$ is onto every layer.

\emph{Injective.} Let $h$ be a jet polynomial with $\pi(h)=0$. Its bihomogeneous components are also killed (the target is
a direct sum over bidegrees), so assume $h$ has bidegree $(p,N)$. If $N-p$ is odd, $h=0$ (no monomials). Otherwise $N=p+2b$ and every
monomial of $h$ has $b$ letters; let $v\in V(c,0)_N$ be the same combination of the corresponding words. Then $\pi(h)=0$ says
$[v]\in F^{p+1}L$, i.e.\ $v=y+z$ with $z$ in the maximal ideal and $y\in C_{p+1}$ (Theorem~\ref{l20:thm:virli}). Project to weight
$N$ (the maximal ideal is graded): the weight-$N$ part of $y$ is a combination of words of weight $N$ and depth $\ge p+1$, hence of
length $\le b-1$ (weight $=$ depth $+\,2\cdot$length). Apply the symbol map $\sigma_b$: it kills $z$ and the shorter words, so
$\sigma_b(v)=0$. But $\sigma_b(v)$ is $h$ read in $\mathcal R_{m+1}$, with $T_{(j)}\mapsto t_j$. So the image of $h$ in
$\mathcal R_{m+1}=\ltwentyJ(\C[T]/(T^{m+1}))$ is zero; that is, $h$ lies in the jet ideal (the comparison map
$\ltwentyJ\C[T]\to \mathcal R_{m+1}$, \lean{psi}, has kernel inside the jet ideal of $(T^{m+1})$, \lean{mem_jetIdeal_of_psi_eq_zero}). \qed

\begin{physics}
In the $(A_1,A_{2m})$ Argyres--Douglas theories, whose chiral algebras are exactly $M(2,2m+3)$ (Beem--Lemos--Liendo--Peelaers--Rastelli--van
Rees; Cordova--Shao), the Higgs branch is a point but the Higgs branch \emph{scheme} $\operatorname{Spec}\C[T]/(T^{m+1})$ is a fat point.
In the heuristic dictionary between the two sides, classical freeness says that the Schur index (the vacuum character) is the
Hilbert series of the space of classical stress-tensor configurations $T(z)$ with $T(z)^{m+1}=0$. On characters, this is the sum side of the
Andrews--Gordon identity: the classical counting of such configurations is the Feigin--Stoyanovsky fermionic formula.
\end{physics}

\begin{status}
\textbf{Proved, unconditional} (isolated-module check; aggregate build not re-run) (no hypothesis; every $m\in\N$). \texttt{\#print axioms} of \lean{classicallyFree},
\lean{classicallyFree_M23}, \lean{classicallyFree_leeYang}, \lean{canonicalEquiv}, \lean{jetMap_injective}:
\lean{[propext, Classical.choice, Quot.sound]} (run 2026-10-27, Section~\ref{l20:sec:verification}).
\end{status}

\begin{theorem}[The full non-reduced $C_2$-algebra of $M(2,2m+3)$; arcs]\label{l20:thm:virc2}
For every $m\in\N$: (1) $T\mapsto[\bar\omega]$ induces an isomorphism $\C[T]/(T^{m+1})\cong R_{L(c_m,0)}$; in particular
$[\omega]^{m+1}=0$ in $R_{L(c_m,0)}$; (2) for every commutative $\C$-algebra $A$ there is a bijection between algebra maps
$\ltwentyJ(\C[T]/(T^{m+1}))\to A$ and algebra maps $R_{L(c_m,0)}\to A[[t]]$, namely the one obtained from (1) and the arc universal
property (Theorem~\ref{l20:thm:arc}). In words, (2) says that the presented arc algebra $\ltwentyJ(\C[T]/(T^{m+1}))$
\emph{is} the arc algebra of the actual $C_2$-algebra $R_{L(c_m,0)}$, so Theorem~\ref{l20:thm:virfree} identifies $\ltwentygr L(c_m,0)$
with the arc algebra of $R_{L(c_m,0)}$ itself.
\end{theorem}
\noindent File \leanpath{VOA/\allowbreak Examples/\allowbreak VirasoroMinimalClassicalFreeness.lean}.
\begin{leancode}
theorem mk₂_omegaBar_pow_eq_zero (m : ℕ) (hJ : TwoSeriesCharacter.HasJetRelations c m) :
    mk₂ (virasoroSimple c) (omegaBar c) ^ (m + 1) = 0 := by

def baseC2AlgEquiv (m : ℕ) :
    (MvPolynomial (Fin 1) ℂ ⧸ relIdeal m) ≃ₐ[ℂ] A₂ (virasoroSimple (cm m)) :=
  baseC2AlgEquiv_of_arc_injective (virasoroSimple (cm m)) wt (genGr (cm m)) (relIdeal m)
    (relIdeal_le_c2 (cm m) m (hasJetRelations m)) (relIdeal_le m) (classicallyFree m).1
    (aeval_c2_surjective (cm m))

def actualC2ArcEquiv (m : ℕ) {A : Type*} [CommRing A] [Algebra ℂ A] :
    (InfiniteJet.PresentedJet (relIdeal m) →ₐ[ℂ] A) ≃
      (A₂ (virasoroSimple (cm m)) →ₐ[ℂ] PowerSeries A) :=
  InfiniteJet.arcEquivOfBaseAlgebraEquiv (relIdeal m) (baseC2AlgEquiv m)

theorem actualC2ArcEquiv_natural (m : ℕ) {A B : Type*} [CommRing A] [Algebra ℂ A]
    [CommRing B] [Algebra ℂ B] (f : InfiniteJet.PresentedJet (relIdeal m) →ₐ[ℂ] A)
    (g : A →ₐ[ℂ] B) :
    actualC2ArcEquiv m (g.comp f) = (PowerSeries.mapAlgHom g).comp (actualC2ArcEquiv m f) :=
\end{leancode}

\begin{dictionary}
\lean{A₂ W} is $R_W=W/C_2(W)$ and \lean{mk₂ W} the class map (Chapter~\ref{ch:L12}). \lean{baseC2AlgEquiv m} is, by its
definition, \lean{baseC2AlgEquiv_of_arc_injective} applied to the generator $[\bar\omega]$: this is the same generic
construction used for the affine \lean{baseC2AlgEquiv} (Theorem~\ref{l20:thm:m8}), i.e.\ \lean{AlgEquiv.ofBijective} of the
canonical map $\C[T]/(T^{m+1})\to R_L$, $T\mapsto[\bar\omega]$. Its inputs are visible in the printed body: the relation
holds in $R_L$ (\lean{relIdeal_le_c2}, from \lean{mk₂_omegaBar_pow_eq_zero}), the arc map is injective
(\lean{(classicallyFree m).1}), and $R_L$ is generated by $[\bar\omega]$ (\lean{aeval_c2_surjective}).
\lean{actualC2ArcEquiv} is \lean{arcEquivOfBaseAlgebraEquiv} (Theorem~\ref{l20:thm:arc}) composed with it.

\textbf{What the types say.} Exactly as for the affine headlines (warning after Theorem~\ref{l20:thm:c2arc}): the type of
\lean{baseC2AlgEquiv} says only ``abstractly isomorphic''; that the isomorphism sends $T$ to $[\omega]$ is in its definition. The type
of \lean{actualC2ArcEquiv} is a bijection of sets for each $A$; its naturality in $A$ is the separate theorem
\lean{VirasoroClassicalFreeness.actualC2ArcEquiv_natural}, printed above, so the presented jet algebra represents the arc
functor of $R_L$. The definitions are the canonical ones.
\end{dictionary}

\textbf{Proof on paper.} $R_L$ is generated by $[\omega]$: words of positive depth lie in $F^1=C_2$, and a word of depth $0$ is
$L_{-2}^b\ltwentyvac$, whose class is $[\omega]^b$. The relation holds (proof of Theorem~\ref{l20:thm:virfree}). Injectivity of
$\C[T]/(T^{m+1})\to R_L$: it is the restriction of the injective $\pi$ to constant arcs (Li degree $0$), and constant arcs embed
in $\ltwentyJ$. Then (2) is Theorem~\ref{l20:thm:arc}. \qed

\begin{physics}
$R_{M(2,2m+3)}=\C[T]/(T^{m+1})$ is the chiral ring of the $(A_1,A_{2m})$ theory restricted to the Higgs-branch sector: the
powers $T,T^2,\dots,T^m$ of the stress tensor survive as nilpotent elements, although the associated variety is a single point
(the VOA is lisse). The physics term is ``non-reduced Higgs branch''.
\end{physics}

\begin{status}
\textbf{Proved, unconditional} (isolated-module check; aggregate build not re-run); \lean{mk₂_omegaBar_pow_eq_zero} carries the hypothesis \lean{hJ}, discharged by
\lean{hasJetRelations m} in the headlines. \texttt{\#print axioms} of all three:
\lean{[propext, Classical.choice, Quot.sound]}.
\end{status}

\subsection{The Poisson vertex structure: the bracket is forced to vanish}

\begin{theorem}[Classical freeness as Poisson vertex algebras]\label{l20:thm:virpva}
Let $\ltwentyJ(\C[T]/(T^{m+1}))$ carry the jet derivation and the \emph{zero} $\lambda$-bracket. For every $m\in\N$ the canonical map
$\pi$ is an isomorphism of Poisson vertex algebras onto $\ltwentygr L(c_m,0)$ with its Li Poisson vertex structure
(Definition~\ref{l20:def:lipva}). Consequently all $(n)$-products ($n\ge0$) of $\ltwentygr L(c_m,0)$ vanish.
\end{theorem}
\noindent File \leanpath{VOA/\allowbreak Examples/\allowbreak VirasoroMinimalPoisson.lean}.
\begin{leancode}
noncomputable def zeroPVA (A : Type*) [CommRing A] [Algebra ℂ A] (T : Derivation ℂ A A) :
    PoissonVertexAlgebra A where
  translation := T
  mode := 0

abbrev arcPVA (m : ℕ) : PoissonVertexAlgebra (InfiniteJet.PresentedJet (relIdeal m)) :=
  zeroPVA _ (InfiniteJet.translation (relIdeal m))

theorem mode_x_x (n : ℕ) :
    (liPoissonVertexAlgebra (virasoroSimple c)).mode n (x c 0) (x c 0) = 0 := by

def arcPoissonEquiv (m : ℕ) :
    PoissonVertexAlgebra.Equiv (arcPVA m) (liPoissonVertexAlgebra (virasoroSimple (cm m))) :=
  PoissonVertexAlgebra.Equiv.ofBijective (arcLiPoissonHom (cm m) m (hasJetRelations m))
    (classicallyFree m)

theorem liPoissonVertexAlgebra_mode_eq_zero (m n : ℕ)
    (a b : liAssociatedGraded (virasoroSimple (cm m))) :
    (liPoissonVertexAlgebra (virasoroSimple (cm m))).mode n a b = 0 := by
\end{leancode}

\begin{dictionary}
\lean{zeroPVA A T} is the Poisson vertex algebra (Definition~\ref{l20:def:pva}) with translation \lean{T} and every mode equal to
\lean{0}; the remaining fields (locality, sesquilinearity, Leibniz, skew-symmetry, Jacobi; not shown) are then trivial and are
discharged by \lean{rfl} or \lean{simp}. \lean{VirasoroClassicalFreeness.arcPVA m} (not to be confused with the affine
\lean{arcPVA m} of Definition~\ref{l20:def:arcpva}) is \lean{zeroPVA} on the presented arc algebra with the jet derivation
$\partial T_{(j)}=(j+1)T_{(j+1)}$. \lean{mode_x_x}: every mode of the Li class of $\omega$ with itself is zero, for every $c$.
\lean{arcPoissonEquiv m} is, by its definition, the canonical map $\pi$ (\lean{arcLiPoissonHom}, a PVA homomorphism) made into an
equivalence by \lean{classicallyFree}. \lean{liPoissonVertexAlgebra_mode_eq_zero}: all modes of $\ltwentygr L(c_m,0)$ vanish, for all
$a,b$.

\textbf{Is the zero bracket the right one?} The textbook arc PVA of a Poisson algebra $R$ has $\{x_\lambda y\}=\{x,y\}$ for
$x,y\in R$. Here $R=R_L=\C[T]/(T^{m+1})$ with Zhu's bracket $\{[a],[b]\}=[a_{(0)}b]$, and $\{[\omega],[\omega]\}=[L_{-3}\ltwentyvac]=0$
because $L_{-3}\ltwentyvac=\omega_{(0)}\omega\in C_2$. So the Poisson bracket of $R_L$ is zero and the textbook arc PVA \emph{is}
\lean{zeroPVA}. The theorem is therefore not a choice of convention: \lean{arcPoissonEquiv} is an isomorphism with an independently
defined target (\lean{liPoissonVertexAlgebra}), and its existence forces the Li structure to be zero, which
\lean{liPoissonVertexAlgebra_mode_eq_zero} states outright. There is no analogue of the affine \lean{arcPVA_unique} because the source
structure is written down explicitly.
\end{dictionary}

\textbf{Proof on paper.} On generators: $x_{(0)}x$ is the degree-$(0+0-0)$ class of $\omega_{(0)}\omega=L_{-1}\omega=L_{-3}\ltwentyvac$,
which lies in $F^1$, so it is $0$ in $F^0/F^1$; for $n\ge1$, $x_{(n)}x$ would have Li degree $0+0-n<0$, so it is $0$ by definition
(Definition~\ref{l20:def:lipva}). A homomorphism of differential algebras from an arc algebra into a PVA is a PVA homomorphism as
soon as it matches the brackets of the generators (the ``generator test'' \lean{poissonVertexHomOfBaseGenerators}, the same tool as in
the affine case); both sides give zero. Bijectivity is Theorem~\ref{l20:thm:virfree}; transporting the zero modes of the source
along the isomorphism gives the vanishing of all modes of the target. There is also a one-line parity argument: in
$\ltwentygr L\cong\ltwentyJ(\C[T]/(T^{m+1}))$ the bidegrees $(p,N)$ that occur all have $N-p$ even, while the $n$-th mode maps
bidegrees $(p,N),(q,M)$ to $(p+q-n,N+M-n-1)$, whose difference $(N-p)+(M-q)-1$ is odd. \qed

\begin{physics}
Recall the grading rule: the $n$-th Li product maps bidegrees $(p,N)$ and $(q,M)$ (Li depth, weight) to $(p+q-n,\,N+M-n-1)$.
The $TT$ OPE, $T(z)T(w)\sim\frac{c/2}{(z-w)^4}+\frac{2T(w)}{(z-w)^2}+\frac{\partial T(w)}{z-w}$, has no classical shadow in Li's
grading: the central term and the $2T$ term lower the derivative count below zero, and the $\partial T$ term has one derivative more
than the leading order. So the classical limit of a minimal model in this sense is a purely commutative, purely kinematic object:
the fat point and its arcs. This differs from the affine case, where the current algebra bracket $\{e_\lambda f\}=h$ survives.
\end{physics}

\begin{status}
\textbf{Proved, unconditional} (isolated-module check; aggregate build not re-run); \texttt{\#print axioms} of \lean{zeroPVA}, \lean{mode_x_x}, \lean{arcPoissonEquiv},
\lean{liPoissonVertexAlgebra_mode_eq_zero}: \lean{[propext, Classical.choice, Quot.sound]}.
\end{status}

\subsection{A worked example: Lee--Yang at low weight}

\begin{example}[{$m=1$: $\ltwentyJ(\C[T]/(T^2))$ against $L(-\frac{22}{5},0)$}]\label{l20:ex:leeyang}
Write $T_j$ for $T_{(j)}$ (Li depth $j$, weight $j+2$) and $Q_k=\sum_{a+b=k}T_aT_b$ for the jet relations of $T^2$:
$Q_0=T_0^2$, $Q_1=2T_0T_1$, $Q_2=2T_0T_2+T_1^2$, $Q_3=2T_0T_3+2T_1T_2$, $Q_4=2T_0T_4+2T_1T_3+T_2^2$. The relation ideal in
weight $N$ is spanned by the products $\mu Q_k$ with $\mu$ a monomial of weight $N-k-4$. A monomial with $b$ letters and weight $N$
has depth $p=N-2b$. Counting, at each weight $N$, monomials minus independent relations, sorted by $b$:
\begin{center}
\begin{tabular}{rlll}
\toprule
$N$ & $b=1$ (depth $N-2$) & $b=2$ (depth $N-4$) & total\\
\midrule
0 & \multicolumn{2}{l}{$1$ (the empty monomial)} & 1\\
1 & -- & -- & 0\\
2 & $T_0$ & -- & 1\\
3 & $T_1$ & -- & 1\\
4 & $T_2$ & $T_0^2=Q_0=0$ & 1\\
5 & $T_3$ & $T_0T_1=\frac12Q_1=0$ & 1\\
6 & $T_4$ & $T_1^2\equiv-2T_0T_2$ (one class) & 2\\
7 & $T_5$ & $T_1T_2\equiv-T_0T_3$ (one class) & 2\\
8 & $T_6$ & $T_0T_4,T_1T_3,T_2^2$ modulo $Q_4$ (two classes) & 3\\
\bottomrule
\end{tabular}
\end{center}
Monomials with $b\ge3$ vanish for $N\le8$: at $N=6$, $T_0^3=T_0Q_0$; at $N=7$, $T_0^2T_1=T_1Q_0$; at $N=8$, $T_0^2T_2=T_2Q_0$,
$T_0T_1^2=\frac12T_1Q_1$ and $T_0^4=T_0^2Q_0$. The totals $1,0,1,1,1,1,2,2,3$ for $N=0,\dots,8$ are the coefficients of the
Rogers--Ramanujan product $\prod_{n\ge0}\bigl((1-q^{5n+2})(1-q^{5n+3})\bigr)^{-1}$, which is the vacuum character of the Lee--Yang
model; so they are the dimensions of $L(-\frac{22}5,0)_N$, as Theorem~\ref{l20:thm:virfree} predicts.

On the VOA side the relation $T_0^2=0$ is the vacuum null vector at weight $4$: for $c=-\frac{22}{5}$,
\begin{equation*}
\chi=\bigl(L_{-2}^2-\tfrac35L_{-4}\bigr)\ltwentyvac
\end{equation*}
is annihilated by $L_1$ and $L_2$ (indeed $L_1\chi=(3-3)L_{-3}\ltwentyvac$ and $L_2\chi=(8+c-\frac{18}{5})L_{-2}\ltwentyvac$, and
$8+c=\frac{18}5$ exactly when $c=-\frac{22}5$). In $L(c,0)$ this gives $L_{-2}^2\ltwentyvac=\frac35L_{-4}\ltwentyvac\in F^2$, so
$[\omega]^2=0$ in $R_L=F^0/F^1$: the depth-$0$ class of $L_{-2}^2\ltwentyvac$ dies, which is $T_0^2=0$. Its leading term is
$\alpha L_{-2}^2\ltwentyvac$ with $\alpha=1$, the $m=1$ case of Definition~\ref{l20:def:virinputs}. The class $T_2\mapsto[L_{-4}\ltwentyvac]$
in depth $2$ survives, which is the weight-$4$ state of the table. At weight $6$ the two surviving classes are $[L_{-6}\ltwentyvac]$ (depth $4$)
and $[L_{-3}^2\ltwentyvac]$ (depth $2$); the universal module has the four states $L_{-6},L_{-4}L_{-2},L_{-3}^2,L_{-2}^3$ and the
null vectors $L_{-2}\chi$ and $L_{-1}^2\chi$ remove two.
\end{example}

\subsection{Comparison with the affine headlines, and a recorded dead end}

\begin{warning}[title={Points a referee should know about the Virasoro companion}]
\begin{enumerate}
\item \textbf{Same type-level caveats as the affine case.} \lean{baseC2AlgEquiv} and \lean{actualC2ArcEquiv} are canonical by
  definition, not by type; naturality of \lean{actualC2ArcEquiv} in $A$ is proved separately
  (\lean{VirasoroClassicalFreeness.actualC2ArcEquiv_natural}).
  The source of every statement is the presented arc algebra \lean{PresentedJet (relIdeal m)}.
\item \textbf{The central charge is reached through the two-series level data.} \lean{cm m} is \lean{(twoSeriesLevels m).c}, and
  only the closed form \lean{cm_eq} connects it to $c_{2,2m+3}$; there is no theorem
  \lean{cm m = minimalCentralCharge 2 (2 * m + 3)} (Definition~\ref{l06:def:cpq}). Harmless (the closed forms agree), but a
  reader who searches for the minimal-model constant will not find it used.
\item \textbf{No conformal-weight statement and no explicit character corollary.} The affine file states
  \lean{arcPoissonEquiv_conformal_iff}; the Virasoro file does not state the analogous fact, nor the resulting identity
  ``character of $L(c_m,0)$ $=$ Hilbert series of $\ltwentyJ(\C[T]/(T^{m+1}))$''. Both follow from the bidegree-preserving
  construction, but they are not Lean theorems in these files.
\item \textbf{Dead end.} An earlier generic surjectivity criterion assumed
  $V=\operatorname{span}(\text{generators}\cup\{\ltwentyvac\})+C_2(V)$ (the hypothesis \lean{hdec}); it fails for
  $M(2,2m+3)$ with $m\ge2$, because $[\omega]^2\neq0$ in $R_V=\C[T]/(T^{m+1})$. The full account is the warning at the end of
  Section~\ref{l11:sec:genwords} (Chapter~\ref{ch:L11}). The route that works is word spanning of Li's filtration
  (Theorem~\ref{l20:thm:virli}).
\end{enumerate}
\end{warning}

%======================================================================
\section{The boundary-affine \texorpdfstring{$p=3$}{p=3} particle and conductor model}\label{l20:sec:p3}
%======================================================================

\subsection{What these files are, and what they are not}

Thirty-two files of finite commutative algebra and combinatorics in \leanpath{VOA/Commutative/} remain from a route to the upper
bound at the first boundary level $p=3$ ($m=1$, $k=-4/3$) that was not completed: a ``rooted Gordon conductor'' argument in the
functional (particle-polynomial) model of Feigin--Stoyanovsky type, pursued before the argument of
Sections~\ref{l20:sec:m1}--\ref{l20:sec:m23} was formalised. They are in the library, and in this chapter, for one reason: five files of the Deligne
carrier (Chapter~\ref{ch:L24}) import five of them: \leanpath{DeligneAuxiliaryCollision} imports
\leanpath{RootedGordonConductorDivisibility}; \leanpath{DeligneAuxiliaryDegree} imports \leanpath{P3ConductorDegree};
\leanpath{QuadraticTwoPointCarrier} imports \leanpath{CollisionCommonFactors}; \leanpath{QuadraticJetCoefficientPairing} imports
\leanpath{ModeParticlePairing}; and \leanpath{QuadraticCurrentCoefficientReindexing} imports \leanpath{ModeParticlePolynomial}. The
other twenty-seven are reached only inside the cluster, or not at all.

\begin{warning}[title={The $p=3$ route is incomplete and superseded; nothing in this chapter's headlines uses it}]
None of the headline theorems of this chapter (Section~\ref{l20:sec:headlines}) depends on these files; the boundary theorem is proved
by the Gordon/partition-filtration argument of Section~\ref{l20:sec:m23}. The $p=3$ route itself was never completed in Lean. Its terminal
statement, \lean{P3ConeCarrierInterface.P3LinearResidualization.finrank_le_expectedDim}, is a fixed-sector dimension bound that takes three
uninstantiated arguments: a datum \lean{D : P3LinearResidualization W E} (two conductor factorisations, a compatibility field and
an injectivity field; its only constructor \lean{ofDivisibility} is never applied, and the coefficient space \lean{E} is an arbitrary
type connected to no vertex algebra), \lean{P3QuadraticConeCarrierIdentification} (the compatible residual window is isomorphic to the
quadratic-cone carrier of Dumanski--Feigin) and \lean{DumanskiFeiginP3GradedDimensionHypothesis} (that carrier has the expected
finite dimension, an external theorem). None is instantiated anywhere in the library, and \lean{E}, \lean{D} are never constructed
from the restricted dual of a VOA. (The file's docstring speaks of ``two remaining inputs''; the statement has three.) Several file docstrings say explicitly that
they do not construct a cone carrier, a common-factor theorem, a Hilbert series, or classical freeness
(\leanpath{P3FunctionalModel.lean}). Read the files as a library of finite lemmas, not as a proof.
\end{warning}

\textbf{The model, on paper (a summary).} This paragraph summarises the construction that the files were written for; none of its
details is needed for the three theorems presented below, which are self-contained statements about polynomials. The reader may
take it as motivation only. The underlying idea is the Feigin--Stoyanovsky \emph{functional realisation}: a state built from $N$ field
modes is paired with a polynomial in $N$ variables (the ``positions'' of the $N$ particles), so that the restricted dual of a graded
piece of a commutative differential algebra becomes a space of polynomials, and each relation of the algebra becomes a vanishing
condition when some variables collide. A \emph{species} is the field a particle comes from; variables of the same species are
interchangeable, so the polynomials are symmetric within each species; and the ``derivative degree'' is the total degree. At $p=3$ normalise the cone equation as $Q=XZ-Y^2$ (with $X,Y,Z$ the three currents) and introduce an
auxiliary field $S(z)=\sum_{r\ge1}S_rz^r$ standing for $z$ times the Casimir field. The auxiliary differential algebra is
$\mathrm{DAlg}(X,Y,Z,S)/(Q,\,XS,\,YS,\,ZS,\,S^2)$. Its restricted dual in a sector with $a,b,c,d$ particles of species $X,Y,Z,S$ and
fixed derivative degree is realised by polynomials in labelled variables $x_1..x_a,\ y_1..y_b,\ z_1..z_c,\ s_1..s_d$, symmetric within
each species. The relations become conditions on these polynomials: divisibility by every $s_\alpha$ (since $S_0$ is absent), vanishing
when two $S$ particles coincide, vanishing when an $S$ particle meets a current particle, and the coupled condition ``$F_{XZ}$ on $x=z=t$
equals $F_{YY}$ on $y_1=y_2=t$''. (Here $\mathrm{DAlg}(X,Y,Z,S)$ is the free commutative differential algebra on four fields, and $\xi_i$ runs over all current
variables $x$, $y$, $z$.) For instance, with $d=1$ and one current particle ($a+b+c=1$, say one $x$), a polynomial $F(x,s)$ must be
divisible by $s$ and vanish at $s=x$; so $s(s-x)$ divides $F$, of degree $2=dN$. In general the first three conditions force the common
factor
\begin{equation}\label{l20:eq:p3conductor}
C_{N,d}=\prod_\alpha s_\alpha\cdot\prod_{\alpha<\beta}(s_\alpha-s_\beta)^2\cdot\prod_{\alpha,i}(s_\alpha-\xi_i),\qquad
\deg C_{N,d}=d+d(d-1)+(a+b+c)d=dN,
\end{equation}
where $N=a+b+c+d$. The plan was to divide by $C_{N,d}$ and identify the residual with a
known carrier. The Lean files prove the finite algebra of this plan. We present individually the three results that Chapter~\ref{ch:L24}
reuses, and the rest collectively.

\subsection{Collision factors}

\begin{definition}[Collision factor, collision map, transposition; increasing pairs]\label{l20:def:collision}
For a field $K$ and variables $(x_i)_{i\in I}$: the \emph{collision factor} of $i,j$ is $x_i-x_j$; the \emph{collision map}
$K[x_I]\to K[x_{I\setminus\{i\}}]$ substitutes $x_i:=x_j$ ($j\ne i$); the \emph{transposition} swaps $x_i$ and $x_j$. For a linearly
ordered $I$, an \emph{increasing pair} is $(i,j)$ with $i<j$, and the Vandermonde is $\prod_{i<j}(x_i-x_j)$.
\end{definition}
\noindent File \leanpath{VOA/\allowbreak Commutative/\allowbreak CollisionCommonFactors.lean}.
\begin{leancode}
def collisionFactor (i j : I) : MvPolynomial I K :=
  X i - X j

def collisionMap (i j : I) (hji : j ≠ i) :
    MvPolynomial I K →ₐ[K] MvPolynomial {a : I // a ≠ i} K :=
  MvPolynomial.rename (collisionIndex i j hji)

def swapVariables (i j : I) : MvPolynomial I K →ₐ[K] MvPolynomial I K :=
  MvPolynomial.rename (Equiv.swap i j)

abbrev IncreasingPair (I : Type v) [LinearOrder I] :=
  {ij : I × I // ij.1 < ij.2}

def vandermonde : MvPolynomial I K :=
  ∏ ij : IncreasingPair I,
    collisionFactor (K := K) ij.1.1 ij.1.2
\end{leancode}

\begin{dictionary}
\lean{MvPolynomial I K} is $K[x_i:i\in I]$ and \lean{X i} the variable $x_i$. \lean{rename f} substitutes $x_a\mapsto x_{f(a)}$.
\lean{collisionIndex i j hji} sends $i$ to $j$ and every other $a$ to itself, landing in the subtype \lean{{a : I // a ≠ i}} of
indices other than $i$; so \lean{collisionMap} is $x_i:=x_j$ with $x_i$ removed from the variable set. \lean{Equiv.swap i j} is the
transposition. \lean{IncreasingPair I} is the subtype of pairs with first entry smaller.
\end{dictionary}

\begin{physics}
In a functional realisation, the labelled variables are positions of quasi-particles, and a collision $x_i=x_j$ is two particles at
the same point. ``Vanishing at a collision'' is the functional form of an OPE relation (``this normal-ordered product is null'').
\end{physics}

\begin{status}
Definitions.
\end{status}

\begin{theorem}[Collision divisibility; square factors from symmetry]\label{l20:thm:collision}
Let $K$ be a field and $p\in K[x_I]$, $i\ne j$.
\begin{enumerate}
\item $x_i-x_j$ divides $p$ iff $p$ vanishes after the substitution $x_i:=x_j$.
\item $x_i-x_j$ is prime, and for two different unordered pairs the collision factors do not divide each other.
\item If $\operatorname{char}K=0$, $p$ is symmetric under $x_i\leftrightarrow x_j$, and $p$ vanishes on $x_i=x_j$, then $(x_i-x_j)^2$
  divides $p$.
\item If $\operatorname{char}K=0$ and $p$ is symmetric under every transposition and vanishes on every diagonal, the square of the
  Vandermonde divides $p$.
\end{enumerate}
\end{theorem}
\noindent File \leanpath{VOA/\allowbreak Commutative/\allowbreak CollisionCommonFactors.lean}.
\begin{leancode}
theorem collisionFactor_dvd_iff (i j : I) (hji : j ≠ i)
    (p : MvPolynomial I K) :
    collisionFactor (K := K) i j ∣ p ↔ collisionMap (K := K) i j hji p = 0 := by

theorem collisionFactor_prime (i j : I) (hji : j ≠ i) :
    Prime (collisionFactor (K := K) i j) := by

theorem collisionFactor_not_dvd_of_unordered_pairs_ne
    {a b c d : I} (hba : b ≠ a) (hcd : c ≠ d)
    (hdirect : ¬(c = a ∧ d = b)) (hreverse : ¬(c = b ∧ d = a)) :
    ¬collisionFactor (K := K) a b ∣ collisionFactor (K := K) c d := by

theorem collisionFactor_sq_dvd_of_swap_invariant [CharZero K]
    (i j : I) (hji : j ≠ i)
    (p : MvPolynomial I K)
    (hsymm : swapVariables (K := K) i j p = p)
    (hvanish : collisionMap (K := K) i j hji p = 0) :
    collisionFactor (K := K) i j ^ 2 ∣ p := by

theorem vandermonde_sq_dvd_of_symmetric_diagonal_vanishing [CharZero K]
    (p : MvPolynomial I K)
    (hsymm : ∀ i j : I, swapVariables (K := K) i j p = p)
    (hvanish : ∀ i j : I, ∀ hji : j ≠ i,
      collisionMap (K := K) i j hji p = 0) :
    vandermonde (K := K) (I := I) ^ 2 ∣ p := by
\end{leancode}

\begin{dictionary}
Namespace \lean{VOA.CollisionCommonFactors}; \lean{K} a field and \lean{I} a type with decidable equality (section variables); the
Vandermonde statement also needs \lean{[Fintype I] [LinearOrder I]}. \lean{a ∣ b} is divisibility, \lean{Prime} the Mathlib notion of
a prime element. The hypotheses \lean{hdirect}, \lean{hreverse} say $\{c,d\}\ne\{a,b\}$ as unordered pairs. \lean{[CharZero K]} is
needed in items~3--4 and is used exactly once, to cancel a factor $2$. These are exactly the textbook statements.
\end{dictionary}

\textbf{Proof on paper.} (1) Isolate $x_i$: $K[x_I]=A[x_i]$ with $A=K[x_{I\setminus i}]$; the substitution is evaluation at
$x_j\in A$, and $X-a$ divides $f\in A[X]$ iff $f(a)=0$. (2) $x_i-x_j$ corresponds to $X-x_j$, prime in $A[X]$. Two collision factors for
different unordered pairs are not associates: under the substitution $x_a:=x_b$, which kills $x_a-x_b$, the factor $x_c-x_d$ becomes
a \emph{nonzero} polynomial (because $\{c,d\}\ne\{a,b\}$, it is $x_c-x_d$ with at most one index renamed, and the two indices remain
different); an associate of $x_a-x_b$ would be killed too. (3) Write $p=(x_i-x_j)q$. Swapping gives
$p=-(x_i-x_j)\,\mathrm{swap}(q)$, so $q$ is antisymmetric; on the diagonal $q=-q$, so $2q|_{x_i=x_j}=0$ and $q$ vanishes there; apply (1)
to $q$. (4) Each squared factor divides $p$ by (3); they are pairwise non-associated primes in a UFD, so their product divides $p$.
The library notes that the right tool is non-association of primes, not comaximality of the principal ideals, which would be false. \qed

\begin{physics}
Item~3 is the familiar fact that for identical bosons a relation forcing a wave function to vanish at coincident points forces a
\emph{double} zero: symmetric functions vanish to even order on diagonals. Item~4 is, heuristically, the origin of the $\Delta(s)^2$ in
quasi-particle character formulas, the $q^{d(d-1)}$-type shifts.
\end{physics}

\begin{status}
\textbf{Proved, unconditional} (isolated-module check; aggregate build not re-run) (the hypotheses are the statement's own); \texttt{\#print axioms} of all five:
\lean{[propext, Classical.choice, Quot.sound]}. \leanpath{CollisionCommonFactors} is one of the five modules reused by
Chapter~\ref{ch:L24} (opening of this section).
\end{status}

\subsection{The rooted Gordon conductor and prime-power assembly}

\textbf{Colours, slots, roots (motivation; the Lean statements below are purely about polynomials).} In the functional model of Gordon-type relations (Feigin--Stoyanovsky), the basic objects are
\emph{quasi-particles of colour $c$}: a quasi-particle of colour $c$ stands for $c$ coincident modes of the same current (a cluster of
size $c$). A \emph{colour profile} is a function $\rho:\{0,\dots,n-1\}\to\N$, read as ``$\rho(j)$ quasi-particles of colour $j$''; the
set of \emph{slots} is the set of pairs $(j,k)$ with $k<\rho(j)$, one variable $x_i$ per slot $i$, and $c_i$ denotes the colour of slot $i$.
The relations of the algebra translate into two kinds of vanishing for the dual polynomials: at the \emph{root} $x_i=0$ to order $c_i$
(the ``rooted'' condition, coming from the absence of the lowest modes) and on each collision $x_i=x_j$ to order $2\min(c_i,c_j)$ (Gordon's
difference condition between clusters). The \emph{rooted Gordon conductor} is the product of all these forced factors. The $p=3$ conductor
$C_{N,d}$ of \eqref{l20:eq:p3conductor} has the same shape for the $S$ particles among themselves (each behaves as a slot of colour one:
root order $1$, pair order $2\min(1,1)=2$), but its current--$S$ factors have order $1$, so it is an analogue, not a special case.

\begin{theorem}[Product of separated prime powers; the rooted Gordon conductor divides]\label{l20:thm:rooted}
(1) In a cancellative commutative monoid with zero, if $f_a$ ($a\in s$, $s$ finite) are primes no two of which divide each other, and
$f_a^{e_a}\mid p$ for each $a$, then $\prod_af_a^{e_a}\mid p$. (2) For a colour profile $\rho$ with slots $i$ of colour $c_i$, the rooted
Gordon conductor $\delta_\rho=\prod_ix_i^{c_i}\prod_{i<j}(x_i-x_j)^{2\min(c_i,c_j)}$ divides every polynomial that is divisible by each
displayed prime power.
\end{theorem}
\noindent File \leanpath{VOA/\allowbreak Commutative/\allowbreak RootedGordonConductor.lean}.
\begin{leancode}
def slotColor {n : ℕ} {rho : Fin n → ℕ} (i : ColoredSlots rho) : ℕ :=
  i.1.1

def rootedGordonConductor {n : ℕ} (rho : Fin n → ℕ) :
    MvPolynomial (ColoredSlots rho) K :=
  rootedGordonRootFactor (K := K) rho * rootedGordonPairFactor (K := K) rho
\end{leancode}
\noindent File \leanpath{VOA/\allowbreak Commutative/\allowbreak RootedIncidence.lean}.
\begin{leancode}
def ColoredSlots {n : ℕ} (rho : Fin n → ℕ) :=
  Σ j : Fin n, Fin (rho j)
\end{leancode}
\noindent File \leanpath{VOA/\allowbreak Commutative/\allowbreak RootedGordonConductorDivisibility.lean}.
\begin{leancode}
theorem prod_pow_dvd_of_pairwise_prime_varying
    {A M : Type*} [DecidableEq A] [CommMonoidWithZero M] [IsCancelMulZero M]
    (s : Finset A) (f : A → M) (e : A → ℕ) (p : M)
    (hprime : ∀ a ∈ s, Prime (f a))
    (hseparate : ∀ a ∈ s, ∀ b ∈ s, a ≠ b → ¬ f a ∣ f b)
    (heach : ∀ a ∈ s, f a ^ e a ∣ p) :
    (∏ a ∈ s, f a ^ e a) ∣ p := by

theorem rootedGordonConductor_dvd_of_root_and_collision_primePowers {n : ℕ} (rho : Fin n → ℕ)
    (p : MvPolynomial (ColoredSlots rho) K)
    (hroot : ∀ i : ColoredSlots rho, MvPolynomial.X i ^ slotColor i ∣ p)
    (hcollision : ∀ ij : IncreasingPair (ColoredSlots rho),
      collisionFactor (K := K) ij.1.1 ij.1.2 ^
        (2 * min (slotColor ij.1.1) (slotColor ij.1.2)) ∣ p) :
    rootedGordonConductor (K := K) rho ∣ p := by
\end{leancode}

\begin{dictionary}
\lean{ColoredSlots rho} is the set of pairs $(j,k)$ with $j<n$ and $k<\rho(j)$: $\rho(j)$ slots of colour $j$, and
\lean{slotColor} reads off $j$. \lean{rootedGordonRootFactor} is $\prod_ix_i^{c_i}$ and \lean{rootedGordonPairFactor} is
$\prod_{i<j}(x_i-x_j)^{2\min(c_i,c_j)}$ (definitions in the same file). The order on slots used to orient pairs is an arbitrary enumeration (a local instance); this is
harmless, since each pair factor carries the even exponent $2\min(c_i,c_j)$, so reversing a pair does not change the product.
\lean{IsCancelMulZero M} is cancellation of nonzero factors; polynomial rings over a field qualify. The separation hypothesis of
the general lemma is discharged inside the conductor theorem (\lean{rootedGordonPrimeFactor_separate}). \lean{rootedGordonConductor_ne_zero}
(same file) proves $\delta_\rho\ne0$.
\end{dictionary}

\textbf{Proof on paper.} (1) Induction on $s$: if $f_a^{e_a}q=p$ and $f_b^{e_b}\mid p$ for $b\ne a$, then $f_b^{e_b}\mid q$: the prime $f_b$ does not divide
$f_a$, hence (a prime dividing a power divides the base) does not divide $f_a^{e_a}$, so all the factors $f_b$ of $p = f_a^{e_a}q$ must
come from $q$. (2) The displayed factors are primes ($x_i$, and $x_i-x_j$ by Theorem~\ref{l20:thm:collision}), pairwise
non-associated; apply (1). \qed

\begin{physics}
Gordon's ``difference-two at distance $k$'' conditions become, in the functional model, prescribed orders of vanishing where quasi-particles
of given colours collide: two quasi-particles of colours $c_i,c_j$ repel with a zero of order $2\min(c_i,c_j)$ (heuristic picture). The conductor collects all
these zeros; dividing it out is the functional form of the shift in a fermionic character (heuristic).
\end{physics}

\begin{status}
\textbf{Proved, unconditional} (isolated-module check; aggregate build not re-run); axioms \lean{[propext, Classical.choice, Quot.sound]} (also \lean{rootedGordonConductor_ne_zero}). The general
lemma (1) is reused by Chapter~\ref{ch:L24} (\leanpath{DeligneAuxiliaryCollision}).
\end{status}

\subsection{The degree of the \texorpdfstring{$p=3$}{p=3} conductor}

\begin{theorem}[The weighted-degree count of the $p=3$ conductor]\label{l20:thm:p3degree}
The number of increasing pairs in $\{0,\dots,d-1\}$ is $\binom d2$, and the counting function
\lean{p3ConductorWeightedDegree} (root weight plus collision multiplicities) of the labelled conductor
\eqref{l20:eq:p3conductor} in multidegree $(a,b,c,d)$ equals $d(a+b+c+d)$.
\end{theorem}
\noindent File \leanpath{VOA/\allowbreak Commutative/\allowbreak P3RestrictedDualCommonFactor.lean}.
\begin{leancode}
abbrev ConductorCollisionIndex (a b c d : ℕ) :=
  IncreasingPair (Fin d) ⊕ (BoundaryCurrentLabel a b c × Fin d)

def conductorCollisionMultiplicity {a b c d : ℕ} :
    ConductorCollisionIndex a b c d → ℕ
  | .inl _ => 2
  | .inr _ => 1
\end{leancode}
\noindent File \leanpath{VOA/\allowbreak Commutative/\allowbreak P3ResidualWindowPreservation.lean}.
\begin{leancode}
def p3ConductorWeightedDegree (a b c d : ℕ) : ℕ :=
  Finsupp.weight
      (particleWeightOne : SpeciesLabel (boundaryMultiplicity a b c d) → ℕ)
      (speciesShiftExponent (boundaryMultiplicity a b c d) boundarySpeciesShift) +
    ∑ i : ConductorCollisionIndex a b c d, conductorCollisionMultiplicity i
\end{leancode}
\noindent File \leanpath{VOA/\allowbreak Commutative/\allowbreak P3ConductorDegree.lean}.
\begin{leancode}
theorem card_increasingPair_fin (d : ℕ) :
    Fintype.card (IncreasingPair (Fin d)) = d.choose 2 := by

theorem p3ConductorWeightedDegree_eq (a b c d : ℕ) :
    p3ConductorWeightedDegree a b c d = d * (a + b + c + d) := by
\end{leancode}

\begin{dictionary}
\lean{BoundaryCurrentLabel a b c} is the disjoint union $\mathrm{Fin}\,a\sqcup\mathrm{Fin}\,b\sqcup\mathrm{Fin}\,c$ of current labels;
\lean{ConductorCollisionIndex} indexes the collision factors of \eqref{l20:eq:p3conductor}: $S$--$S$ pairs (multiplicity $2$) and
current--$S$ pairs (multiplicity $1$). \lean{SpeciesLabel (boundaryMultiplicity a b c d)} is the set of all $a+b+c+d$ labelled
particles; \lean{speciesShiftExponent … boundarySpeciesShift} is the exponent vector with $1$ on each $S$ label and $0$ on currents (the
monomial $\prod_\alpha s_\alpha$), and \lean{Finsupp.weight particleWeightOne} its total degree, which is $d$
(\lean{weight_boundarySpeciesShift}). So \lean{p3ConductorWeightedDegree} is $d+2\binom d2+(a+b+c)d$ written as a sum over the factors,
and the theorem evaluates it. The definition counts factors and does not mention the polynomial. That the conductor polynomial \lean{p3CommonFactor} is
weighted-homogeneous of this degree is the separate theorem \lean{p3CommonFactor_isWeightedHomogeneous}
(\leanpath{P3ResidualWindowPreservation.lean}, not quoted here); ``has degree exactly $dN$'' would also need the conductor to be
nonzero, which that theorem does not state. So the theorem printed here is the arithmetic of the count, and the title promises no more.
\end{dictionary}

\textbf{Proof on paper.} Increasing pairs in a $d$-element ordered set are $2$-subsets, so there are $\binom d2$ (the library goes
through off-diagonal points of the symmetric square). Then $d+2\binom d2+(a+b+c)d=d+d(d-1)+(a+b+c)d=d(a+b+c+d)$. \qed

\begin{physics}
$dN$ is the total ``energy shift'' that the $d$ auxiliary particles cost in a sector with $N$ particles: each $S$ particle repels every
other particle once (and other $S$ particles twice) and cannot sit at the origin. In a fermionic formula this is the exponent of $q$
multiplying the free part of the sector.
\end{physics}

\begin{status}
\textbf{Proved, unconditional} (isolated-module check; aggregate build not re-run); axioms \lean{[propext, Classical.choice, Quot.sound]}. \lean{card_increasingPair_fin} is reused by
Chapter~\ref{ch:L24} (\leanpath{DeligneAuxiliaryDegree}).
\end{status}

\subsection{The remaining files, collectively}

All of the following are finite statements about labelled polynomials, finite sets and group actions, with no hypothesis beyond
those displayed in each statement; none mentions a vertex algebra. Every file's docstring states what it does \emph{not} prove.
\begin{center}
\begingroup\small
\begin{longtable}{p{0.33\textwidth}p{0.47\textwidth}p{0.11\textwidth}}
\toprule
Files & What is proved, in words & Used by L24\\
\midrule\endhead
\leanpath{ModeParticlePairing}, \leanpath{ModeParticlePolynomial}, \leanpath{MultispeciesParticlePolynomial} &
The dictionary of the functional model: monomials in field modes are paired with symmetric polynomials in particle positions,
the pairing is perfect (a duality), and it is compatible with taking several species at once. &
first two\\
\leanpath{AuxiliaryIdealOrthogonal} & The relations $S^n=0$, $A\cdot S^n=0$ and $XZ=Y^2$ translate, under that duality, into the
vanishing (resp.\ agreement) conditions on diagonals listed in ``The model, on paper''. & --\\
\leanpath{BoundaryTypedSubstitution}, \leanpath{BoundaryCoefficientFiber}, \leanpath{BoundaryCollisionDivisibility} & How the
coupled cone condition ``$F_{XZ}$ on $x=z$ equals $F_{YY}$ on $y_1=y_2$'' is expressed by substitutions of variables; it yields
an \emph{equality} of two restrictions, not a vanishing, so no factor of the conductor is deduced from it. & --\\
\leanpath{P3RestrictedDualCommonFactor}, \leanpath{P3AuxiliaryCollisionBridge}, \leanpath{P3CommonFactorInvariant},
\leanpath{P3CommonFactorSourceInvariant}, \leanpath{P3BoundaryRootEquiv}, \leanpath{P3InvariantPolynomialRealization} & The
conductor $C_{N,d}$ divides every polynomial satisfying the collision vanishings; it and the quotients by it are symmetric under
relabelling particles of one species; bookkeeping of labels. & --\\
\leanpath{P3ResidualWindowPreservation}, \leanpath{P3ResidualConeBridge}, \leanpath{P3ConeCarrierInterface}, \leanpath{P3FunctionalModel} &
Dividing by the conductor lowers the degree by its degree and keeps the symmetry and the cone condition; the conditional dimension
bound \lean{finrank_le_expectedDim} (three uninstantiated arguments; see the warning above). & --\\
\leanpath{RootedMatching}, \leanpath{RootedProfile}, \leanpath{RootedTypedProfile}, \leanpath{RootedIncidence},
\leanpath{RootedIncidenceAction}, \leanpath{RootedIncidenceOrbit}, \leanpath{RootedLeafRelabeling}, \leanpath{RootedSourceRelabeling},
\leanpath{RootedGordonConductor} & Finite combinatorics of how particles can coincide (which particles sit at the same point, and
with which colours), and the statement that two such coincidence patterns with the same colour counts differ only by a relabelling
of particles. & --\\
\leanpath{RootedGordonConductorDivisibility} & Theorem~\ref{l20:thm:rooted}. & yes\\
\leanpath{CollisionCommonFactors} & Theorem~\ref{l20:thm:collision}. & yes\\
\leanpath{P3ConductorDegree} & Theorem~\ref{l20:thm:p3degree}. & yes\\
\leanpath{OrbitSectionAssembly}, \leanpath{OrbitEvaluationDescent}, \leanpath{LiteralSubstitutionNewton} & Generic tools: a
symmetric family of linear maps has a symmetric common kernel; a function constant on fibres factors through the quotient; a count of
exponent vectors with restricted support. & --\\
\bottomrule
\end{longtable}
\endgroup
\end{center}

\begin{physics}
Collectively: the functional (quasi-particle) model of the $p=3$ auxiliary algebra, in which states are dual to symmetric
polynomials in particle positions and OPE relations become vanishing conditions at collisions. The intended physics output, never
reached in Lean by this route, was the fermionic character of $L_{-4/3}(\ltwentysl)$, which Section~\ref{l20:sec:m4} obtains by the
Gordon route.
\end{physics}

\begin{status}
No file declares an axiom or contains \lean{sorry} or \lean{native_decide} (text search of the 32 files). The \texttt{\#print axioms}
spot checks above are isolated-module runs on \leanpath{P3ConductorDegree} and \leanpath{RootedGordonConductorDivisibility}, all
\lean{[propext, Classical.choice, Quot.sound]}; aggregate build not re-run (Section~\ref{l20:sec:verification}). The route's terminal bound is conditional on three uninstantiated arguments (the datum \lean{D} and two named
hypotheses; see the warning at the head of this section) and is not used by any headline.
\end{status}

%======================================================================
\section{Verification record}\label{l20:sec:verification}
%======================================================================

\textbf{What the machine confirmed (2026-10-13; not re-run for this edition).} On the working tree of that date, before the
campaign files were committed:
\begin{plaincode}
lake build                                   # aggregate: green, 13150 jobs
lake build VOA.ArakawaDai                    # umbrella: green (ledger: 4,513 jobs)
scripts/audit/ArakawaDaiAxiomAudit.lean      # 160 declarations: PASS, zero discrepancies
  headline axioms: [propext, Classical.choice, Quot.sound]   (no VOA.pbw)
\end{plaincode}
The audit file lists, among others, all ten headline declarations, the \lean{...OfCharacter} helpers, and the
supporting $q$-series, Gordon, determinant and PVA lemmas. The campaign ledger also records a repository check: all Lean
files reachable, no new axiom, no \lean{sorry}/\lean{admit}/\lean{native_decide}.

\textbf{What this chapter's author checked by hand.} (1) Every \texttt{leancode} block above was compared with the working-tree
source (a script replicating \texttt{sync.py}'s verbatim-chunk test against the working tree, including untracked files).
(2) The paper's displayed Theorem~1.2 (arXiv:2609.33348v1) is the unshifted, unflavoured character; Lean's \lean{character}
is the same. (3) The PVA axioms in \lean{PoissonVertexAlgebra} were compared with the $\lambda$-bracket axioms
(Definition~\ref{l20:def:pva}). (4) Coefficients of both sides of Theorem~1.2 for $m=1,2$, $N\le14$
(Example~\ref{l20:ex:numbers}).

\textbf{Sections~\ref{l20:sec:vir} and~\ref{l20:sec:p3} (2026-10-27; not re-run for this edition).} The aggregate
\texttt{lake build} was \emph{not} re-run for these sections. The author ran, from the repository root, on the existing compiled
modules, one scratch file per module with \texttt{lake env lean -DmaxSynthPendingDepth=3}:
\begin{plaincode}
import VOA.Examples.VirasoroMinimalPoisson      # 18 declarations
import VOA.Commutative.P3ConductorDegree        # 7 declarations
import VOA.Commutative.RootedGordonConductorDivisibility   # 3 declarations
  every one: [propext, Classical.choice, Quot.sound]
\end{plaincode}
The 18 Virasoro declarations are \lean{liFiltration_eq_virWordPiece}, \lean{simpleLiFiltration_eq_virWordPiece},
\lean{classicallyFree}, \lean{classicallyFree_M23}, \lean{classicallyFree_leeYang}, \lean{canonicalEquiv}, \lean{baseC2AlgEquiv},
\lean{actualC2ArcEquiv}, \lean{cm_eq}, \lean{cm_one}, \lean{jetMap_injective}, \lean{mk₂_omegaBar_pow_eq_zero}, \lean{arcPoissonEquiv},
\lean{liPoissonVertexAlgebra_mode_eq_zero}, \lean{mode_x_x}, \lean{virSymbol_twoSeries}, \lean{twoSeries_hasJetRelations_all},
\lean{zeroPVA}. The project's internal build record of 2026-10-25 reports green targeted builds of the Virasoro files (9707 jobs) and of
\texttt{VOA.ArakawaDai VOA.Examples} (10562 jobs), and a 38-declaration axiom probe with the same three axioms (not re-run for this chapter). Every new
\texttt{leancode} block was checked verbatim by \texttt{sync.py check}. The same record (2026-10-25, the
classical-freeness audit) is where the description of \lean{virSymbol_twoSeries} as ``classical freeness'' was corrected.

\textbf{Referee revision (an earlier commit; not re-run for this edition).} The fidelity referee could not load the
Virasoro modules (a missing \texttt{.olean}); at a later commit the compiled modules were present, and the reviser re-ran, with the same command on one scratch
file importing \leanpath{VOA/\allowbreak Examples/\allowbreak VirasoroMinimalPoisson.lean}:
\begin{plaincode}
'VOA.VirasoroClassicalFreeness.classicallyFree' depends on axioms: [propext, Classical.choice, Quot.sound]
'VOA.VirasoroClassicalFreeness.canonicalEquiv' depends on axioms: [propext, Classical.choice, Quot.sound]
'VOA.VirasoroClassicalFreeness.baseC2AlgEquiv' depends on axioms: [propext, Classical.choice, Quot.sound]
'VOA.VirasoroClassicalFreeness.actualC2ArcEquiv' depends on axioms: [propext, Classical.choice, Quot.sound]
'VOA.VirasoroClassicalFreeness.arcPoissonEquiv' depends on axioms: [propext, Classical.choice, Quot.sound]
'VOA.VirasoroClassicalFreeness.liPoissonVertexAlgebra_mode_eq_zero' depends on axioms: [propext, Classical.choice, Quot.sound]
'VOA.TwoSeries.DSSlice.virSymbol_twoSeries' depends on axioms: [propext, Classical.choice, Quot.sound]
'VOA.TwoSeries.twoSeries_hasJetRelations_all' depends on axioms: [propext, Classical.choice, Quot.sound]
'VOA.VirasoroLi.liFiltration_eq_virWordPiece' depends on axioms: [propext, Classical.choice, Quot.sound]
\end{plaincode}
(line breaks of the output joined). The aggregate \texttt{lake build} was again not run.

\textbf{Theorem~1.1 in the paper's form (2026-10-28; not re-run for this edition).} The new file
\leanpath{VOA/\allowbreak Examples/\allowbreak AffineSl2BoundaryTheorem11.lean} (Section~\ref{l20:sec:ad11}) is committed, but its
compiled module was \emph{missing} from the main checkout's \texttt{.lake} (no \texttt{lake build} was run for this revision). The
reviser therefore ran \texttt{lake env lean -DmaxSynthPendingDepth=3} on a scratch file importing
\texttt{VOA.Examples.AffineSl2BoundaryTheorem11} from the root of the \texttt{release-cheap} worktree, whose compiled modules
include it (built 2026-10-27) and whose copy of the file is byte-identical to the main checkout's (\texttt{cmp}). Output (line
breaks joined; peak RSS 3.7\,GB):
\begin{plaincode}
'VOA.AffineSl2BoundaryClassicalFreeness.arakawaDai_theorem_1_1' depends on axioms: [propext, Classical.choice, Quot.sound]
'VOA.AffineSl2BoundaryClassicalFreeness.canonicalC2Map' depends on axioms: [propext, Classical.choice, Quot.sound]
'VOA.AffineSl2BoundaryClassicalFreeness.baseC2AlgEquiv_apply' depends on axioms: [propext, Classical.choice, Quot.sound]
'VOA.AffineSl2BoundaryClassicalFreeness.canonicalC2Map_mk_X' depends on axioms: [propext, Classical.choice, Quot.sound]
'VOA.AffineSl2BoundaryClassicalFreeness.canonicalC2Map_bijective' depends on axioms: [propext, Classical.choice, Quot.sound]
'VOA.AffineSl2BoundaryClassicalFreeness.c2UniversalArc' depends on axioms: [propext, Classical.choice, Quot.sound]
'VOA.AffineSl2BoundaryClassicalFreeness.c2UniversalArc_universal' depends on axioms: [propext, Classical.choice, Quot.sound]
'VOA.AffineSl2BoundaryClassicalFreeness.c2ToArc' depends on axioms: [propext, Classical.choice, Quot.sound]
'VOA.AffineSl2BoundaryClassicalFreeness.c2ToArc_eq_constantCoeff' depends on axioms: [propext, Classical.choice, Quot.sound]
'VOA.AffineSl2BoundaryClassicalFreeness.arcPoissonEquiv_c2ToArc' depends on axioms: [propext, Classical.choice, Quot.sound]
'VOA.AffineSl2BoundaryClassicalFreeness.eq_arcPoissonEquiv_of_c2ToArc' depends on axioms: [propext, Classical.choice, Quot.sound]
'VOA.AffineSl2BoundaryClassicalFreeness.presentedJet_algHom_ext' depends on axioms: [propext, Classical.choice, Quot.sound]
\end{plaincode}
The build record for this file is the project's (\texttt{Notes/RELEASE\_PLAN.md}, ``Caveats removed in Lean'', 2026-10-28:
merged from \texttt{release-cheap}, umbrellas green, axioms standard). The aggregate \texttt{lake build} of the release commit is recorded in Chapter~\ref{ch:15},
Section~\ref{ch15:sec:record}.

\begin{warning}[title={Book tooling caveat (historical) and superseded statements elsewhere}]
\texttt{scripts/textbook/sync.py check} builds its verbatim corpus from files tracked at \texttt{HEAD}. The campaign
files were untracked on 2026-10-13, when this chapter was first written; they were committed the same day, and \texttt{sync.py check} now verifies this chapter's quotes against the committed sources. Several earlier chapters (L09, L11, L16) contain statements now superseded by this campaign: the boundary
presentation is no longer conditional (L09, Section~\ref{l09:sec:boundary}); Li's PVA structure is now defined (L11,
which carries an update note); the
Li telescope \lean{sum_liWeightLayer_finrank_eq_grading} is hypothesis-free (as L11 already quotes it); and for $m=1$
the content of \lean{Sl2ExactSymbol} (L16) is now true on paper, though no Lean bridge exists.
\end{warning}

%======================================================================
\textbf{Release audit (2026-10-27).} Release audit 2026-10-27: see Notes/Textbook/referee/RELEASE\_AUDIT\_L20.md. The
auditor re-ran \texttt{\#print axioms} on the ten headlines (\lean{[propext, Classical.choice, Quot.sound]}, no project
axiom), traced every definition to its leaves, and recomputed the product coefficients for $m=1,2$; verdict: no blocker,
the character theorems announceable as stated, Theorem~1.1 in a two-part wording, since at the time no single
theorem had the paper's shape. That wording is superseded by \lean{arakawaDai_theorem_1_1} (Theorem~\ref{l20:thm:ad11}),
which also discharges the audit's findings F2 (no single theorem) and F4 (no \lean{_apply} lemma for \lean{baseC2AlgEquiv}). The audit used compiled modules from 2026-10-25 and did not re-run the aggregate
\texttt{lake build}; for the release build see Chapter~\ref{ch:15}, Section~\ref{ch15:sec:record}.

%======================================================================
\section{Findings, oddities and Lean defects}\label{l20:sec:defects}
%======================================================================

No statement-level defect was found: the headlines say what the paper says, about the actual objects, with no hidden
hypothesis. The following are recorded for the referee.

\begin{enumerate}
\item \textbf{(Type-level weakness; resolved 2026-10-28) \lean{actualC2ArcEquiv}.} Its type is a bijection of sets for each
  $A$ and does not itself state naturality in $A$. Naturality is now proved as a separate theorem in all three places:
  \lean{AffineSl2BoundaryClassicalFreeness.actualC2ArcEquiv_natural},
  \lean{AffineSl2BoundarySingularFamily.Candidate.actualC2ArcEquiv_natural} and
  \lean{VirasoroClassicalFreeness.actualC2ArcEquiv_natural} (with the helper
  \lean{AffineSl2BoundaryClassicalFreeness.actualC2ArcEquivOfCharacter_natural}). Together with its naturality theorem,
  \lean{actualC2ArcEquiv} is a representability statement, and the docstrings now say so. The bundled form is
  \lean{c2UniversalArc_universal} (Theorem~\ref{l20:thm:univarc}), which names the bijection ($f\mapsto f_*\circ\gamma$).
\item \textbf{(Documentation; resolved 2026-10-28) \lean{baseC2AlgEquiv}'s type is an abstract isomorphism}; that it is
  the canonical map $x_i\mapsto[\text{current}_i]$ was by definition only, and is now a theorem (\lean{baseC2AlgEquiv_apply},
  \lean{canonicalC2Map_mk_X}; Theorem~\ref{l20:thm:canonc2}). The literal kernel statement \lean{Candidate.baseIdeal_eq_simpleSymbol_kernel}
  exists only in a form with the hypothesis \lean{hinj}; an unconditional corollary is not stated in the headline file.
  Attribution: the paper displays only the surjection $\ltwentyA\twoheadrightarrow R_V$; the isomorphism is a corollary
  of Theorem~1.1, not a displayed claim of the paper.
\item \textbf{(Wording; resolved 2026-10-28) Theorem~1.1 had no single Lean statement.} It was formalised as
  \lean{baseC2AlgEquiv} plus \lean{boundary_classically_free}/\lean{arcPoissonEquiv} on the presented arc algebra. It is now
  the single theorem \lean{arakawaDai_theorem_1_1} (Theorem~\ref{l20:thm:ad11}); cite that. What remains is the honest caveat
  that $\ltwentyJ R_V$ is realised as a presented arc algebra with a universal arc (no $\ltwentyJ$ functor on abstract
  algebras), characterised up to unique isomorphism by \lean{c2UniversalArc_universal}. The two older headlines
  \lean{canonicalMap_bijective} and \lean{boundary_classically_free} still state the same fact on two isomorphic
  presentations.
\item \textbf{(Docstring) \lean{character_eq_chain}} is described as ``the literal descending-chain expression'' of
  Theorem~1.2; the literal display is \lean{character_eq_fermionic_chain}; \lean{hilbertBoundSeries} is the paper's
  intermediate $S_m$.
\item \textbf{(Stale docstrings).} \leanpath{AffineSl2BoundaryVacuumCharacter.lean} calls the character equality ``the
  remaining character gate'' / ``the remaining input'' (lines 83, 95); \leanpath{AffineSl2SimpleCharacterRank.lean}
  (line 12) calls the full character product ``unresolved''; \leanpath{AffineInfiniteClassicalFreeness.lean} (line 739)
  speaks of a ``remaining input''. All are now proved.
\item \textbf{(Naming)} Two pairs of declarations share a short name with different meanings:
  \lean{sourcePiece_finrank_le_product} (index $m+1$ in \leanpath{GordonConformalCharacter.lean}, index $m$ in
  \leanpath{AffineSl2BoundaryVacuumCharacter.lean}) and \lean{polynomialMap} (\lean{BoundaryArcLiMap} versus
  \lean{BoundaryArcConformalSqueeze}, different arguments). Also, \lean{coeff_product_le_character} carries the unused
  hypothesis \lean{hm} (harmless), and only the inclusion \lean{baseIdeal_le_boundaryIdeal}, not the equality, is stated.
\item \textbf{(Not a defect; noted)} \lean{polynomialPVA} is built by transport from $\ltwentygr V^0(\ltwentysl)$ (level $0$),
  which looks arbitrary; \lean{arcPVA_unique} shows the result is canonical.
\item \textbf{(Not a defect; noted)} \lean{uniformCandidate} is a \lean{Classical.choose}; the headline maps do not depend
  on the choice.
\item \textbf{(Not a defect; noted)} \lean{hilbertBoundSeries} and \lean{fermionicTrace} use \lean{m - 1} in \lean{ℕ}, which
  is meaningless at $m=0$; all uses assume $1\le m$.
\end{enumerate}

\textbf{Virasoro companion and $p=3$ files (added 2026-10-27).} Again no statement-level defect: the Virasoro headlines are
unconditional and about the actual simple quotient. Recorded for the referee:
\begin{enumerate}\setcounter{enumi}{9}
\item \textbf{(Type-level, as item 1; naturality resolved)} \lean{VirasoroClassicalFreeness.actualC2ArcEquiv} is natural in $A$
  (\lean{VirasoroClassicalFreeness.actualC2ArcEquiv_natural}), and
  \lean{VirasoroClassicalFreeness.baseC2AlgEquiv} is canonical only by definition (Theorem~\ref{l20:thm:virc2}).
\item \textbf{(Missing bridge)} \lean{VirasoroClassicalFreeness.cm m} is \lean{(twoSeriesLevels m).c}; only the closed form
  \lean{cm_eq} relates it to $c_{2,2m+3}$. No theorem equates it with \lean{minimalCentralCharge 2 (2 * m + 3)} of Chapter~\ref{ch:L06}.
\item \textbf{(Missing corollaries)} No Virasoro analogue of \lean{arcPoissonEquiv_conformal_iff} (weight preservation) and no character
  identity is stated, although both follow from the construction.
\item \textbf{(Naming)} \lean{VirasoroClassicalFreeness.arcPVA}, \lean{.arcPoissonEquiv}, \lean{.baseC2AlgEquiv} and
  \lean{.actualC2ArcEquiv} share short names with the affine \lean{arcPVA}, \lean{arcPoissonEquiv}, \lean{baseC2AlgEquiv},
  \lean{actualC2ArcEquiv}, with different types; only the namespace distinguishes them.
\item \textbf{(Scope, honest docstrings)} The 32 $p=3$ files prove finite lemmas only; the terminal bound
  \lean{P3LinearResidualization.finrank_le_expectedDim} takes the uninstantiated hypotheses \lean{P3QuadraticConeCarrierIdentification} and
  \lean{DumanskiFeiginP3GradedDimensionHypothesis}. The docstrings say so. \leanpath{RootedGordonConductor.lean} and
  \leanpath{RootedGordonConductorDivisibility.lean} end with \texttt{\#print axioms} commands (an \lean{AxiomAudit} section) that print at
  every build; cosmetic.
\end{enumerate}

%======================================================================
\section{Referee's checklist}\label{l20:sec:checklist}
%======================================================================

\begingroup\small
\begin{longtable}{p{0.30\textwidth}p{0.30\textwidth}p{0.22\textwidth}p{0.10\textwidth}}
\toprule
Lean declaration & Intended meaning & Hypotheses / status & Fidelity check\\
\midrule\endhead
\lean{AffineSl2BoundaryCharacter.}\allowbreak\lean{character_}\allowbreak\lean{eq_}\allowbreak\lean{product} & $\operatorname{ch}L_{k_m}=(q^p;q^p)^3_\infty/(q)^3_\infty$ & $1\le m$; unconditional, \lean{pbw}-free & Thm.~\ref{l20:thm:character}\\
\lean{character_}\allowbreak\lean{eq_}\allowbreak\lean{chain} & $\operatorname{ch}L_{k_m}=S_m$ (paper (5.12)) & $1\le m$; unconditional & Thm.~\ref{l20:thm:character}\\
\lean{character_}\allowbreak\lean{eq_}\allowbreak\lean{fermionic_}\allowbreak\lean{chain} & paper Thm.~1.2 literally: $T_m/(q)^2_\infty$ & $1\le m$; unconditional & Thm.~\ref{l20:thm:character}\\
\lean{coeff_}\allowbreak\lean{product_}\allowbreak\lean{le_}\allowbreak\lean{character} & lower bound for the character (a lemma, not a paper theorem) & $1\le m$; unconditional & Thm.~\ref{l20:thm:lower}\\
\lean{coeff_}\allowbreak\lean{character_}\allowbreak\lean{le_}\allowbreak\lean{product} & Hilbert-series upper bound & $1\le m$; unconditional & Thm.~\ref{l20:thm:upper}\\
\lean{AffineSl2BoundaryClassicalFreeness.}\allowbreak\lean{arakawaDai_}\allowbreak\lean{theorem_}\allowbreak\lean{1_}\allowbreak\lean{1} & \textbf{Arakawa--Dai Thm.~1.1 as one theorem}: $\kappa\colon\ltwentyA\to R_V$ bijective with the current-class formula; $(\ltwentyJ\ltwentyA,\gamma)$ is the arc algebra of $R_V$; the PVA isomorphism is the canonical $\ltwentyJ R_V\to\ltwentygr V$ and preserves conformal weight & $1\le m$; unconditional; $\ltwentyJ R_V$ realised as presented arc algebra with universal arc & Thm.~\ref{l20:thm:ad11}\\
\lean{canonicalC2Map}, \lean{baseC2AlgEquiv_}\allowbreak\lean{apply}, \lean{canonicalC2Map_}\allowbreak\lean{mk_}\allowbreak\lean{X}, \lean{canonicalC2Map_}\allowbreak\lean{bijective} & $\kappa$; \lean{baseC2AlgEquiv} is $\kappa$ (by theorem); $\kappa(x_i)=[\bar x_i]$ & $1\le m$; unconditional & Def.~\ref{l20:def:canonc2}, Thm.~\ref{l20:thm:canonc2}\\
\lean{c2UniversalArc}, \lean{c2UniversalArc_}\allowbreak\lean{universal}, \lean{c2ToArc} & universal arc of $R_V$; its universal property ($\exists!$); zero jets & $1\le m$; unconditional & Def.~\ref{l20:def:canonc2}, Thm.~\ref{l20:thm:univarc}\\
\lean{presentedJet_}\allowbreak\lean{algHom_}\allowbreak\lean{ext}, \lean{arcPoissonEquiv_}\allowbreak\lean{c2ToArc}, \lean{eq_}\allowbreak\lean{arcPoissonEquiv_}\allowbreak\lean{of_}\allowbreak\lean{c2ToArc} & the PVA isomorphism extends $R_V\hookrightarrow\ltwentygr V$ and is the unique such translation-compatible map & unconditional & Thm.~\ref{l20:thm:pvacanon}\\
\lean{AffineSl2BoundaryClassicalFreeness.}\allowbreak\lean{canonicalMap_}\allowbreak\lean{bijective} & $\ltwentyJ\ltwentyA\to\ltwentygr L_{k_m}$ bijective ($\ltwentyA$ presented; packaged with the rest as Thm.~\ref{l20:thm:ad11}) & $1\le m$; unconditional & Thm.~\ref{l20:thm:freeness}\\
\lean{boundary_}\allowbreak\lean{classically_}\allowbreak\lean{free} & same map, PVA-packaged source & $1\le m$; unconditional & Thm.~\ref{l20:thm:freeness}\\
\lean{arcPoissonEquiv} & PVA isomorphism (inverse compatible) & $1\le m$; unconditional & Thm.~\ref{l20:thm:freeness}, Def.~\ref{l20:def:pva}\\
\lean{arcPoissonEquiv_}\allowbreak\lean{conformal_}\allowbreak\lean{iff} & both directions preserve conformal weight & $1\le m$; unconditional & Thm.~\ref{l20:thm:freeness}\\
\lean{baseC2AlgEquiv} & $\ltwentyA\cong R_{L_{k_m}}$ (canonical by theorem: \lean{baseC2AlgEquiv_}\allowbreak\lean{apply}) & $1\le m$; unconditional & Thm.~\ref{l20:thm:c2arc}\\
\lean{actualC2ArcEquiv} & for each $A$, a bijection $\mathrm{Hom}(\ltwentyJ\ltwentyA,A)\simeq\mathrm{Hom}(R_{L_{k_m}},A[[t]])$; natural in $A$ by \lean{actualC2ArcEquiv_}\allowbreak\lean{natural}, so a representability statement (bundled: \lean{c2UniversalArc_}\allowbreak\lean{universal}) & $1\le m$; unconditional; naturality a separate theorem & Thm.~\ref{l20:thm:c2arc}\\
\lean{character}, \lean{product}, \lean{vacuum}, \lean{simple}, \lean{universal}, \lean{level}, \lean{denominator} & objects & definitions & Defs.~\ref{l20:def:level}--\ref{l20:def:character}\\
\lean{euler}, \lean{reciprocalEuler}, \lean{fermionicTrace}, \lean{hilbertBoundSeries}, \lean{LocallyFinitePowerSeriesSum.}\allowbreak\lean{sum} & $q$-series & definitions; finiteness proved & Defs.~\ref{l20:def:euler}--\ref{l20:def:chain}\\
\lean{BasePoly}, \lean{baseIdeal}, \lean{InfiniteJet.}\allowbreak\lean{PresentedJet}, \lean{jetIdeal}, \lean{arcEquiv} & $\C[e,h,f]$, $I_m$, $\ltwentyJ$, arc property & definitions; arc property proved & Def.~\ref{l20:def:jet}, Thm.~\ref{l20:thm:arc}\\
\lean{PoissonVertexAlgebra}, \lean{.}\allowbreak\lean{Hom}, \lean{.}\allowbreak\lean{Equiv}, \lean{liPoissonVertexAlgebra}, \lean{arcPVA} & PVAs & definitions; all laws proved & Defs.~\ref{l20:def:pva}, \ref{l20:def:lipva}, \ref{l20:def:arcpva}\\
\lean{uniformCandidate}, \lean{canonicalMap}, \lean{arcLiPoissonHom} & null vector; the map $\pi$ & constructions & Defs.~\ref{l20:def:candidate}, \ref{l20:def:canonical}\\
\lean{arcConformalPiece}, \lean{sourcePiece}, \lean{liConformalPiece} & conformal pieces & definitions & Def.~\ref{l20:def:conformal}\\
\lean{admissible_}\allowbreak\lean{first_}\allowbreak\lean{null}, \lean{Candidate.}\allowbreak\lean{boundaryIdeal_}\allowbreak\lean{le_}\allowbreak\lean{kernel} & boundary relations, null vector & unconditional & Thms.~\ref{l20:thm:firstnull}, \ref{l20:thm:symbol}\\
\lean{sourcePiece_}\allowbreak\lean{finrank_}\allowbreak\lean{le_}\allowbreak\lean{auxiliaryPiece} & degeneration to the auxiliary arc algebra & unconditional & \S\ref{l20:sec:m1}\\
\lean{sector_}\allowbreak\lean{finrank_}\allowbreak\lean{le_}\allowbreak\lean{profile_}\allowbreak\lean{sum}, \lean{auxiliaryPiece_}\allowbreak\lean{finrank_}\allowbreak\lean{le_}\allowbreak\lean{source} & Gordon divisibility, partition filtration & unconditional & Thms.~\ref{l20:thm:sectorbound}, \ref{l20:thm:auxbound}\\
\lean{hilbertBoundSeries_}\allowbreak\lean{product}, \lean{fermionicTrace_}\allowbreak\lean{eq}, \lean{jacobi_}\allowbreak\lean{cube} & $q$-series identities & unconditional ($1\le m$) & Thm.~\ref{l20:thm:qseries}\\
\lean{finrank_}\allowbreak\lean{ker_}\allowbreak\lean{le_}\allowbreak\lean{rootMultiplicity}, \lean{universal_}\allowbreak\lean{grading_}\allowbreak\lean{finrank_}\allowbreak\lean{le_}\allowbreak\lean{simple_}\allowbreak\lean{add_}\allowbreak\lean{rootOrder} & lower bound: rank bound & unconditional & Lem.~\ref{l20:lem:rootorder}, Thm.~\ref{l20:thm:rootbound}\\
\lean{vacuum_}\allowbreak\lean{determinant_}\allowbreak\lean{order_}\allowbreak\lean{add_}\allowbreak\lean{reflected}, \lean{vacuumRootOrderTotal_}\allowbreak\lean{add_}\allowbreak\lean{product_}\allowbreak\lean{eq_}\allowbreak\lean{universal_}\allowbreak\lean{finrank} & lower bound: determinant count & unconditional & Thms.~\ref{l20:thm:vacdet}, \ref{l20:thm:resonances}\\
\lean{sourceMap_}\allowbreak\lean{injective_}\allowbreak\lean{of_}\allowbreak\lean{hilbert_}\allowbreak\lean{bound} & conformal Hilbert squeeze & needs \lean{hbound} (discharged) & Thm.~\ref{l20:thm:squeeze}\\
\lean{poissonVertexAlgebra_}\allowbreak\lean{unique}, \lean{baseC2Map_}\allowbreak\lean{injective_}\allowbreak\lean{of_}\allowbreak\lean{arc_}\allowbreak\lean{injective} & PVA structure, intrinsic arc equivalence & unconditional / \lean{hinj} (discharged) & Thms.~\ref{l20:thm:pvahom}, \ref{l20:thm:m8}\\
\midrule
\multicolumn{4}{l}{\emph{Virasoro companion (Section~\ref{l20:sec:vir}); namespace \lean{VirasoroClassicalFreeness} unless noted}}\\
\lean{VirasoroLi.}\allowbreak\lean{liFiltration_}\allowbreak\lean{eq_}\allowbreak\lean{virWordPiece}, \lean{simpleLiFiltration_}\allowbreak\lean{eq_}\allowbreak\lean{virWordPiece} & $F^pV(c,0)=C_p$; $F^pL(c,0)=$ image of $C_p$ & every $c$; unconditional & Thm.~\ref{l20:thm:virli}\\
\lean{classicallyFree} (\lean{_M23}, \lean{_leeYang}), \lean{canonicalEquiv} & $\ltwentyJ(\C[T]/(T^{m+1}))\cong\ltwentygr L(c_{2,2m+3},0)$ & every $m\in\N$; unconditional & Thm.~\ref{l20:thm:virfree}, Def.~\ref{l20:def:vircanon}\\
\lean{baseC2AlgEquiv}, \lean{actualC2ArcEquiv}, \lean{mk₂_}\allowbreak\lean{omegaBar_}\allowbreak\lean{pow_}\allowbreak\lean{eq_}\allowbreak\lean{zero} & $R_L\cong\C[T]/(T^{m+1})$; arcs & unconditional; canonical by definition; naturality a separate theorem & Thm.~\ref{l20:thm:virc2}\\
\lean{arcPoissonEquiv}, \lean{liPoissonVertexAlgebra_}\allowbreak\lean{mode_}\allowbreak\lean{eq_}\allowbreak\lean{zero}, \lean{zeroPVA} & PVA isomorphism; all Li modes vanish & unconditional & Thm.~\ref{l20:thm:virpva}\\
\lean{twoSeries_}\allowbreak\lean{hasJetRelations_}\allowbreak\lean{all}, \lean{virSymbol_}\allowbreak\lean{twoSeries}, \lean{jetMap_}\allowbreak\lean{injective} & singular-vector inputs; symbol $\Rightarrow$ injective & unconditional (inputs proved) & Def.~\ref{l20:def:virinputs}, Thm.~\ref{l20:thm:virinputs}\\
\midrule
\multicolumn{4}{l}{\emph{$p=3$ conductor files (Section~\ref{l20:sec:p3}); finite algebra, reused by Chapter~\ref{ch:L24}}}\\
\lean{collisionFactor_}\allowbreak\lean{dvd_}\allowbreak\lean{iff}, \lean{collisionFactor_}\allowbreak\lean{prime}, \lean{collisionFactor_}\allowbreak\lean{sq_}\allowbreak\lean{dvd_}\allowbreak\lean{of_}\allowbreak\lean{swap_}\allowbreak\lean{invariant}, \lean{vandermonde_}\allowbreak\lean{sq_}\allowbreak\lean{dvd_}\allowbreak\lean{of_}\allowbreak\lean{symmetric_}\allowbreak\lean{diagonal_}\allowbreak\lean{vanishing} & collision divisibility; square factors & char $0$ where stated; unconditional & Thm.~\ref{l20:thm:collision}\\
\lean{prod_}\allowbreak\lean{pow_}\allowbreak\lean{dvd_}\allowbreak\lean{of_}\allowbreak\lean{pairwise_}\allowbreak\lean{prime_}\allowbreak\lean{varying}, \lean{rootedGordonConductor_}\allowbreak\lean{dvd_}\allowbreak\lean{of_}\allowbreak\lean{root_}\allowbreak\lean{and_}\allowbreak\lean{collision_}\allowbreak\lean{primePowers} & separated prime powers multiply & unconditional & Thm.~\ref{l20:thm:rooted}\\
\lean{card_}\allowbreak\lean{increasingPair_}\allowbreak\lean{fin}, \lean{p3ConductorWeightedDegree_}\allowbreak\lean{eq} & $\binom d2$; $\deg C_{N,d}=dN$ & unconditional & Thm.~\ref{l20:thm:p3degree}\\
\lean{P3LinearResidualization.}\allowbreak\lean{finrank_}\allowbreak\lean{le_}\allowbreak\lean{expectedDim} & route's terminal bound & conditional on \lean{D} and two hypotheses, none instantiated; unused & \S\ref{l20:sec:p3}\\
\bottomrule
\end{longtable}
\endgroup

%======================================================================
\section{Notes and references}\label{l20:sec:notes}
%======================================================================

\textbf{The theorems.} T.~Arakawa and X.~Dai, ``Classical freeness of affine vertex algebras of $\mathfrak{sl}_2$ at boundary
admissible levels'', arXiv:2609.33348 (2026); Theorems~1.1 and 1.2, equations (2.15), (5.8), (5.12), (5.16), (5.33), (5.37)
as cited in the Lean docstrings. The conjecture: G.~E.~Andrews, J.~van~Ekeren and R.~Heluani, ``The singular support of the
Ising model'', \emph{Int.\ Math.\ Res.\ Not.} 2023, and J.~van~Ekeren and R.~Heluani, ``Chiral homology of elliptic curves and
the Zhu algebra'', \emph{Comm.\ Math.\ Phys.} 386 (2021), 495--550 (classical freeness, singular support).

\textbf{Li filtration, arcs, associated varieties.} H.~Li, ``Abelianizing vertex algebras'', \emph{Comm.\ Math.\ Phys.} 259
(2005), 391--411; T.~Arakawa, ``A remark on the $C_2$-cofiniteness condition on vertex algebras'', \emph{Math.\ Z.} 270
(2012), 559--575; T.~Arakawa and A.~Moreau, ``Arc spaces and chiral symplectic cores'', \emph{Publ.\ RIMS} 57 (2021),
795--829; T.~Arakawa, C.~Jiang and A.~Moreau, ``Simplicity of vacuum modules and associated varieties'', \emph{J.\ \'Ec.\
polytech.\ Math.} 8 (2021), 169--191 (the $C_2$ relations at boundary levels).

\textbf{Admissible representations and characters.} V.~G.~Kac and M.~Wakimoto, ``Modular invariant representations of
infinite-dimensional Lie algebras and superalgebras'', \emph{Proc.\ Nat.\ Acad.\ Sci.\ USA} 85 (1988), 4956--4960;
V.~G.~Kac and M.~Wakimoto, ``On rationality of $W$-algebras'', \emph{Transform.\ Groups} 13 (2008), 671--713 (boundary levels);
V.~G.~Kac and D.~A.~Kazhdan, ``Structure of representations with highest weight of infinite-dimensional Lie algebras'',
\emph{Adv.\ Math.} 34 (1979), 97--108 (the determinant formula); D.~Adamovi\'c and A.~Milas, ``Vertex operator algebras
associated to modular invariant representations for $A_1^{(1)}$'', \emph{Math.\ Res.\ Lett.} 2 (1995), 563--575.

\textbf{Functional realisations and fermionic formulas.} B.~Feigin and A.~Stoyanovsky, ``Quasi-particles models for the
representations of Lie algebras and geometry of flag manifold'', arXiv:hep-th/9308079 (1993); G.~E.~Andrews, ``An analytic
generalization of the Rogers--Ramanujan identities for odd moduli'', \emph{Proc.\ Nat.\ Acad.\ Sci.\ USA} 71 (1974),
4082--4085; B.~Gordon, ``A combinatorial generalization of the Rogers--Ramanujan identities'', \emph{Amer.\ J.\ Math.} 83
(1961), 393--399; C.~G.~J.~Jacobi, \emph{Fundamenta nova theoriae functionum ellipticarum} (1829) (the cube identity).

\textbf{Physics.} C.~Beem, M.~Lemos, P.~Liendo, W.~Peelaers, L.~Rastelli and B.~C.~van~Rees, ``Infinite chiral symmetry in four
dimensions'', \emph{Comm.\ Math.\ Phys.} 336 (2015), 1359--1433, arXiv:1312.5344; C.~Beem and L.~Rastelli, ``Vertex operator
algebras, Higgs branches, and modular differential equations'', \emph{JHEP} 2018(8), 114, arXiv:1707.07679; C.~Cordova and
S.-H.~Shao, ``Schur indices, BPS particles, and Argyres--Douglas theories'', \emph{JHEP} 2016(1), 040, arXiv:1506.00265;
T.~Creutzig, ``W-algebras for Argyres--Douglas theories'', \emph{Eur.\ J.\ Math.} 3 (2017), 659--690, arXiv:1701.05926.

\textbf{The Virasoro $(2,2m+3)$ series (Section~\ref{l20:sec:vir}).} J.~van~Ekeren and R.~Heluani, ``Chiral homology of elliptic
curves and the Zhu algebra'', \emph{Comm.\ Math.\ Phys.} 386 (2021), 495--550 (classical freeness of these minimal models);
A.~Rocha-Caridi, ``Vacuum vector representations of the Virasoro algebra'', in \emph{Vertex Operators in Mathematics and Physics},
MSRI Publ.\ 3 (1985), 451--473 (the characters); G.~E.~Andrews, ``An analytic generalization of the Rogers--Ramanujan identities for
odd moduli'' (cited above) and B.~Gordon (cited above) for the sum sides; B.~Feigin and A.~Stoyanovsky (cited above) for the
quasi-particle basis; M.~Benoit and Y.~Saint-Aubin, ``Degenerate conformal field theories and explicit expressions for some null
vectors'', \emph{Phys.\ Lett.\ B} 215 (1988), 517--522 (the singular-vector recursion). For the physics: Beem--Lemos--Liendo--Peelaers--Rastelli--van~Rees
and Cordova--Shao (cited above) identify $M(2,2m+3)$ with the chiral algebra of the $(A_1,A_{2m})$ Argyres--Douglas theory.
Project record: \texttt{Notes/CLASSICAL\_FREENESS\_STATUS.md}.

\textbf{The $p=3$ route (Section~\ref{l20:sec:p3}).} \texttt{Notes/P3\_FUNCTIONAL\_CARRIER\_AUDIT.md},
\texttt{Notes/BOUNDARY\_AFFINE\_CLASSICAL\_FREENESS.md} (historical). The external input it would have needed is the finite-freeness
theorem of I.~Dumanski and E.~Feigin for global Demazure modules (type $A_1$, level two), as cited in those notes.

\textbf{Project records.} \texttt{Notes/ARAKAWA\_DAI\_2609\_33348\_CAMPAIGN.md} (dependency ledger); the umbrella
\leanpath{VOA/ArakawaDai.lean}; the audit \texttt{scripts/audit/ArakawaDaiAxiomAudit.lean}.

\endgroup
